\documentclass[12pt]{article}
\usepackage[backend=biber,natbib=true,style=alphabetic,maxbibnames=10]{biblatex}
\usepackage[utf8]{vietnam}
\addbibresource{/home/nqbh/reference/bib.bib}
\usepackage{tocloft}
\renewcommand{\cftsecleader}{\cftdotfill{\cftdotsep}}
\usepackage[colorlinks=true,linkcolor=blue,urlcolor=red,citecolor=magenta]{hyperref}
\usepackage[dvipsnames]{xcolor}
\usepackage{amsmath,amssymb,amsthm,float,graphicx,mathtools,soul,subfig}
\allowdisplaybreaks
\newtheorem{assumption}{Assumption}
\newtheorem{conjecture}{Conjecture}
\newtheorem{corollary}{Corollary}
\newtheorem{definition}{Definition}[section]
\newtheorem{dinhnghia}{Định nghĩa}[section]
\newtheorem{example}{Example}
\newtheorem{hypothesis}{Hypothesis}
\newtheorem{lemma}{Lemma}
\newtheorem{notation}{Notation}
\newtheorem{note}{Note}
\newtheorem{principle}{Principle}
\newtheorem{problem}{Problem}
\newtheorem{proposition}{Proposition}
\newtheorem{question}{Question}
\newtheorem{remark}{Remark}
\newtheorem{Rule}{Rule}
\newtheorem{theorem}{Theorem}
\usepackage[margin=2cm]{geometry}

\title{Some Thoughts On Learning, Teaching, {\it\&} Research\\Vài Suy Nghĩ Về Việc Học, Việc Dạy, {\it\&} Nghề Nghiên Cứu\\{\Large\sf A Short Novel -- 1 Tiểu Thuyết Ngắn}}
\author{Nguyễn Quản Bá Hồng\footnote{A Scientist {\it\&} Creative Artist Wannabe. E-mail: {\tt nguyenquanbahong@gmail.com}. Bến Tre City, Việt Nam.}\and Nguyễn Quản Trung Nhân\footnote{In the Void: Classified information.}}
\date{\today}

\begin{document}
\maketitle
\begin{abstract}
	This short novel consists of some pieces of my writing, which are able to be shared \& I am willing to share, in order to sharpen my flows of thoughts, to balance my scientific work via various aesthetic forms, \& to track psychologically \& mentally my transitions from boyhood to manhood if there is any.

	Tiểu thuyết ngắn này bao gồm 1 số bài viết của tôi, những bài có thể chia sẻ được \& tôi tự nguyện chia sẻ chúng, để mài bén các dòng suy nghĩ của tôi hơn, để cân bằng với các công việc nghiên cứu khoa học thông qua muôn vàn các hình thái nghệ thuật, \& để đánh dấu các bước chuyển mình về mặt tâm lý \& trí tuệ của tôi từ 1 cậu nhóc trở thành 1 người đàn ông trưởng thành nếu có bất cứ sự thay đổi nào xảy ra.
\end{abstract}
\setcounter{secnumdepth}{4}
\setcounter{tocdepth}{3}
\tableofcontents

%------------------------------------------------------------------------------%

\section{Preliminaries -- Các Khâu Chuẩn Bị}
Lúc còn làm về động cơ đốt trong (combustion engines), có 1 giai đoạn tôi, tác giả thứ nhất, phải học luật \& các vấn đề đạo đức (ethical issues) phát sinh phòng trường hợp research output của đề tài mà tôi làm bị ứng dụng sai vào mục đích quân đội (military purposes) hoặc dân dụng (civil purposes), làm gây thiệt hại nghiêm trọng. Lúc đó tôi chỉ thấy phiền mà không hiểu cho lắm. Nhưng trải qua vài chuyện thì thấy nó cần thiết thật, với nhiều cái hay lắm (thảo nào mấy cuốn sách luật ở Đức dày 1 cách quái dị). Nên tôi nghĩ điều đó cũng cần phải nêu rõ ở đây, khi mà các vấn đề về tâm lý chưa bao giờ là dễ dàng để tiếp cận, để đề phòng các tình huống xấu nhất. Có lẽ thời gian tôi phải đọc tài liệu về luật \& làm các workshop về scientific ethics (i.e., nonfiction) triền miên, ngốn cả thời gian làm nghiên cứu raw của tôi, cuối cùng sau vài năm cũng được bù đắp theo cách dị hợm nhất có thể: tôi dùng nó cho fiction novel -- tiểu thuyết hư cấu. What a joke in the endless series of {\it Infinite Jest} \cite{Wallace2011} viết bởi cây bút thiên tài người Mỹ {\sc David Foster Wallace}.

\subsection{Disclaimer -- Tuyên bố từ chối trách nhiệm}
We, the authors, in the role of the narrators of this novel, clearly state that:
\begin{itemize}
	\item Any character names mentioned in this piece of writing is purely imaginary. If there is any coincidence, so that a person in real life in the three-dimensional physical space, denoted by $\mathbb{R}^3$, feels disturbed or deeply offended. Well, this century is the optimal era of the offensive, we guess. We can apologize, if we feel necessary, but we take no responsibility.
	\item {\it Warning}: This novel contains sexuality, violence, \& profanities. Consider carefully before you decide to read. If you feel disturbed \& even depressed, We will take no responsibility.
\end{itemize}
Chúng tôi, các tác giả, trong vai trò người kể chuyện, tuyên bố 1 cách rõ ràng rằng:
\begin{itemize}
	\item Bất kỳ tên nhân vật nào trong câu chuyện này đều là thuần hư cấu. Nếu xảy ra bất cứ sự trùng hợp nào, để mà 1 người nào đó ngoài đời thực trong không gian vật lý 3 chiều, tạm ký hiệu là $\mathbb{R}^3$, cảm thấy khó chịu hoặc bị xúc phạm sâu sắc. Chà, thế kỷ này là thời đại tối ưu của giới dễ bị xúc phạm, chúng tôi đoán thế. Chúng tôi có thể xin lỗi nếu cảm thấy cần thiết nhưng từ chối nhận trách nhiệm.
	\item {\it Cảnh báo}: Tiểu thuyết có chứa các yếu tố tình dục, bạo lực, \& thô tục. Cân nhắc cẩn thận trước khi bạn quyết định đọc. Nếu bạn cảm thấy khó chịu, thậm chí là trầm cảm, chúng tôi không chịu trách nhiệm.
\end{itemize}

\subsection{Notation \& convention -- Ký hiệu \& quy ước}

\begin{itemize}
	\item $\mathbb{R}^3$: three-dimensional physical space -- không gian vật lý 3 chiều.
	\item $\mathbb{C}^3$: three-dimensional imaginary space -- không gian ảo, tưởng tượng 3 chiều.
	
	Có thể xem đây là không gian nội tâm đầy màu sắc của các bệnh nhân trong viện tâm thần, e.g., quyển {\it Thiên Tài Bên Trái, Kẻ Điên Bên Phải} của tác giả người Trung Quốc {\sc Cao Minh}, hoặc bộ phim tâm lý nổi tiếng \href{https://www.imdb.com/title/tt0073486}{\it One Flew Over the Cuckoo's Nest} (1975) của đạo diễn {\sc Milos Forman}.
	\item $\Psi$: Psychology -- tâm lý học. {\it Lý do}: 3 chữ cái đầu của từ `psychology' là `psy' giống với `psi'\footnote{Không liên tới hãng đồ uống giải khát Pepsi.} -- cách phát âm của 3 ký hiệu Hy Lạp $\psi,\Psi,\varPsi$, see, e.g., \href{https://en.wikipedia.org/wiki/Psychology}{Wikipedia{\tt/}psychology}.
	\item $\Phi$: Philosophy -- triết học. {\it Lý do}: 3 chữ cái đầu của từ `philosophy' là `phi' -- cũng là cách phát âm của 4 ký hiệu Hy Lạp $\varphi,\phi,\Phi,\varPhi$, see, e.g., \href{https://en.wikipedia.org/wiki/Philosophy}{Wikipedia{\tt/}philosophy}.
	\item bhr: Behavior function ${\rm bhr}(t)$ -- hàm hành vi của 1 người $P$ tại thời điểm $t$.
	\item env: Environment function ${\rm env}(t)$ -- hàm môi trường mà 1 con người $P$ hiện đang sinh sống \& phát triển tại thời điểm $t$.
	\item $P$: a typical person -- 1 con người điển hình. P cũng là viết tắt của Point, có thể hiểu như 1 điểm trong không gian Euclidean $\mathbb{R}^3$. Kiểu hiểu này có lợi khi xét tương tác giữa 2 hay nhiều người với nhau cũng tương tự như tương tác giữa các hạt trong lý thuyết kinetics\footnote{1st author được học môn này của lớp Master 2 ở Rennes, Pháp.}, see, e.g., \cite{Tartar2008}.
	\item $P[t;j;\psi;\phi]$: Với 1 người $P$, $t$ là tuổi, $j$ là nghề nghiệp (job) hoặc chức vụ (title), $\psi$ là các đặc điểm tính cách, tâm lý đặc trưng (characters \& psychological characteristics), $\phi$ là các quan điểm hay trường phái triết học (philosophical schools) của người $P$ đó. Ở đây sự phụ thuộc của 3 hàm $j,\psi,\phi$ vào biến thời gian $t$ \& biến không gian (spatial variable) ${\bf x}\in{\rm env}(t)\subset\mathbb{R}^3$, tức môi trường đang làm việc ở thời điểm $t$ được ngầm hiểu. 1 ký hiệu toán học ngắn gọn cho 1 mô hình toán học đơn giản:
	\begin{equation}
		P[t;j;\psi;\phi] = P[t;j(t,{\rm env}(t));\psi(t,{\rm env}(t),j(t,{\rm env}(t)));\phi(t,{\rm env}(t),j(t,{\rm env}(t)))].
	\end{equation}
	{\it Cắt nghĩa}: Công việc của 1 người tại 1 thời điểm $t$ (thường sử dụng là tuổi từ đây trở đi nếu không nói gì thêm) phụ thuộc vào môi trường ${\rm env}(t)$ người đó sống, \& cả môi trường sống lẫn công việc của người đó trong môi trường sống đó sẽ có ảnh hưởng đến tâm lý \& hệ giá trị triết học của người đó. Tất cả các biến nghề nghiệp, môi trường, tâm lý, hệ giá trị triết học sẽ thay đổi theo thời gian $t$ \& không gian theo biến ${\bf x}\in{\rm env}(t)$, i.e., khi người đó thay đổi nơi sinh sống học tập, \& làm việc.
	
	Có 1 trường hợp đặc biệt là nếu người đó có hệ giá trị triết học độc lập với môi trường \& công việc, hoặc người đó chỉ sống trong cái đầu của họ thì có thể ký hiệu gọn lại là $\phi(t)$ để chỉ rõ sự độc lập vào công việc \& môi trường. Nhưng chỉ cho phép điều này với phần triết học $\phi(t)$, vì nếu tâm lý $\psi_P(t)$ của 1 người $P$ mà hoàn toàn độc lập với môi trường ${\rm env}_P(t)$ \& công việc $j_P(t)$ thì chắc người đó đã siêu thoát, hoặc nếu còn sống thì cũng đã đạt tới cảnh giới giác ngộ cỡ level niết bàn trong Phật giáo thì có lẽ nên cân nhắc việc phong thánh hoặc sản xuất xá lị{\tt/}xá lợi là vừa. Còn nếu bất cứ yếu tố nào trong 3 yếu tố gồm công việc $j(t)$, tâm lý $\psi(t)$, \& hệ giá trị triết học $\phi(t)$, hoàn toàn độc lập với biến thời gian $t$ thì phải ngó lại lại đối tượng $P$ đang xem xét có phải là con người, hay thậm chí là vật thể sống hay không để tiết kiệm thời gian nghiên cứu \& chuyển sang 1 đối tượng (sống) khác.
\end{itemize}
{\it Quy ước về ngôi \& cách xưng hô.} Trong các đoạn dẫn chuyện, nếu người dẫn chuyện (the narrator) sử dụng ``Tôi'', tức ám chỉ tác giả đầu (the first author), nếu sử dụng ``chúng tôi'' tức ám chỉ cả 2 tác giả (both authors). Các xưng hô ``Tui'', ``mình'', ``ta'', ``mi'', ``tao'', ``mày'', etc., các cách xưng hô của các nhân vật trong tiểu thuyết. Các quy ước này được tham khảo 1 phần từ quyển {\it First Person Singular: Stories} \cite{Murakami_1st_person} hay bản dịch tiếng Việt {\it Ngôi thứ nhất số ít} \cite{Murakami_ngoi_1} của nhà văn nổi tiếng người Nhật Bản {\sc Haruki Murakami}.

\subsection{Outline}
3 phần chính của tiểu thuyết này là bàn về việc học, bàn về việc dạy, \& bàn về nghiên cứu, theo thứ tự trưởng thành dần theo chiều phát triển của 1 người. Nhưng chúng tôi chợt nghĩ có khi thứ tự đó lại không chuẩn cho lắm, thứ tự mà chúng tôi cho là ``đúng'' \& ``chuẩn'' chính là bàn về nghiên cứu, rồi bàn về việc dạy, \& cuối cùng là bàn về việc học. Nghe có vẻ ngược ngạo, nhưng từ những trải nghiệm của tác giả trên con đường nghiên cứu chuyên nghiệp, chúng tôi nhận ra có nhiều thứ quý giá hơn ngoài văn phòng nghiên cứu. Nên sau khi nêu 1 số suy nghĩ về nghiên cứu ở Sect. \ref{sect: research}, chúng tôi xem lại về việc dạy ở Sect. \ref{sect: teaching}, cuối cùng là bàn về các sự tinh chỉnh không ngừng trong nghiệp\footnote{`Nghiệp'' ở đây có thể hiểu là ``sự nghiệp'' hoặc ``nghiệp chướng'', e.g., các bài báo \href{https://vnexpress.net/huy-chuong-vang-toan-tre-nhat-lam-toan-vi-nghiep-chuong-2393718.html}{VNExpress{\tt/}Huy chương vàng toán trẻ nhất làm toán vì 'nghiệp chướng' }, \href{https://dantri.com.vn/giao-duc/hcv-toan-hoc-tre-nhat-lam-toan-vi-nghiep-chuong-1354447180.htm}{Dân Trí{\tt/}HCV Toán học trẻ nhất: Làm toán vì ``nghiệp chướng''}.} học vấn trong cả cuộc đời của 1 con người trong Sect. \ref{sect: learning}.

\section{Ignition -- Mồi Lửa}

\subsection{When things fall apart -- Khi mọi thứ sụp đổ}
Tên của mục này lấy từ tựa đề của quyển {\it When Things Fall Apart: Heart Advice for Difficult Times} \cite{Chodron2002} của sư mục {\sc Pema Ch\"odr\"on} với bản dịch tiếng Việt {\it Khi Mọi Thứ Sụp Đổ: Lời Khuyên Chân Thành Trong Những Thời Điểm Khó Khăn} \cite{Chodron2021}.
\begin{flushright}
	Christmas Eve, Winter 2020. Berlin, Germany.
	
	In the {\it imaginary world} $\mathbb{C}^3$ instead of the {\it real physical world} $\mathbb{R}^3$.
\end{flushright}
Hắn thấy mình đang mệt nhoài, toàn thân rã rời, lê bước trở về phòng riêng trong 1 student studio ở Alexanderflatz\footnote{ {\it Alexander Square}: a large public square \& transport hub in the central Mitte district of Berlin, reputedly the most visited area of Berlin, see \href{https://en.wikipedia.org/wiki/Alexanderplatz}{Wikipedia{\tt/}Alexanderplatz}.}, Berlin--Mitte, sau buổi tối thân mật cùng đồng nghiệp chung team nghiên cứu của hắn. Nói là đồng nghiệp nhưng thật ra không hắn làm việc chung 1 đề tài. Nếu là 1 công việc bên mảng công nghệ hoặc kinh doanh chắc có lẽ sẽ khác. Đây là công việc nghiên cứu. Mỗi người trong team sẽ chịu trách nhiệm 1 mảng nghiên cứu riêng, dù chung 1 research theme, nhưng khó mà làm chung với nhau được. Đấy là cái thất bại đầu tiên của hắn: hắn không tương tác được với đồng nghiệp để tạo ra ý tưởng mới. Nhưng có nên trách hắn không? Hắn chỉ là 1 nghiên cứu sinh bậc tiến sĩ, đồng nghiệp của hắn toàn postdoc -- nghiên cứu sau tiến sĩ, khi mà họ đã rành mảng của họ \& sẵn sàng tương tác với các đồng nghiệp trong nhóm lẫn những đồng nghiệp khác nhóm, thậm chí các giáo sư dạng Head of Research Groups, thì đây là 1 mảng hoàn toàn mới với hắn. Hắn cần thời gian để cày kiến thức nền (background knowledge) -- nhưng đúng ra thì thời gian không cho phép. Hắn đang trong giai đoạn cuối của 1 đề tài sắp nghiệm thu mà các người tiền nhiệm của hắn đã bỏ hẳn. He felt that this project is like a fucking death star which he was still trying to hold on his weak shoulder:
\begin{figure}[H]
	\centering
	\includegraphics[scale=.7]{spiral}
	\caption{Rust Cohle's hallucination on The Spiral before the last fight with the Yellow King. \href{https://www.imdb.com/title/tt2356777/}{True Detective} (2014--) [S1.E8].}
\end{figure}
Ít ra thì hắn cũng có 1 bài toán nhỏ với thầy của hắn được xuất bản trên 1 tạp chí toán phổ thông nổi tiếng của Đức, tựa tựa tạp chí Toán học \& Tuổi trẻ ở quê nhà của hắn. Hắn mừng lắm. Chẳn phải xuất bản khoa học, nhưng cũng được tính KPI (key performance indicator) để giảm bớt áp lực công việc.

Sáng Giáng sinh năm ấy hắn dậy sớm, ra siêu thị mua đồ ăn ngon. 

There have been many broken pieces of memories in his head. Good memories. Bad memories.

Hắn nhớ lại cha hắn. Chợt hiểu ra vài thứ. Có lẽ hắn không có khả năng để tạo hạnh phúc cho người khác. Nếu có khả năng về mặt thuần tình cảm thì cũng chả có vị thế về vật chất. Hắn chấp nhận rồi buông bỏ.

A dark student room in , opposite the main Police station. No sound. No sign of life. He can't feel anything from himself. He is not presented.

Hắn chợt choàng tỉnh dậy sau cơn mê man, đầu óc mụ mị, nhức như búa bổ. Mình phải tính tay cho xong -- hắn nấu 1 ít rồi ngồi vào bàn \& làm việc. Hắn không muốn những cảm xúc. Hắn cho những thứ đó là cấm kỵ của 1 người làm khoa học. Quá ủy mị. Quá mềm yếu. Hắn cứ lờ đi mặc dù tiềm thức cảnh báo hắn.

\begin{quote}
	God. It's just a fucking simple work. Just get it done. Move to something else.
\end{quote}
Hắn ước gì phiên bản tương lai của hắn sẽ quay trở về giúp hắn, cho hắn 1 lời khuyên hay ho nào đấy hoặc đấm 1 cú trời giáng thẳng vào bộ não ngây thơ của hắn để hắn tỉnh ngộ, thậm chí thức tỉnh. Chả có ai giúp hắn cả. A contradiction in both his drinks \& his subconscious mind. He works in mathematics so long to avoid contradiction, now it comes to his lifestyle, his psyche. Fucking ironic.

Hắn say xưa làm việc. Phòng bên hì hục làm tình. Sự thăng hoa hắn tự cho là đã tìm thấy trong công việc \& sự khoái cảm của bộ trai gái phòng bên hòa hợp, át đi tiếng siren còi hú inh ỏi của đồn cảnh sát. Nhịp nhàng. Giống như vũ điệu reo hò soran, 1 nét văn hóa đẹp của người Nhật. Hắn thấy chị Trinh, 1 junior software developer tại Amazon, nhắn. Chắc lại nhờ vả gì đó, chị ta chỉ có thế, chỉ xuất hiện khi cần cái gì đó hoặc tỏ ra hơn cái gì đó rồi tốc biến không báo trước -- hắn thoáng nghĩ. Hắn chưa bao giờ hiểu tại sao bạn bè hắn lại cảnh báo chị ta nhiều đến thế. Hắn theo chủ nghĩa ôn hòa, nên dẫu có bao nhiêu red flags hắn vẫn mặc kệ. Nhưng ai ai cũng bảo nên cẩn thận với chị ta. Hắn nhớ lại trước đây, lúc còn năm 2 Đại học, ở Ký Túc Xá Đại Học Quốc Gia ở Thủ Đức, bạn chung phòng của hắn cũng cảnh báo hắn về 1 đứa chung lớp hắn như thế.

He knock a drink of beer. He didn't like wines. Then he drank a lot of black coffee, the pure type. He always has a very bad eating \& drinking habit.

Beer so he can forget bad shits, coffee so he can focus on the good. Actually a side effect of coffee helps him shit well too. The cure for all the sittings long hours labor work.

Thay vì chịu sự cô đơn, hắn thỏa hiệp với vài con người thuần lợi dụng. 

What is missing? What pieces of the picture of life do I miss? A little self-esteem? A little understanding how life works?

Hắn tự huyền hoặc bản thân với 1 trí nhớ tốt như 1 cỗ máy điện tử, sự ảo tưởng sức mạnh về món quà mà hắn được ban tặng.

Nah, these things so difficult. I get back to my mathematics, my calculus.

\subsection{Why bad things always happen to good people? -- Tại sao người tốt luôn gặp chuyện xấu?}
What is the point of all of this?

Oh oh, so if I can do all these math, my father's liver cancer will disappear right? he can live right? He can be normal again like nothing happens right?

Fucking bullshit. All of this. Purely fucking bullshit.

Hắn tự hỏi liệu Chí Phèo đã cảm giác thế nào vào cái lúc hắn vừa đi vừa chửi. Cái làng Vũ Đại ngày ấy. Vãi đạn thật.


Then he felt awake. Something inside him awoke. He was not sure about it but totally aware of its presence.

This is a very different battle - the one I can't win now. If you lost a battle within yourself then how the hell you can win any other battle?

He had worked so hard to be able to take all the responsibilities of the only job offered to him, then he got judged to be lazy \& irresponsible in many senses.



This is the battle he will lose no matter how.

Do I want to die as a good man? Like my father? No, I want to die as a wise man. 1st, I need to confront my stupidities.

The man who lost his faith in God finds a way to get insight. The journey of revelation begins.

\cite{Tolle2022}

Oppenheimer Homocunlus

You mean Newton is the fuck boy of science? Yo, no! The player. Like Mozart - the player of musical instruments. Here I mean the player of formulas \& concepts.

My conscious mind has still been working with this idea no matter how I deny \& destroy all pieces of my writings. The day my teacher die, the day I remember my father die. Consonant echoes from the past into the present. 

It is about the passion.

Power of concentration -- black fist in JJK.

\begin{center}\it
	Feeds the dying light, \& brings me back to life\\
	Let the darkness lead us into the light 
	-- Ignite, {\sc Alan Walker}
\end{center}
Then the most dangerous intrusive thought won:

Hay là mình tạm ngưng làm toán mà thử làm nhà văn nhỉ? -- You're fucking kidding me right? Don't fool yourself. Joke on all of us.

Winter that year so cold, colder than before. He felt the damn cold in his own mind \& own heart.

He has no literary gift, not any gift for conceptual thoughts, just some average mathematical computation skills, not abstract enough to become a pure mathematician, not complex or useful enough to become an applied mathematician. Who can he become then?

But the day his son, if he has any, asks him about these, what kind of father he is then?

I don't know what the fuck do you want. I really have no idea. \& I am so fucking tired to be pretentious. Listen. Pick your pieces up.

Solitude is inevitable.

Đắm mình trong đại dương trầm cảm đủ lâu, hắn cảm thấy 1 sự tự do \& thanh thoát tột cùng, như kiểu đả thông kinh mạch trong \href{https://www.imdb.com/title/tt0373074/}{\it Kung Fu Hustle} (2004) với tựa Việt {\it Tuyệt Đỉnh Kungfu}, nhưng ở đây là đã thông thế giới quan \& nền tảng giá trị của hắn. Nhìn mọi thứ rõ thế này thích thật.

Khi bạn rời khỏi 1 địa hạt nào đó, những người sùng bái, những tín đồ của địa hạt, tôn giáo đó có thể coi bạn như kẻ bại trận (loser), đồ súc vật (animal), thứ rác rưởi (trash), quân phản trắc (traitor), etc., để có lý do tự cho mình cái quyền chà đạp lên bạn. Chả sao cả. Bạn không còn lãnh trách nhiệm mình là 1 mắc xích trong cái luồng công việc ở địa hạt đó nữa. Bạn dành nhiều thời gian hơn cho bản thân, để phát triển bản thân, tự do làm những điều mình thích, bên cạnh việc kiếm tiền trong trường hợp việc kiếm tiền không dính dáng đến sở thích hiện tại của bạn. This is the kind of the ultimate freedom mentioned in \href{https://www.imdb.com/title/tt0137523}{\it Fight Club} (1999).
\begin{quote}\it
	``\& then $\ldots$ something happened. I let go. Lost in oblivion. Dark \& silent \& complete. I found freedom. Losing all hope was freedom.'' -- Fight Club (1999)
\end{quote}
Naturally \& inevitably, in the autumn of 2021, after the death of his mathematics teacher in high school, in another dark room in The Student Hotel in Vienna, Austria, he decided to fight his primal fears, in the realms of literary -- word \& emotion analysis, \& psychology -- behavior \& emotion analysis, with the help \& guidance of philosophy -- value \& meaning analysis -- so that he can be lost in his flow(s?) of thoughts less frequently, or at least being lost in a positive way this time.

Maybe he never nails it, the kind of thoughts \& writings he always wants to pursue. Is it important though? Just another lose. He is used to it since he is a loser. But it is still a good try, a way to live aesthetically, spiritually, which means to be able to live properly -- he supposed \& still believes.

\begin{quotation}\it
	``In this world, is the destiny of mankind controlled by some transcendental entity or law? Is it like the hand of God hovering above? At least it is true that man has no control, even over his own will. Man takes up the sword in order to shield the small wound in his heart sustained in a far-off time beyond remembrance. Man wields the sword so that he may die smiling in some far-off time beyond perception.'' -- {\sc Kentaro Miura}, {\it Berserk}, Vol. 1
\end{quotation}
Tạm dịch: {\it Trên thế giới này, vận mệnh của loài người có phải do 1 thực thể hay quy luật siêu việt nào đó điều khiển không? Có giống như bàn tay của Chúa đang lơ lửng trên cao không? Ít nhất thì đúng là con người không thể kiểm soát được, ngay cả ý chí của chính mình. Con người cầm kiếm để che chắn vết thương nhỏ trong lòng mình đã phải chịu đựng từ 1 thời gian xa xôi không thể nào quên được. Con người sử dụng thanh kiếm để có thể chết với 1 nụ cười ở 1 thời điểm xa xôi nào đó ngoài tầm nhận thức}.

\begin{figure}[H]
	\centering
	\includegraphics[scale=.2]{Berserk_behelit_color}
	\caption{The Hand of God holds the Behelit in Berserk.}
\end{figure}

\begin{quotation}\it
	``Providence may guide a man to meet 1 specific person, even if such guidance eventually leads him to darkness. Man simply cannot forsake the beauty of his own chosen path. When will man learn a way to control his soul?'' – {\sc Kentaro Miura}, {\it Berserk}
\end{quotation}
Tạm dịch: {\it Mệnh trời có thể chỉ đường 1 gã đàn ông gặp 1 người cụ thể, ngay cả khi sự hướng dẫn đó cuối cùng lại dẫn hắn ta vào bóng tối. Gã đàn ông đơn giản là không thể từ bỏ vẻ đẹp của con đường mình đã chọn. Khi nào gã sẽ học được cách kiểm soát linh hồn của chính mình?}
	
\begin{quote}
	{\sf Hồng [?; humble philosopher ready to die]}: Do you know what is one of the biggest regrets in life?
	
	-- Cậu có biết 1 trong những điều đáng tiếc nuối nhất trên đời này là gì không?
	
	{\sf Hồng [23--25; depressed mathematician wannabe]}: I have no idea. How the hell can I know it? There are too many regrets in this fucking life in all aspects. What is it according to you?
	
	-- Tôi không biết. Thế đéo nào mà tôi biết được cơ chứ. Có quá nhiều điều để đáng tiếc trong cái cuộc sống chó chết này theo nhiều phương diện. Thế theo ông là cái quái gì?
	
	{\sf Hồng [?; humble philosopher ready to die]}: A gifted, talented or even multi-talented man without any piece of art, with no contribution \& no legacy for the next generation(s) after him. What a terrible \& accursed waste of talents \& giftedness!
	
	-- 1 gã có tài, thậm chí nhiều tài hay nhiều loại năng khiếu, nhưng lại không chịu{\tt/}không thể (?) để lại bất cứ tác phẩm nghệ thuật nào, không có 1 chút cống hiến hay truyền lại bất cứ di sản nào cho (các) thế hệ tiếp theo sau hắn. 1 sự phí phạm tài năng đến mức kinh khủng \& đáng bị nguyền rủa.
\end{quote}

\section{An Initial Configuration-- 1 Thiết Lập Ban Đầu}
Bằng cách đặt ra 1 số nguyên lý (principles) \& quy tắc (rules), chúng tôi tin là `cuộc thập tự chinh' này sẽ có chỗ để bám níu vào, chúng có tác dụng như những cái mỏ neo, ít nhất để không bị lạc giữa đại dương tri thức mênh mông với những con quái vật của sự vô tri (ignorance monsters) \& những con quỷ của sự ngu dốt \& ác độc (devils of stupidities \& cruelties) chầu chực chờ sẵn dưới đáy đại dương sâu thẳm -- the deep bottom of a dark ocean -- có thể vồ túm lấy ta bất cứ lúc nào khi mà nhận thức của ta bị đánh lừa \& rơi vào trạng thái u mê mất cảnh giác.

\subsection{Rules -- Các quy tắc}
Mỗi người $P$ (abbr., person) có một xuất phát điểm $\{P(t)\}_{0\le t\le t_0}$ khác nhau, được hưởng hoặc bị ép nhồi các nền tảng giáo dục khác nhau, sự tương tác với những người khác nhau, cùng vô vàn những chuyện \& những biến cố họ gặp trong suốt 1 cuộc đời hoàn toàn khác nhau, thành ra nền tảng nhận thức \& xu hướng phát triển nhận thức, cùng sự hình thành các cấu trúc niềm tin \& các hệ giá trị cơ bản cùng thế giới quan của mỗi người hoàn toàn khác nhau.

Quy tắc đầu tiên ở đây là:
\begin{Rule}[On judgment]
	Không phán xét, công kích, e.g., dí trên mạng xã hội, bất cứ ai. Cũng không áp đặt ai, thậm chí cả việc áp đặt ai đó không được áp đặt người khác. Tạo cho người khác 1 cảm giác thoải mái tối thiểu khi tiếp xúc.
\end{Rule}

\begin{Rule}[On stalking]
	Không quá tò mò vào cuộc sống cá nhân của người khác, e.g., stalk in social media -- rình mò trên các nền tảng mạng xã hội, xâm phạm tài khoản  riêng tư cá nhân bất hợp pháp. Keep healthy boundaries for both.
\end{Rule}

\begin{Rule}[On system reset]
	Một phản tư xa hơn trong tương lai có lẽ là chẳng có hành trình phát triển tự thân nào mà đủ sức chống chọi 1 cách hiệu quả với các tương tác xã hội cả, đặc biệt là các tương tác xấu \& các mối quan hệ độc hại (toxic relationships) cả. Khi đó thì tất cả các ghi chú ở đây sẽ bị xóa. Mọi thứ trở về cấu hình sống nhiều mặt phổ dụng để che giấu bản thân.
\end{Rule} 

\subsection{Goals -- Các mục tiêu}
This writing activity is one of many ways, which is likely to become the main one, to balance between my scientific work \& personal life. I believe some arts will be the tool.

Việc viết lách, theo mình nghĩ, bằng cách này hay cách khác, một lúc nào đó \& theo 1 cách tự nhiên nào đó, cũng sẽ tìm tới những kẻ thích suy nghĩ, những kẻ hay nghĩ nhiều, \& những kẻ mệt mỏi vì cái tật đó, e.g., nhà nghiên cứu, nghiên cứu sinh, các học giả, nói chung là những người làm trong mảng học thuật hoặc phải tiếp xúc nhiều với chữ. Tật hay tài thì chưa biết nhưng ắt hẳn việc viết dùng để sắp xếp mọi thứ trong đầu cho ngăn nắp thì không thể tránh khỏi đối với những người làm việc đầu óc nhiều.

\section{On Writing: Literary writing for a literary retard -- Bàn về việc viết: Học viết văn cho kẻ dốt đặc văn}

\begin{quote}
	{\sf Hồng [4--23.5; typical theoretical agreeable giver]}: But you have to be useful to other people right? You have to write something nice, something educated, something makes people happy, even in a fake way. Whether you are happy when writing or not does not matter.
	
	{\sf Hồng [26--?; critically practical disagreeable writer]}: Fuck off. I write because I like to write, because I need to write to be able to understand myself, \& because I will be good at writing no matter how. I write for me 1st. Only if I find it useful for other people, I will share to the class of people who need it. You fucking slave! Go fuck yourself.
	
	{\sf Hồng [4--23.5; typical theoretical agreeable giver]}: But you have to be useful to other people \& forget yourself. You have to sacrifice yourself when needed. Only in that way, you can live properly as a good person.
	
	{\sf Hồng [26--?; critically practical disagreeable writer]}: I don't give a single fuck \cite{Manson_giving_fuck,Manson_giving_fuck_vn}. Fucking Jesus Christ!
	
	{\sf Hồng [4--23.5; typical theoretical agreeable giver]}: Hey man. Let me put this matter like this: I know he has been being fixed in the Cross. So he cannot do anything to react or response to your physical action. But please for the mercy of God, don't $\ldots$ Do not do it you $\ldots$ filthy animal.
	
	{\sf Hồng [26--?; critically practical disagreeable writer]}: Oh my fucking gosh! It is true that I am an atheist but what in the bloody hell are you even thinking about my sexuality? Let me alone. Get the fuck out of my head. I need to write.
\end{quote}

\subsection{I want to be a writer -- Tôi muốn trở thành 1 nhà văn}
Tên phần này bắt chước tựa đề của quyển {\it I want to be a mathematician: An Automathography} \cite{Halmos1985,Halmos1985_3_parts} (tạm dịch: {\it Tôi muốn trở thành 1 nhà toán học: tự truyện cho nghề làm toán}) của nhà toán học \& nhà triển lãm nghệ thuật về khoa học nổi tiếng người Áo--Hungary {\sc Paul Halmos}.
\begin{quote}
	{\sf Hồng [25; writer wannabe, literary retard]}: Tôi muốn học viết, anh có cách nào hay sách nào chỉ tôi với.
	
	{\sf Hồng [28--?; writer]}: Nếu anh mới bắt đầu thì tôi nghĩ anh nên đọc quyển {\it The Elements of Style} \cite{Strunk1918} (tạm dịch: {\it Các yếu tố của phong cách}) của tác giả người Mỹ {\sc William Strunk Jr.} hoặc bản tái bản có bổ sung \cite{Strunk_White2019} của ổng \& đệ tử {\sc E. B. White} để hiểu tầm quan trọng của phong cách viết chuẩn mực, ngắn gọn, súc tích.
	
	{\sf Hồng [25; writer wannabe, literary retard]}: Sau đó thì đến quyển nào? Tại tôi thấy quyển này khá ngắn, \& tôi có thể đọc khá nhanh. Nên tôi nghĩ tôi cần nhiều hơn trong lần `bàn giao tri thức' đầu tiên này. Tôi đoán thế.
	
	{\sf Hồng [28--?; writer]}: Đồng ý là quyển {\it The Elements of Style} khá ngắn, nhưng cần thời gian để cảm thụ các quy tắc. Quyển đấy nhỏ mà có võ. Anh sẽ hiểu ý của tôi sớm thôi khi bắt tay vào viết 1 thứ gì đó của riêng anh.
	
	{\sf Hồng [25; writer wannabe, literary retard]}: Anh thông cảm. Tôi hơi tham vọng về mặt tri thức. Anh có thể gọi là tham lam cũng được. Ambitious \& intellectually greedy. What's so different though? Anh có thể cho tôi thêm tên vài quyển nữa được không. Phòng trường hợp tôi đọc quyển đầu nhanh quá nên xong, hoặc chán hay khó quá nên (tạm) ngưng.
	
	{\sf Hồng [28--?; writer]}: Chiều ý anh luôn. Tiếp theo là quyển {\it On Writing: A Memoir of the Craft} \cite{King2000,King2010} của nhà văn nổi tiếng về truyện kinh dị người Mỹ {\sc Stephen King}, quyển {\it On Writing Well: The Classic Guide to Writing Nonfiction} \cite{Zinsser2001,Zinsser2016} của {\sc William Zinsser}, đệ tử của {\sc E. B. White}, mà {\sc E. B. White} lại là đệ tử của {\sc William Strunk Jr.}, nên đây là bộ 3 sư phụ--đồ đệ của 3 thế hệ liên tiếp. Sẵn tiện nếu anh muốn viết về quá khứ, kiểu về các kỷ niệm tuổi thơ đẹp đẽ hoặc thậm chí là các ám ảnh tuổi thơ, mà 1 khi trưởng thành anh muốn hiểu, thì nên đọc quyển {\it Writing About Your Life: A Journey into the Past} \cite{Zinsser2005} cũng của {\sc William Zinsser}. Quyển {\it Still Writing: The Perils \& Pleasures of a Creative Life} của {\sc Dani Shapiro} tôi chưa đọc hết nhưng anh cũng nên thử. Tóm lại là anh nên xem các quyển sau:
\end{quote}
\noindent\textbf{\textsf{Resources -- Tài nguyên.}}
\begin{itemize}
	\item \cite{Can_dtnv}. Nguyễn Duy Cần. {\it Để Trở Thành Nhà Văn}.
	\item \cite{Strunk1918}. William Strunk Jr. {\it The Elements of Style}.
	\item \cite{Strunk_White2019}. William Strunk Jr., E. B. White. {\it The Elements of Style}.
	\item \cite{King2000,King2010}. Stephen King. {\it On Writing: A Memoir of the Craft}.
	\item \cite{Zinsser2001,Zinsser2016}. William Zinsser. {\it On Writing Well: The Classic Guide to Writing Nonfiction}.
	\item \cite{Zinsser2005}. William Zinsser. {\it Writing About Your Life: A Journey into the Past}.
	\item \cite{Shapiro2014}. Dani Shapiro. {\it Still Writing: The Perils \& Pleasures of a Creative Life}.
\end{itemize}

\begin{quote}
	{\sf Hồng [25; writer wannabe, literary retard]}: Biết là hơi ngu \& quá gấp, nhưng có quyển nào nặng đô không? Để tôi bào từ từ.
	
	{\sf Hồng [28--?; writer]}: Các quyển tôi vừa kể là bàn về việc học cách viết cơ bản. Còn nếu anh muốn kiểu thực chiến, bay thẳng vào trận mạc, kiểu trầy da tróc vảy, bầm dập nhừ tử tương để biết cách chiến đấu, thì anh nên tìm đọc tiểu thuyết {\it The Fountainhead} của {\sc Ayn Rand} với bản dịch tiếng Việt {\it Suối Nguồn} \cite{Rand_fountainhead} \& tiểu thuyết {\it Infinite Jest} \cite{Wallace2011} của {\sc William Foster Wallace}. Nói ngắn gọn cho anh dễ hiểu: {\it Uproarious \& close to madness}.
\end{quote}

\subsection{Authenticity -- Tính xác thực}

\begin{quote}
	{\sf Hồng [25; writer wannabe, literary retard]}: Lúc tôi mới từ Đức về thì tôi nản lắm. Anh Quang Vinh, nghệ danh Vinhmath, trước tôi 1 khóa có nhắn đểu cho tôi. Anh này cấp 3 chuyên Tin, nhưng lên Đại học thì học ngành Toán-Tin giống tôi. Nếu có thể diễn tả ngắn gọn về ảnh thì tôi nghĩ là: ảnh tính xác suất của 1 sự kiện ra mấy triệu, mấy tỷ mà vẫn tự tin khoe chiến tích tính toán rầm rầm trên mạng xã hội cho các học viên là các bé cấp 2, cấp 3 dễ dụ, dễ bị thao túng. Anh ta không phân biệt được 2 phân môn Xác suất (Probability) \& Thống kê (Statistics) có thể do 2 môn này thường được gom chung thành môn Xác suất thống kê cho các ngành không sử dụng Toán học nhiều. Nhưng anh Vinh là sinh viên ngành Toán Tin giống tôi cơ mà. Nói chung chả hiểu. Anh ta nói khá nhiều, về nhiều lĩnh vực \& mỗi lĩnh vực thì nói cũng nhiều. Nói chung là viết dài lắm mà ý hông bao nhiêu. Như thằng Tử, bạn học chung lớp chuyên Toán của tôi hồi cấp 3 hay nói là: ``Xuất tinh ra 1 đống mà toàn là nước, chả có con nòng nọc nào trong đó''. Kiểu kiểu vậy, phép ví von hơi kỳ nhưng chắc anh hiểu ý tôi.
	
	{\sf Hồng [28--?; writer]}: Tôi hiểu. Nhưng ý chính là gì?
	
	{\sf Hồng [25; writer wannabe, literary retard]}: Tôi có thấy 1 lần anh Vinh viết về bà Phượng Hằng \& thói cổ xúy. Những đứa dễ bị thao túng bởi dư luận hò reo tán thưởng inh ỏi. Trong khi những người có xíu phản biện thì ném đá ảnh ầm ầm, lăng nhục đủ kiểu, khiến ảnh phải xóa bài đăng trên mạng xã hội Facebook, bị report nhiều đến nỗi đến mất cả nick, \& còn cạo cả đầu nữa. Chỉ vì 1 cuộc chiến võ mồm trên mạng xã hội mà phải chịu làn sóng phẫn nộ kinh khủng thế. Thế theo anh tại sao anh Vinh lại bị nhiều người ném đá trên mạng \& sỉ nhục không thương tiếc như vậy?
	
	{\sf Hồng [28--?; writer]}: Vấn đề nằm ở {\it tính xác thực -- authenticity} ({\sf[n] the quality of being true or what somebody claims it is}). Anh phải biết anh đang nói gì với người khác. Nếu anh tỏ ra trịch thượng, kiểu đứng top, đứng nhất là ta đây, còn tụi bây toàn lũ hạ đẳng, thì coi như xong. Chả có cuộc thảo luận nào diễn ra cả. Chỉ là 1 sự phô trương kiến thức phù phiếm, chưa kể kiến thức đó có đúng hay không, như kiểu xác suất mà lớn hơn $1$, à nhầm, xác suất mà có thể đạt tới mấy triệu, mấy tỷ như anh vừa kể, thì chả ai có kiến thức nể cả. 1 trò hề đúng nghĩa. Bị ném đá, sỉ nhục là chuyện tất yếu.
	
	{\sf Hồng [25; writer wannabe, literary retard]}: Mà điều hóm hỉnh nhất anh biết là gì không? Chính là anh Vinh đăng hình đi ăn với thầy Minh -- người nổi tiếng bắt bớ các giáo viên Toán soạn đề không đúng ngữ pháp tiếng Việt theo ý thầy ấy. Nói nôm na anh cứ hiểu thầy Minh là chuyên gia chuyên trị những ông thầy, bà cô dạy Toán sơ cấp, tức Toán cấp 2, cấp 3, sai, thế đéo nào mà 2 người lại thân nhau, đăng hình bữa tối thân mật như thế lên mạng xã hội?
	
	{\sf Hồng [28--?; writer]}: Anh hiểu nghĩa của câu thành ngữ {\it ``Không vào hang hùm, sao bắt được cọp con''} không? Chỉ đơn giản vậy thôi. Để củng cố quyền lực 1 cách ám mụi bằng cách tỏ ra gần gũi với những tai to mặt lớn để nâng tầm bản thân lên thay vì tự trau dồi bản thân bằng cách tính xác suất lại cho kỹ \& phân biệt rành rọt lại 2 phân môn Xác suất \& Thống kê, để truyền đạt kiến thức đúng cho các học viên chịu theo học mình vì mình giỏi, vì kiến thức của mình, chứ không phải theo vì dễ bị thao túng, quay mòng mòng theo cái dòng dư luận trên mạng xã hội chả biết đi đâu \& về đâu nữa.
	
	{\sf Hồng [25; writer wannabe, literary retard]}: Có thể là 1 tiệm cắt tóc xịn xò để cạo trọc hay nối tóc lại cho đẹp chẳng hạn?
	
	{\sf Hồng [28--?; writer]}: Anh tập trung viết những vấn đề khác giùm tôi \& ít quan tâm mạng xã hội hơn đi. Theo tôi, để viết tốt, anh cần 1 nơi yên tĩnh, tách biệt với thế giới bên ngoài, nơi những tiếng ồn xe cộ, tiếng con người tất bật mưu sinh làm anh mất tập trung vào dòng chảy văn chương trong anh nếu giả sử có dòng chảy văn chương nào đó chịu chảy qua anh. Vào trong rừng ở 1 thời gian không phải là 1 ý tồi. Tùy anh quyết.
	
	{\sf Hồng [25; writer wannabe, literary retard]}: Nhưng tôi sợ ở 1 mình. Không, ý tôi là ở 1 mình đôi khi cũng tốt. Nhưng nếu cô độc quá lâu, thì con người ta sẽ tự nhiên xuất hiện thêm 1 giọng nói trong đầu không kiểm soát được phải hông? Có khi nhiều giọng nói nữa. Chắc anh thừa biết ở 1 mình cô đơn trong phòng kín quá lâu sẽ có tác hại thế nào tới sức khỏe tâm lý.
	
	{\sf Hồng [28--?; writer]}: Bingo! Exactly.
	
	{\sf Hồng [25; writer wannabe, literary retard]}: Tôi chưa hiểu?
	
	{\sf Hồng [28--?; writer]}: Let solitude invade your entire physical body \& the whole of your mind \& then try to master that voice or these inner voices in your head. Turn those from inner enemies into your allies. You will get it when you have already mastered it.
	
	Cứ để sự cô độc gặm nhắm toàn bộ cơ thể vật lý của anh \& toàn bộ tâm trí anh, rồi sau đó cố gắng làm chủ (các) giọng nói trong đầu anh. Biến chúng từ kẻ thù bên trong anh thành đồng minh của anh. Anh rồi sẽ hiểu ý tôi thôi khi đã làm chủ được chúng.
\end{quote}

\subsection{Into the wild \& a peaceful life in the woods -- Tìm về chốn hoang dã \& 1 cuộc sống yên ả trong rừng sâu}
\textbf{\textsf{Resources -- Tài nguyên.}}
\begin{itemize}
	\item \cite{Thoreau2014}. Henry David Thoreau. {\it Walden or Life in the Woods} với bản dịch tiếng Việt {\it Walden: Một Mình Sống Trong Rừng} \cite{Thoreau2023}.
	\item \cite{Krakauer1997}. Jon Krakauer. {\it Into The Wild} với chuyển thể thành phim \href{https://www.imdb.com/title/tt0758758}{\it Into the Wild} (2007) của đạo diễn người Mỹ {\sc Sean Penn}.
\end{itemize}
Trong phim \href{https://www.imdb.com/title/tt0097165}{\it Dead Poets Society} (1989) có nhắc đến đoạn này trong giờ giảng văn:
\begin{quote}\it
	``I went to the woods because I wished to live deliberately, to front only the essential facts of life, \& see if I could not learn what it had to teach, \& not, when I came to die, discover that I had not lived. I did not wish to live what was not life, living is so dear; nor did I wish to practise resignation, unless it was quite necessary. I wanted to live deep \& suck out all the marrow of life, to live so sturdily \& Spartan-like as to put to rout all that was not life, to cut a broad swath \& shave close, to drive life into a corner, \& reduce it to its lowest terms, and, if it proved to be mean, why then to get the whole \& genuine meanness of it, \& publish its meanness to the world; or if it were sublime, to know it by experience, \& be able to give a true account of it in my next excursion.'' -- {\sc Henry David Thoreau}, ``Where I Lived, \& What I Lived For'', Walden {\rm\cite{Thoreau2014}}
\end{quote}
{\it Tạm dịch}: Tôi vào rừng sâu vì tôi muốn sống có chủ ý, chỉ đối mặt với những sự thật thiết yếu của cuộc sống, \& liệu xem tôi có thể học được những gì nó [đời sống trong rừng sâu] dạy hay không, \& khi chết đi, tôi mới phát hiện ra rằng mình chưa sống. Tôi không muốn sống 1 cuộc sống không phải là cuộc sống, cuộc sống thật quý giá; tôi cũng không muốn thực hành sự cam chịu, trừ khi điều đó thực sự cần thiết. Tôi muốn sống sâu sắc \& hút hết tủy sống, nhựa sống của cuộc đời, sống thật kiên cường \& giống Spartan đến mức đánh tan tất cả những gì không phải là cuộc sống, cắt 1 đường rộng \& cạo sát, dồn cuộc đời vào 1 góc, \& giảm nó xuống mức thấp nhất, \& nếu nó được chứng minh là hèn hạ thì tại sao lại lấy toàn bộ ý nghĩa đích thực của nó \& công bố ý nghĩa đó của nó với thế giới; hoặc nếu nó là cao siêu, hãy biết nó bằng kinh nghiệm \& có thể tường thuật chân thực về nó trong chuyến du ngoạn tiếp theo của tôi.

FOMO vs. JOMO. Fear of missing out vs. Joy of missing out.

%------------------------------------------------------------------------------%

\section{Combustion: Order {\it\&}{\tt/}vs. Chaos -- Bùng Cháy: Trật tự {\it\&}{\tt/}vs. Hỗn loạn}
\noindent\textbf{\textsf{Resources -- Tài nguyên.}}
\begin{itemize}
	\item {\sc Jordan B. Peterson}. {\it 12 Rules for Life: An Antidote to Chaos} \cite{Peterson2018} với bản dịch {\it 12 Quy Luật Cuộc Đời: Thần Dược Cho Cuộc Sống Hiện Đại} \cite{Peterson2022a}
	\item {\sc Jordan B. Peterson}. {\it Beyond Order: 12 More Rules for Life} \cite{Peterson2021} với bản dịch {\it Vượt Trên Trật Tự: 12 Quy Tắc Cho Cuộc Sống} \cite{Peterson2022b}.
\end{itemize}

\subsection{The ultimate split? -- Sự chia tách cuối cùng?}
\begin{flushright}
	\href{https://www.youtube.com/watch?v=dIbeazAlxM4}{\small[{\sc Chopin} {\it Ballade No. 1 in G Minor}]}
\end{flushright}

\begin{quote}
	{\sf Unknown entity [?]}: Mày làm ơn ngưng ngay trò dối gạt bản thân mình đi. Mày coi lại tổng thể cái dự án nghiên cứu này đi.
	
	{\sf Hồng [25; depressed mathematics PhD student]}: Nhưng tao mang vai trò tiên phong, tao phải chứng tỏ thực lực của tao cho họ không còn khinh thường tao, màu da của tao, hay dân tao nữa.
	
	{\sf Unknown entity [?]}: Mày nghĩ ông thầy phụ mày để mày làm tiếp sau khi mày vạch trần ổng là kẻ phá hoại nhóm à? Chưa kể ông ta còn tới mức 1 psychopath in the workplace? Sao tới tận bây giờ mà mày vẫn ngây thơ quá vậy?
	
	{\sf Hồng [25; depressed mathematics PhD student]}: Nhưng không làm Toán tiếp thì tao biết làm gì nữa?
	
	{\sf Unknown entity [?]}: Bộ mày chỉ giỏi có món đó chắc? Đừng tự dối gạt bản thân mình nữa. Mày biết điều mày muốn nhất là gì mà. Just say it clearly. Say it.
	
	{\sf Hồng [25; depressed mathematics PhD student]}: Trong sâu thẳm thâm tâm của tao, cho tới tận lúc này, tao vẫn không chắc 1 tấm bằng Tiến sĩ Toán (Mathematics PhD), dẫu thuần túy hay ứng dụng đi chăng nữa (pure or applied), có thể đo được hết khả năng của tao. Tao muốn được công nhận theo cách riêng của tao trên nhiều lĩnh vực, trong địa hạt Toán học \& cả bên ngoài địa hạt Toán học. Tao muốn bất cứ thứ gì tao quan tâm \& đụng tới đều trở thành 1 dạng nghệ thuật cả, kiểu {\it creative artist} như nhà Toán học người Anh nổi tiếng {\sc Godfrey Harold Hardy} trong quyển {\it A Mathematician's Apology} \cite{Hardy1940,Hardy1992,Hardy2022} (tạm dịch: {\it Lời Xin Lỗi Của 1 Nhà Toán Học}) hay kiểu {\it dilettante}\footnote{\sf a person who does or studies something but is not serious about it \& does not have much knowledge; a person with a general but superficial interest in any art or a branch of knowledge.} như {\sc Blaise Pascal} được mô tả trong quyển {\it Pens\'ees} \cite{Pascal2003} (tạm dịch: {\it Suy tư{\tt/}nghĩ}) vậy.
	
	{\sf Hồng's honesty devil [?]}: Chuẩn. Từ giờ trở đi tự mày biết phải làm gì rồi đấy.
\end{quote}

\begin{quote}\it
	``Perfectly balanced. As all things should be.'' -- {\sc Thanos}, Avenger: End Game
\end{quote}



\subsection{Art \& life -- Nghệ thuật \& cuộc sống}
We need arts for life. An advise for artist by Sir {\sc Ian McKellen}:
\begin{quote}\it
	``Practice any art to experience becoming. To find out what's inside you \& make your soul grow'' -- Sir {\sc Ian McKellen} reads Kurt Vonnegut's inspiring letter to a group of school students.
\end{quote}

\begin{quote}\it
	``We don't read \& write poetry because it's cute. We read \& write poetry because we are members of the human race. \& the human race is filled with passion. \& medicine, law, business, engineering, these are noble pursuits \& necessary to sustain life. But poetry, beauty, romance, love, these are what we stay alive for. To quote from Whitman, ``O me! O life! $\ldots$ of the questions of these recurring; of the endless trains of the faithless $\ldots$ of cities filled with the foolish; what good amid these, O me, O life?'' Answer. That you are here - that life exists, \& identity; that the powerful play goes on \& you may contribute a verse. That the powerful play \emph{goes on} \& you may contribute a verse. What will your verse be?'' --  {\sc N. H. Kleinbaum}, \href{https://www.imdb.com/title/tt0097165}{\it Dead Poets Society} (1989)
\end{quote}

\begin{example}[Music \& Mathematics]
	Dr. {\sc Phạm Tuấn Huy} -- chàng trai vàng của làng Toán học Việt Nam trưởng thành với Âm nhạc.
\end{example}

%------------------------------------------------------------------------------%

\section{Towards the $\Psi$-Flow: Optimal Experiences -- Hướng Đến Dòng Chảy Tâm Lý Học: Các Trải Nghiệm Tối Ưu}
\noindent\textbf{\textsf{Resources -- Tài nguyên.}}
\begin{itemize}
	\item {\sc Mihaly Csikszentmihalyi}. {\it Flow: The Psychology of Optimal Experience} \cite{Csikszentmihalyi2008} với bản dịch {\it Dòng Chảy: Tâm Lý Học Hiện Đại Trải Nghiệm Tối Ưu}.
	\item {\sc Mihaly Csikszentmihalyi}. {\it Creativity: Flow \& the Psychology of Discovery \& Invention} \cite{Csikszentmihalyi2013}.
\end{itemize}

\begin{definition}[$\Psi$-flow]
	``\emph{Flow} in \href{https://en.wikipedia.org/wiki/Positive_psychology}{positive psychology}, also known colloquially as being \emph{in the zone} or \emph{locked in}, is the \href{https://en.wikipedia.org/wiki/Mental_state}{mental state} in which a person performing some activity is fully immersed in a feeling of energized focus, full involvement, \& enjoyment in the process of activity. In essence, flow is characterized by the complete absorption in what one does, \& a resulting transformation in one's sense of time. Flow is the melting together of action \& \href{https://en.wikipedia.org/wiki/Consciousness}{consciousness}; the state of finding a balance between a skill \& how challenging that task is. It requires a high level of concentration. Flow is used as a \href{https://en.wikipedia.org/wiki/Coping}{coping} skill for stress \& anxiety when productively pursuing a form of leisure  that matches one's skill set.'' -- \href{https://en.wikipedia.org/wiki/Flow_(psychology)}{Wikipedia{\tt/}flow (psychology)}
\end{definition}

\begin{dinhnghia}[Dòng chảy trong tâm lý học]
	\emph{Dòng chảy} trong tâm lý học tích cực, còn được gọi thông tục là \emph{ở trong vùng} hoặc \emph{bị nhốt}, là trạng thái tinh thần trong đó 1 người thực hiện 1 số hoạt động hoàn toàn đắm chìm trong cảm giác tập trung tràn đầy năng lượng, tham gia trọn vẹn, \& thích thú trong quá trình hoạt động. Về bản chất, dòng chảy được đặc trưng bởi sự tập trung hoàn toàn vào những gì 1 người làm \& dẫn đến sự biến đổi trong nhận thức về thời gian của 1 người. Dòng chảy là sự hòa tan của sự hành động \& ý thức; trạng thái tìm kiếm sự cân bằng giữa 1 kỹ năng \& mức độ thách thức của nhiệm vụ đó. Dòng chảy đỏi hỏi mức độ tập trung cao. Dòng chảy được sử dụng như 1 kỹ năng đối phó với căng thẳng \& lo lắng khi theo đuổi 1 hình thức giải trí phù hợp với kỹ năng của 1 người 1 cách hiệu quả. 
\end{dinhnghia}

\subsection{Passion -- Niềm đam mê}

%------------------------------------------------------------------------------%

\section{On Research -- Bàn Về Nghiên Cứu}
\label{sect: research}

\subsection{Dirty tricks -- Các thủ đoạn bẩn thỉu}
{\it Phạm vi áp dụng.} Các thủ đoạn dưới đây cũng áp dụng trong môi trường văn phòng, \& các môi trường làm việc trong phòng kín (why?).

\noindent\textbf{\textsf{Resources -- Tài nguyên.}}
\begin{itemize}
	\item \cite{Phipps_Gautreys_muu_hen_ke_ban_tap_1}. Mike Phipps, Colin Gautreys. {\it Mưu Hèn Kết Bẩn Nơi Công Sở. Tập 1: Nghệ Thuật Nhận Biết \& Phòng Tránh ``Tiểu Nhân'' Trong Công Việc}.
	\item \cite{muu_hen_ke_ban_tap_2}. Alpha Books. {\it Mưu Hèn Kết Bẩn Nơi Công Sở. Tập 2: Nghệ Thuật Thăng Tiến Trong Sự Nghiệp}.
\end{itemize}


\subsubsection{Pretend to borrow documents -- Giả vờ xin tài liệu}
1 văn phòng nghiên cứu viện Weierstrass, Berlin, Đức. Hắn thấy mình đang cặm cụi tính, chợt đồng nghiệp người Anh gốc Ấn của hắn hỏi mượn tài liệu. Hắn chả nghĩ nhiều, cứ gửi qua bản tiếng Đức, kèm luôn cả bản dịch tiếng Anh mà hắn tự soạn, thêm cả ghi chú cá nhân vào đó. Hắn chợt nhớ tới chị Thương, senpai trước hắn 2 khóa, lúc hắn học Master 2. Hắn may mắn được học bổng của 1 viện nghiên cứu Pháp, Henri Lebesgue centre de math\'ematiques\footnote{\url{https://www.lebesgue.fr/en}.}. Chị ta cũng hay xin đề thi \& tài liệu của hắn. Hắn chả nghĩ nhiều nên cho mượn tuốt. Mỗi lần đưa là chỉ sẽ nói đề khó quá, có khi chỉ làm không nổi, cười mỉm, rồi sau lưng hỏi điểm từng đứa. Hắn lúc đó tự hỏi:
\begin{quote}
	{\sf Hồng [22; mathematics Master 2 student]}: Sao chỉ làm nổi nhỉ? Mình thấy chuyện hốt giải nhất Olympic sinh viên toàn quốc còn dễ chịu hơn mấy cái đề master 2 này. Toàn mánh khóe, calculus tricks đủ kiểu. Phải cày liên tục suốt vài năm mới nhớ đủ trick để làm nổi. Trong khi chỉ còn chưa biết là hàm bình phương khả tích $L^2$, thậm chí $L^p$ với $p\in[1,\infty]$, thì được không cho giá trị tại 1 vài điểm, nói đúng hơn là 1 tập có độ đo không (sets with Lebesgue measure zero). Chắc chỉ chỉ đùa cho vui. Chắc thế.
\end{quote}
Tất cả tài liệu hắn đưa cho chị ta, chỉ đều phán: ``Chị thấy cũng dễ mà.'' Hắn khá rành kiểu này, nên chỉ cười mỉm rồi cho qua. Hắn chưa bao giờ muốn gây chuyện.
\begin{quote}
	{\sf Thương [25; 2nd year mathematics PhD student in numerical analysis, sucked at mathematics]}: Chị sẽ nói thầy của tụi mình để em khỏi làm luận văn luôn.
	
	{\sf Hồng [23; mathematics Master 2 student]} nghĩ trong đầu: Đ.M. cái con khốn học dốt tới mức không biết mình dốt mà hám danh, hám quyền. Tối ngày làm mấy cái chuyện nói xấu, đâm thọt sau lưng người khác mà tỏ ra tốt đẹp. Lủng background giải tích thua cả 1 đứa licence năm nhất Đại học mà tự cho mình cái quyền đì, quyền hành xác cái thằng tự kiếm học bổng Master làm từng khâu từ làm hồ sơ {\it \'etudes en France} tới phỏng vấn {\it Campus France}. Làm bẻ mặt dân Việt Nam trước cả lớp toàn mấy thằng Pháp hệ {\it École normale supérieure} (Paris)\footnote{\href{https://en.wikipedia.org/wiki/Ecole_normale_superieure_(Paris)}{Wikipedia{\tt/}École normale supérieure (Paris)}.} (abbr., ENS). Má nó cái con khốn. {\it Fucking stupid soulless talentless eyebrowless ass-licking controlling bitch}!
\end{quote}
Đoạn, hắn im lìm, mặt cúi gầm, không chịu ngẩng lên. Mà đúng ra hắn không thể ngẩng lên. Có cái gì đó đang chết dần chết mòn \& mục rữa bên trong hắn. Không phải khối ung thư kiểu tế bào sinh học, mà là 1 loại `ung thư' khác bên trong tâm trí hắn: {\it Không giống bệnh ung thư ở thể vật lý, sinh học tế bào, bệnh trầm cảm là 1 dạng ung thư về tinh thần \& trí óc}.
\begin{quote}
	{\sf Hồng [23; mathematics Master 2 student]}: \& anh có biết điều khiến tôi vừa sợ, vừa phát điên, \& ám ảnh hơn là gì không?
	
	{\sf Hồng [27--?; psychotherapist -- nhà trị liệu tâm lý]}: Anh cứ việc nói. Tôi sẽ lắng nghe.
	
	{\sf Hồng [23; mathematics Master 2 student]}: Đó là khi tôi nhìn xung quanh cái bữa tiệc ấy, đợi ai đó cản chị ta hay nhắc khéo để ngăn cái hành động ngu xuẩn của chị ta lại. Thì ai nấy cũng nhìn tôi như cái việc tôi đáng bị như thế là hiển nhiên nhất trên trần đời. A fucking crowd of bullies. Tôi có làm gì sai chứ? Tôi chỉ tốt với sai người, \& nhẫn nhịn với nhầm người thôi mà.
\end{quote}

\subsubsection{Verbal bullying -- Bắt nạt lời nói}

\begin{quote}
	{\sf Dương [26; 2nd year mathematics PhD student in mathematical analysis, sucked at coding \& numerics]}: Mày lấy cái cứt gì mà qua được tới đây hả thằng kia?
	
	Ai thèm thuê cái chó như mày? Mẹ chả có cái chó gì cả mà xin được học bổng.
\end{quote}
Hắn chợt hiểu ra cái chuyện hắn giấu nhẹm mấy giải Olympic toán sinh viên toàn quốc là sai. Trước những đứa dốt tới mức không biết họ dốt, hắn hoàn toàn bất lực. {\it Giữ hòa khí á? Chung tay góp sức xây dựng cộng đồng đoàn kết lành mạnh á?} Pure fucking bullshits.

Dần dần, hắn căm ghét những buổi ăn chung tưởng chừng như đoàn kết nhưng toàn mấy thủ đoạn tìm điểm yếu, vạch lá tìm sâu, tỏ vẻ thông minh thượng đẳng, rồi bợ đít xu nịnh nhau như 1 lũ ô hợp.
\begin{quote}
	{\sf Hồng [23; depressed mathematics Master 2 student]}: Trừ vài người tốt bụng \& lành tính như chị Châu, anh Tuấn, anh Hợi, chị Anh ra thì toàn 1 lũ bộ tịch, lúc nào cũng tỏ vẻ bộ tịch đầy giả tạo. Mãi như những đứa con nít chưa lớn nhưng giả vờ trưởng thành để gieo rắc cái ác thuần khiết -- cái thể loại ác không biết điểm dừng của mấy đứa con nít khi làm đau người khác \& tỏ ra ngày càng thích thú khi thấy người khác ngày càng đau \& đứa con nít đó cảm thấy ngày càng có quyền lực đàn áp mọi thứ trên đường đi của chúng. Hệt như trong phim \href{https://www.imdb.com/title/tt0213338/}{\it Cowboy Bebop} (1998--1999). 1 cái ác sơ khai, thuần khiết, non nớt không biết điểm dừng nếu không chịu bất cứ sự trừng phạt nào.
\end{quote}
-- hắn phán hệt như nhân vật chính không tên trong tiểu thuyết {\it The Catcher In The Rye} \cite{Salinger_catcher_in_rye} của {\sc Jerome David Salinger} với bản dịch tiếng Việt {\it Bắt Trẻ Đồng Xanh} \cite{Salinger_btdx}.

Từ 1 đứa cố gắng sống tốt, trở nên có ích để giúp đỡ mọi người, bất cứ ai cần giúp hắn đều sẵn sàng giúp, giờ đây hắn dần hắc hóa thành 1 kẻ khó chịu. Điều mà sau này hắn nhìn lại thì thấy đó là 1 lẽ tất yếu. Hắn phải giết chết hoặc ít nhất là hạn chế cái phần agreeable tới mức quá lành của hắn, bằng cách đặt ra ngưỡng tối đa -- maximum threshold -- thì hắn mới phát triển 1 cách lành mạnh được.

\subsubsection{Micro-envy \& micro-greedy -- Đố kỵ vi mô \& tham lam từng ly từng tý}

\begin{quote}
	{\sf Hồng [27--?; psychiatrist]}: Tại sao anh không tham gia ăn cùng với mọi người mà toàn ăn 1 mình? Những anh chị nghiên cứu sinh Tiến sĩ chắc hẳn sẽ phải đối tốt với học viên Thạc sĩ như anh mới phải lẽ chứ hả?
	
	{\sf Hồng [23; depressed mathematics Master 2 student]}: Họ hơn thua, ganh tỵ từng miếng ăn. Lúc nào tôi ăn họ cũng nhìn vào đĩa thức ăn của tôi cả. Ngắm xem miếng thịt của họ to hơn hoặc bằng miếng của tôi hay không. Nếu miếng thịt của họ to hơn thì họ tạm bằng lòng rồi khịt khịt mũi, chách lưỡi chành chạch, tằn hắn, liếc nhìn nhau tỏ vẻ khinh bỉ tôi. Nếu miếng thịt của họ nhỏ hơn tí xíu hoặc thực tế không hề nhỏ hơn nhưng nhận thức của họ khiến họ cảm thấy phần của họ nhỏ hơn của tôi thì họ phát tiết lên. Họ khiến tôi có cảm giác tôi không xứng đáng với đĩa thức ăn trước mặt, không xứng đáng ngồi chung với họ ở cái bàn ăn thiêng liêng chỉ dành cho những người đủ thông minh đó. Tôi mất hẳn vị giác.
	
	{\sf Hồng [27--?; psychiatrist]}: Khoan đã, ý anh nói ``miếng thịt'' là theo nghĩa bóng phải hông? Tức học bổng Master của anh?
	
	{\sf Hồng [23; depressed mathematics Master 2 student]}: Theo nghĩa đen, họ hơn thua cái miếng thịt mà nhà ăn sinh viên phát theo phần cho mỗi người.
	
	{\sf Hồng [27--?; psychiatrist]}: ``Họ'' bao nhiêu tuổi?
	
	{\sf Hồng [23; depressed mathematics Master 2 student]}: Hơn tôi 2--3 tuổi, nếu hơn 3 thì bằng tuổi chị ruột tôi. Nhưng tính tình chả khác gì những đứa con nít to xác với những cái mồm độc ác cả.
	
	{\sf Hồng [24; depressed mathematics PhD student]}: Trường hợp này cũng xảy ra với tôi. Gã người Pháp làm postdoc chung team nghiên cứu với thân hình to lớn lúc nào cũng liếc tôi cả. Mỗi lần tôi đến khu ăn chung với đồng nghiệp là hắn sẽ thét lên ``Dangerous here! Dangerous here!'' để cảnh báo tất cả đồng nghiệp chung phe với hắn hoặc những đồng nghiệp cũng cách ly hắn nhưng hắn ảo tưởng là họ chung phe bắt nạt tụi nghiên cứu sinh Tiến sĩ. Tôi nhận ra là những thành viên của đám bắt nạt cũng chả ưa gì nhau cả, toàn soi mói, chờ thời cơ hạ bệ lẫn nhau. Khi bị đồng nghiệp ngó lơ, bị bơ đi thì hắn sau đó trách tôi tại sao không hòa đồng, không trổ tài nấu ăn cho mọi người xem như cách hắn nướng cái bánh sinh nhật siêu khéo cho đồng nghiệp cùng thưởng thức.
	
	{\sf Hồng [27--?; psychiatrist]}: Gã người Pháp anh nói bao nhiêu tuổi?
	
	{\sf Hồng [24; depressed mathematics PhD student]}: Hơn 40 tuổi, có khi gần 50. Mà tính tình như 1 đứa con nít to xác, cũng với cái mồm ác độc.
	
	{\sf Hồng [27--?; psychiatrist]}: Thế 2 anh có nghĩ là sẽ có thể học hỏi gì từ những kẻ hơn thua miếng ăn theo nghĩa đen này không? Không, hơn thế nữa, 2 anh có nghĩ là có thể trưởng thành trong môi trường mà những kẻ này tỏ vẻ cầm trịch không?
	
	{\sf Hồng [23; depressed mathematics Master 2 student]} \& {\sf Hồng [24; depressed mathematics PhD student]}: Thế chẳng lẽ chúng tôi phải bỏ Toán mới trưởng thành được à? Chúng tôi hoàn toàn có thể hủy diệt cái sự vô học đó của bọn họ. Họ không giỏi tiếng Anh, phát âm bập bẹ không nghe rõ được, kỹ năng máy tính thì cùi bắp, kiến thức Toán của tôi lủng phần Master 1 do phân môn Giải Tích (Mathematical Analysis) đòi hỏi kiến thức nền (background) nhiều hơn hẳn các phân môn khác như Đại Số \& Thống Kê, nhưng họ thì còn tệ hơn: họ lủng cả phần Bachelor; chưa kể học đến năm 2, năm 3 nghiên cứu sinh Tiến sĩ mà không sửa được mấy cái lỗi \LaTeX\ đơn giản, trong khi chúng tôi tự học trong 1 thời gian ngắn cũng sửa trong vài phút là xong. Chúng tôi thừa sức để hủy diệt cái mồm vô học của bọn họ. Nhưng chả bao giờ chúng tôi chọn cách thô lỗ \& vô học đó cả. Đơn giản vì bản chất con người chúng tôi không phải thế. Nếu làm thế, chúng tôi sợ sẽ biến thành những người mà chúng tôi không ưa hoặc căm ghét mất.
	
	{\sf Hồng [27--?; psychiatrist]}: Ít nhất 2 anh đi tới nơi nào đó mà có thể thưởng thức bữa ăn 1 cách trọn vẹn, không bị soi mói kích cỡ của miếng thịt trên đĩa ăn của anh nữa, thì anh mới có cảm giác ăn ngon miệng lại. Trốn vào chốn rừng sâu để ăn beef steak như các video nấu ăn ngoài trời ở gần nơi có suối hoặc rừng rậm cũng không phải là 1 ý tồi. Bon apetit!
	
	{\sf Hồng [23; depressed mathematics Master 2 student]}: Thằng Dương còn mượn 1 tháng học bổng của tôi nữa, 1000 euro, để gom mua máy tính gì đấy. Hắn đổ thừa cái laptop của hắn. Nhưng muốn code chạy nhanh thì phải cải thiện thuật toán (algorithm) chứ, mà hắn có lương của nghiên cứu sinh mà, nhiều hơn học bổng của tôi, chắc 1700--1800 euro{\tt/}tháng trước thuế, sau thuế còn tầm 1400 euro{\tt/}tháng. Nhưng hắn bắt phải mượn của tụi Master bọn tôi.
	
	{\sf Hồng [27--?; psychiatrist]}: Nhưng tại sao anh lại cho người lạ mượn nhiều tiền đến thế?
	
	{\sf Hồng [23; depressed mathematics Master 2 student]}: Anh Tuấn học chung Master với tôi cũng hỏi y như thế. Mà tôi biết lý do rồi. 1 phần tôi muốn thân với cái nhóm các anh chị nghiên cứu sinh Tiến sĩ, có chị Phương giúp tôi các thủ tục tiếng Pháp lúc mới qua, nên tôi thấy giúp được gì giúp. Nhưng cái thằng chó này chả giúp gì tôi cả, mỗi lần đi chợ chung là hắn thó 1 vài món gì đấy từ túi đi chợ của tôi, mượn tiền mà cái giọng như cướp giật. Mà nguyên nhân chính á? Chắc là chị Thương có lần cảnh báo tôi là: ``Nếu em sang đây mà không chơi với ai là chết chắc đấy.'' Chơi thì chơi rồi đấy. Tặng quà, tặng đồ ăn đủ kiểu. Còn bắt phải làm cái đéo gì nữa? Cống nạp tài liệu, bóc lột công sức thì mới gọi là chơi hợp với đám bắt nạt à? Cái con khốn ngu lồn đấy. Cái thể loại gì đéo có chân mày. Xin lỗi tôi hơi xúc động, anh thông cảm.
	
	{\sf Hồng [27--?; psychiatrist]}: Không việc gì, anh cứ thành thật với cảm xúc bản thân.
	
	{\sf Hồng [23; depressed mathematics Master 2 student]}: Học thì dốt mà toàn hăm he dọa nạt tụi Master. Chả hiểu sao được sang đây học PhD. Ngay cả thầy chị ta cũng chả thể chịu nổi cái sự ngu dốt về mặt Toán học của chị ta. Chị ta than với tôi suốt là không hiểu sao chị ta ở cái tòa nhà nghiên cứu Toán lý thuyết (Pure Mathematics) đấy. Cái đó là chuyện tất yếu, chả có background thì sao làm việc ở level nghiên cứu nổi? Rồi cuối cùng thì anh biết đấy. Do không thể tự bản thân xử lý với cái mớ phức cảm tự ti (inferiority complex) mà trong tâm lý học Adler lấy làm 1 trong 2 tâm điểm nghiên cứu, con khốn \& thằng khốn này phóng chiếu (psychological projection) tất cả những gì tự bản thân biết mình tệ hại nhưng quá hèn nhát nên không dám tự mình thừa nhận, đối diện khuyết điểm mà đổ thẳng lên đầu tôi. Tôi có làm gì sai chứ? Chẳng lẽ tôi phải dâng hết học bổng mà tự tôi kiếm được cho đám đó thì tụi nó mới hết cảnh cáo tôi à? Chẳng lẽ tôi phải dùng đến vũ lực như hồi năm nhất Đại học do phải ở chung với 1 thằng siêu khốn nạn ngay lúc cha tôi vừa mới mất để rồi lại bị mấy ông thầy Đại học chớm lấy \& dùng làm điểm yếu lần nữa à? Lại mấy cái chuyện giảng đạo nhảm nhí về bản tính của mấy con vật gì đấy trên kênh Discovery Channel khám phá động vật trong thiên nhiên để rồi 1 đứa tốt tánh như tôi lại phải bị ví như 1 con rắn hay 1 con sói lang nham hiểm, còn con rắn, con sói thật sự thì lại trở thành con nai tơ ngơ ngác à? Pure fucking bullshit.
	
	{\sf Hồng [27--?; psychiatrist]}: Cái sai của anh là anh không thiết lập 1 ranh giới lành mạnh (healthy boundary).
\end{quote}

\begin{Rule}[On humanity development -- Bàn về phát triển nhân cách]
	Dẫu cho bạn làm bất cứ ngành nghề nào, đừng quên nhiệm vụ chính của việc làm người là phát triển nhân cách 1 cách toàn diện. Đừng phát triển nhân cách theo xu hướng của 1 kẻ khốn nạn, thích bắt nạt bất cứ ai mà bạn cho là dưới cơ hay yếu thế hơn bạn.
\end{Rule}

\subsubsection{Steal books, delete files -- Trộm sách, xóa tập tin}

\begin{question}
	Bạn sẽ làm gì khi phát hiện bạn bè, đồng nghiệp, thậm chí là người quản lý hoặc sếp nếu làm trong môi trường công sở hoặc chính thầy{\tt/}cô, người hướng dẫn, thậm chí giáo sư nếu làm trong mảng học thuật xóa bài hoặc hủy hoại các công trình của bạn?
\end{question}

\subsubsection{Divide to control groups -- Chia rẽ khiến lục đục nội bộ để dễ dàng kiểm soát}
Thay vì phương pháp {\it chia để trị -- divide to conquer} như trong mảng {\it Bất đẳng thức -- inequality} trong địa hạt của Toán học, phần này bàn về {\it divide to control \& micro-managing -- phương pháp chia để kiểm soát \& quản lý vi mô}.

\begin{question}
	Trong cuộc chạy đua vũ trang để mưu sinh \& để khẳng định thực lực bản thân, con người ta cần phải tốn hay từ bỏ 1 cách phí phạm bao nhiêu phần nhân tính tốt thuần khiết để có thể chiến thắng trên đấu trường vật chất, đấu trường vị thế, cùng nhiều thể loại đấu trường khác, bằng các phương pháp bạo lực, vô đạo đức, các thủ đoạn bẩn thỉu, hèn hạ, etc. để rồi lại thua cuộc trên đấu trường về nhân tính \& lương tâm?
\end{question}

\subsection{Standards -- Các tiêu chuẩn}
It is kind of funny, ironic, \& sarcastic that the 1st author of this writing is a dropout PhD student from one of the best research institutes of Applied Mathematics in Germany, Europe. Anyhow, it is also a good idea to see from the outside. The perspective of an outsider sometimes may shine some light \& reveal some insight to a dark room.

Khá là hài hước, mỉa mai, \& châm biến khi mà tác giả đầu tiên của bài viết này lại là 1 nghiên cứu sinh bậc Tiến sĩ bỏ cuộc từ 1 trong các viện nghiên cứu tốt nhất về Toán ứng dụng của Đức, châu Âu. Dù gì đi nữa, cũng sẽ là 1 ý hay khi mà nhìn mọi thứ từ bên ngoài (từ tầm mắt của 1 đứa bỏ học, nghỉ việc ngang). Tầm nhìn của 1 kẻ ngoại lai đôi khi lại có thể chiếu vài tia sáng \& tiết lộ vài cái nhìn sâu sắc hay sự hiểu biết vào 1 căn phòng tối tăm. Đầu tiên chúng ta cần 1 khẩu hiệu cho cuộc diễu hành qua các dãy văn phòng u ám, hôi hám, \& tối tăm.

\vspace{2mm}
\noindent\textsf{\textbf{Slogan -- Khẩu hiệu}}: {\it A single bad publication will lead to endless public humiliations. -- Chỉ cần 1 bài báo, 1 công bố tồi cũng sẽ dẫn đến các cuộc công kích làm nhục công khai không có hồi kết}.

\begin{quote}
	{\sf Thọ [26; 3rd year theoretical mathematics PhD student]}: Anh muốn làm nghiên cứu khoa học ``thực chất'' á.
\end{quote}
Ý ám chỉ các xuất bản mà hắn đã làm là bẩn, mà thực ra tự bản thân hắn thấy nó không đủ chuẩn như hắn mong muốn.

Hắn cũng chả hiểu anh ta cho lắm, chỉ biết nhiều anh chị lớn tuổi hơn than phiền việc bị anh này nhìn đểu. Nhìn đểu là sao cơ chứ? Thôi, không liên quan tới hắn, nên hắn cứ kệ. Mà cái éo gì bất cứ course nào hắn học, thì anh ta đều bảo là dễ cả. Dễ á? Hắn cày muốn bụp mắt mà chưa thấy có cửa cạnh tranh với mấy anh chị ENS Paris chung lớp để điểm hắn không bị đôn xuống quá đáng. Hắn cũng để ý là ông anh này lúc nào cũng trốn ở 1 góc nào đó để rình rập \& quan sát hắn cả. \textit{Tại sao cái con người lúc nào cũng tỏ ra là 1 đàn anh đáng kính với lứa đàn em đi sau, đồng thời là 1 người chồng mẫu mực, thương yêu quan tâm vợ hết mực, dẫu công việc nghiên cứu có gian khổ thế nào đi chăng nữa, lại có đặc tính của 1 kẻ rình rập (stalker) mờ ám \& nham hiểm?}

\subsection{Philosophical methodologies}
Khi phải đối đầu với những thứ thật sự khó nhằn, hoàn toàn nằm ngoài hiểu biết hiện tại của 1 cá nhân, thì 1 cách khá đơn giản là bám víu vào những thứ đã biết rõ, dù có thể lặp đi lặp lại 1 cách đơn điệu \& nhàm chán, nhưng lại có trật tự để cân bằng với hỗn loạn -- tượng trưng cho những điều chưa biết \cite{Peterson2018,Peterson2022a,Peterson2022b}.

\begin{example}
	Dạy học bậc phổ thông trở xuống thì ``nhàn'', theo nghĩa là không cần phải nạp quá nhiều kiến thức mới, nhưng phải chú trọng về phương pháp dạy \& truyền đạt kiến thức 1 cách hiệu quả tới các học sinh. Nếu học sinh giỏi, tiếp thu nhanh thì khỏe. Gặp học sinh dốt hoặc đầu gấu thì mệt, đâm ra chán chường, cảm thấy phí phạm thời gian \& nguồn sức lực hạn chế của bản thân.
	
	Nghiên cứu thì lại khác. Trách nhiệm của nghiên cứu là phải đọc thật nhiều, nạp thật nhiều kiến thức để trau dồi bản thân mỗi ngày.***
\end{example}
[...]

Tạm phân loại học giả, theo ý cá nhân (sẽ bổ sung thêm):
\begin{itemize}
	\item Học giả làm các mảng, lĩnh vực năng động, với năng suất xuất bản ấn phẩm khoa học cao, thường được trích dẫn nhiều nhờ sự năng động của cộng đồng khoa học tương ứng.
	\item Học giả làm các mảng khó nhằn, trừu tượng, nên tần suất xuất bản ấn phẩm khoa học khá thấp, nhưng các bài này đều ở dạng nặng đô (hardcore), thường ít được trích dẫn vì kén độc giả. Nếu bài báo đó trở thành cornerstone thì lại được trích dẫn nhiều đến rất nhiều, na ná dạng benchmark cases for industrial purposes của loại 1 (data mẫu chuẩn để các người làm nghiên cứu R\&D ở các lĩnh vực công nghiệp dùng).
\end{itemize}
Ưu điểm của loại 1 là đi hội nghị thường xuyên. Mà đa số mấy hội nghị này giàu do dính đến công nghiệp hoặc dịch vụ số hóa (Artificial Intelligence{\tt/}Deep Learning{\tt/}Machine Learning) nên chắc đồ ăn nhiều \& ngon, ít nhất cũng ăn đứt mấy bữa tiệc giản đơn gồm trà, cafe máy cùng vài cái bánh quy như các hội nghị toán lý thuyết ở Pháp mà hồi mình học Master (hay chỉ có mấy chỗ nằm ở rìa của Pháp là vậy nhỉ?). Mà thực ra lúc mấy giáo sư Toán thảo luận với nhau, thay vì nhắm nháp cafe \& ăn bánh quy, vài người lại say xưa thảo luận mà ăn (nhầm?) phấn trắng.

Chắc mình thuộc loại 2, hoặc ít nhất là mình tự ép bản thân thuộc loại 2 (nên gọi là {\it giả học -- fake scholar} thì hợp hơn). Trong khi loại 1 thì tạo cảm giác năng động, tràn trề của sức trẻ, thì loại 2 hoàn toàn ngược lại, mà phần lớn là phải cày background khá nhiều \& nặng, \& 1 trong những cái mệt nhất nhưng rewarding nhất của loại 2 là làm các công trình khoa học liên ngành, kết nối các kết quả mạnh nhất của các lĩnh vực lý thuyết với nhau.

Có 1 bài viết phân loại học giả hay của GS. Nguyễn Tiến Zũng của ĐH Toulouse. Tiếc là sau khi GS Zũng hỗn chiến với bác  Phùng Xuân Nhạ thì website cá nhân \url{http://zung.zetamu.net/} của GS trước bị lỗi font \& giờ có lẽ đã bay màu.

\subsection{Trends \& choices -- Các xu hướng \& lựa chọn}	

\begin{quote}
	{\sf Nhân [23]}: Thế anh có biết những sở thích thời học sinh của 1 người ảnh hưởng thế nào đến xu hướng các lựa chọn chuyên ngành trong tương lai của họ không?
	
	{\sf Hồng [28]}: Tôi không rõ lắm. Cụ thể sao?
	
	{\sf Nhân [23]}: Tui sẽ lấy ví dụ về ngành Toán. Vì nó là cái duy nhất tui rành, ít hơn là rành hơn ối thứ còn lại.
	
	Những học sinh thích giải phương trình, hệ phương trình ở Toán Sơ Cấp nhưng không thích Tin học thường sẽ có xu hướng chọn các ngành lý thuyết trừu tượng, như Đại Số, Hình Học Đại Số. 
	
	Những người thích bất đẳng thức ở Toán Sơ Cấp thường sẽ có xu hướng chọn hướng Giải tích, đặc biệt là hướng Phương Trình Vi Phân Đạo Hàm Riêng (Partial Differential Equations, abbr., PDEs) vì hướng này chủ yếu đánh giá (estimation), chặn (bound), i.e., các bất đẳng thức giữa các không gian hàm. Như vậy, xu hướng thích đánh giá các đại lượng liên quan tới các hàm sơ cấp ở Toán Sơ Cấp thường sẽ phát triển thành niềm đam mê việc đánh giá các đại lượng liên quan đến hàm hoặc các đối tượng toán học trừu tượng hơn.

	1 câu hỏi điển hình của các nhà Giải tích học (mathematical analysts) khi thảo luận các vấn đề toán học liên quan đến PDEs là:
	
	- {\it Do you think it is smooth (or regular) enough? -- Anh nghĩ nó có đủ trơn (hay nhớt) không?}
	
	- {\it It seems a little rough at the initial phase. But it will be smoother later. Oh, now it's already smooth enough for us. Let's do{\tt/}play with it. -- Nhìn có vẻ hơi thô trong giai đoạn đầu (màn dạo đầu?). Nhưng rồi nó sẽ trơn hơn thôi. Ồ nhìn này, nó đủ trơn rồi kìa. Nào, chúng ta cùng xử{\tt/}quất{\tt/}chơi nó (vấn đề giải tích này) thôi}.
	
	Hàm đối tượng trơn chưa đủ, để đặt tốt 1 bài toán, miền xác định, i.e., nơi hàm đó sống, phải đủ trơn nữa, tức là cái mép (boundary $\Gamma\coloneqq\partial\Omega$) của cái miền $\Omega$ phải đủ trơn để xài các công thức tích phân từng phần (integration by parts formulas or Green's identities) để tạo ra dạng yếu (weak formulation or variational formulation). Những miền quá thô, e.g., có các góc nhọn (rough boundaries with corners), kỳ dị (singularities), chỗ nhọn dễ bị đâm (cusps), có nhiều lỗ (holes) hoặc gai (thorns) sẽ không thích hợp để làm chỗ chơi đối với các nghiệm trơn, dẫu mấy cái nghiệm đó có trơn chùi cỡ nào đi chăng nữa, vẫn không đảm bảo an toàn để chơi với chúng. Safety 1st.
	
	Ngoài lề, dù hay thắc mắc với việc đòi hỏi các nghiệm trơn, nghiệm nhớt của phương trình vi phân đạo hàm riêng có đủ trơn, đủ nhớt hay không để mà có thể vô tư chơi với chúng, tuổi thơ của các nhà giải tích cho thấy họ không có liên quan đến bất kỳ về tình dục sớm kiểu con nít quỷ hoặc sống thử, hay lạm dụng tình dục nào cả. Cho nên việc đề xuất những khẳng định kiểu như của {\sc Sigmund Freud}, e.g., các nhà toán học loay hoay với câu hỏi đủ trơn thường có tuổi thơ liên quan đến các vấn đề tình dục sớm do cha mẹ hoặc người tình của họ không quan hệ kín đáo để cho con cái vô tình bắt gặp hoặc các sang chấn tâm lý do chịu lạm dụng tình dục từ sớm; hoặc lý luận kiểu {\sc Malcolm Gladwell} trong quyển {\it Outliers: The Story of Success} \cite{Gladwell2008} hay bản dịch {\it Những Kẻ Xuất Chúng: Cái Nhìn Mới Lạ Về Nguồn Gốc Của Thành Công} \cite{Gladwell_outlier} ngụ ý việc tiếp xúc 1 cách vô thức với các từ gợi hình (gợi dục) tác động đến tiềm thức sâu bên dưới ý thức dẫn đến xu hướng chỉ thích làm với các đối tượng đủ trơn hoặc cuồng với các khái niệm đủ nhớt, etc. là hoàn toàn không có sơ sở.
	
	{\sf Hồng [28]}: What is so wrong with you?
\end{quote}

\subsection{Signs -- Các dấu hiệu}

\subsubsection{Personal systems of notations, abbreviations, \& conventions}
Bộ (tuple), tập hợp (set), hay hệ thống các ký hiệu, cách viết tắt, \& các quy ước cá nhân -- a personal set{\tt/}system of notations, abbreviations, \& conventions -- của 1 nhà khoa học tự nhiên thiên về lý thuyết hơn là về tính toán engineering thuần ứng dụng, e.g., nhà toán học (mathematicians), nhà vật lý (physicists), nhà khoa học máy tính (computer scientist), etc. là dấu hiệu đầu tiên cho biết trình độ của họ. Đơn giản vì các môn khoa học này có 1 đặc thù là đòi hỏi độ nhất quán (consistency) cực kỳ cao cho nên 1 hệ thống ký hiệu nhất quán, không mâu thuẫn, tiện dụng, không tạo ra bất kỳ sự mơ hồ, mờ mập (confusion) sẽ phản ánh phần nào trình độ của họ. Đấy là dấu hiệu dễ nhận biết đầu tiên -- nhưng còn xa so với mức phán xét -- của 1 người làm khoa học giỏi hoặc ít nhất là có 1 người thầy, người hướng dẫn giỏi.

Riêng các nhà hóa học (chemists) thì có lẽ họ được quy định chung bởi các danh pháp quốc tế như International Union of Pure \& Applied Chemistry (abbr., IUPAC)\footnote{\url{https://en.wikipedia.org/wiki/International_Union_of_Pure_and_Applied_Chemistry}.} nên không{\tt/}chưa thể dùng hệ thống ký hiệu cá nhân để đánh giá sơ bộ. Có lẽ mình nên kết thêm vài đứa bạn chuyên ngành Hóa để hiểu thêm (vừa đủ).

Thus, a good advice for young science students: Build, polish, \& perfect endlessly your personal system of notations \& conventions so well that it will fit perfectly to any of, or at least most of, your research fields. Then you can effortlessly attack each of them, connect them, play with the interaction between them \& beyond, \& even foresee the hidden structure in the realm of abstractness.

Lời khuyên (tự thân) này na ná câu trích dẫn sau của Abraham Lincoln về việc đầu tư khâu chuẩn bị kỹ lưỡng:
\begin{quote}\it
	``Give me 6 hours to chop down a tree \& I will spend the 1st 4 sharpening the axe.'' -- {\sc Abraham Lincoln} (1809--1865) -- 16th President of the United States (1861--1865)
\end{quote}

\subsubsection{Consistency -- Sự nhất quán}

\begin{question}
	Liệu có nên (dấn thân) theo 1 nghề cố định, không chịu{\tt/}thèm nhảy nghề không? 
\end{question}

\begin{figure}[H]
	\centering
	\includegraphics[scale=.15]{high_salary}
	\caption{Credit: \href{https://www.linkedin.com/posts/judy-soloai_i-didnt-learn-this-in-school-i-learned-activity-7089014937494179840-jTze/}{Linkedin{\tt/}Judy Soloai{\tt/}I didn't learn this in school}.}
\end{figure}
Hiển nhiên 1 câu hỏi khó muôn thở. Khó chịu lẫn khó nhằn theo nhiều nghĩa. Nghĩa thứ nhất là nó không rõ ràng, \& sự không rõ ràng đến từ việc bản thân nó phụ thuộc vào quá nhiều yếu tố không thể xác định hết như các yếu tố về phương diện vật chất, e.g., lương, tài chính; cũng như các yếu tố về phương diện tinh thần, e.g., ý nghĩa công việc, cân bằng công việc--cuộc sống (work-life balance), sự phát triển cá nhân, cùng sự tương tác lẫn nhau giữa các yếu tố trong 2 phương diện đó; \& nếu lùi xa hơn nữa về quá khứ thì chúng cũng phụ thuộc vào nhiều yếu tố ban đầu của 1 cá nhân như điểm xuất phát mà bố mẹ mang lại, hoàn cảnh như khả năng tài chính của gia đình, sự ủng hộ từ dòng họ, \& ảnh hưởng của các mâu thuẫn, xung đột, lục đục nội bộ trong 2 môi trường nền tảng đó.

Tôi không hề nghĩ sẽ cố trả lời 1 cách hoàn hảo câu hỏi này hay giải quyết vấn đề này. Đồng ý là tôi ngu, nhưng chưa ngu đến mức vậy. Chưa kể có bất kỳ câu trả lời nào không (no guarantee of existence), nếu có thì cũng không hề có câu trả lời duy nhất (even if the existence is assumed, the nonuniqueness is still valid), cũng như chưa \& sẽ không chả có câu trả lời nào sẽ thỏa mãn hết tất cả các phương diện giá trị được suy xét ở {\it biểu diễn phân hoạch các giá trị \& ý nghĩa của cuộc đời} ({\it a decomposition of values \& meanings in life}) sẽ được xét đến trong Sect. \ref{sect: bullshit theory on live}: {\it A Bullshit Theory on Living -- 1 Lý Thuyết Nhảm Nhí Về Việc Sống}.


\subsubsection{Accuracy{\tt/}Precision -- Tính chính xác}

\subsubsection{Simplicity -- Sự giản đơn}
We use simplicity to fight difficulty. We do not add more unnecessary complexities \& redundancies to the war because if we do so, we will have to fight ourselves, our entanglements.

\subsubsection{Minimality -- Sự tối giản}
\noindent\textbf{\textsf{Resources -- Tài nguyên.}}
\begin{itemize}
	\item {\sc Chi, Nguyễn -- The Present Writer}. {\it Một Cuốn Sách về Chủ Nghĩa Tối Giản} \cite{Chi2022}.
\end{itemize}

\begin{quote}
	{\sf Nhân [26--?; Jeet Kune Do practitioner -- môn sinh Tiệt Quyền Đạo]}: Ngôi sao võ thuật {\sc Bruce Lee -- Lý Tiểu Long} (tên khai sinh: {\sc Lý Chấn Phiên}) rất giỏi trong việc đưa ra cú đấm \& cú đá chuẩn xác nhất, hiệu quả nhất, với lực mạnh nhất. {\sc Lý Tiểu Long} học đủ thứ loại võ \& chọn lọc những tinh chất tốt nhất của mỗi loại võ cho riêng môn võ của ông, i.e., Tiệt Quyền Đạo (Jeet Kune Do), tinh giản đến mức tối đa. Chẳng hạn, {\sc Lý Tiểu Long} cho rằng cú đấm xoáy của Karate phí quá nhiều lực do lợi dụng việc xoáy cả cánh tay trong khi trực quyền, i.e., cú đấm thẳng, thì chuyển gần như toàn bộ lực của người ra đòn vào diện tích tiếp xúc (contact area).
	
	{\sf Hồng [25; fluid mechanics researcher -- nghiên cứu sinh Cơ Học Chất Lỏng ]}: Bản chất của điều này là do lãng phí lực phát quyền{\tt/}cước do phải truyền 1 phần vào việc xoáy (rotation) $\boldsymbol{\omega} = \operatorname{curl}{\bf v} = \operatorname{rot}{\bf v} = \nabla\times{\bf v}$ với curl\footnote{\href{https://en.wikipedia.org/wiki/Curl_(mathematics)}{Wikipedia{\tt/}curl (mathematics)}.} là toán tử vector diễn tả độ xoay hay xoáy, còn ký hiệu là rot (i.e., rotation), còn ${\bf v}$ là tốc độ của cú đấm hoặc đòn đá.
	
	{\sf Nhân [26--?; Jeet Kune Do practitioner]}: Đúng đấy, thì ra có thể hiểu được nguyên lý cú đấm của {\sc Lý Tiểu Long} theo cách nhìn của Cơ Học Chất Lỏng của Vật Lý.
	
	{\sf Hồng [25; fluid mechanics researcher]}: Nếu tôi muốn hiểu hơn về môn võ, hay đúng hơn là trường phái võ thuật này (martial art school) thì nên tìm đọc quyển nào?
	
	{\sf Nhân [26--?; Jeet Kune Do practitioner]}: Quyển {\it Tao of Jeet Kune Do} \cite{Lee2011} (tạm dịch: {\it Đạo của Tiệt Quyền Đạo}) của chính {\sc Bruce Lee} viết, tôi chưa nghiệm được tí nào.
	
	{\sf Hồng [26--?; boxing practitioner -- môn sinh đấm bốc]}: Anh có thể tham khảo thêm việc cuộn người (body roll) để di chuyển trọng tâm cơ thể (gravity of body) để điều khiển dòng năng lượng (energy flow) vào cú đấm trong môn Quyền Anh. Nó khá là giống với bài Khởi Động cuộn 2 cánh tay di chuyển như làn sóng trong Tiệt Quyền Đạo.
	
	{\sf Hồng [25; fluid mechanics researcher]}: Mục đích của các bài tập chuyển động cơ thể này là gì?
	
	{\sf Nhân [26--?; Jeet Kune Do practitioner] \& Hồng [27; boxing practitioner]}: Dò tìm dòng chảy khí lực \& dòng chảy năng lượng trong cơ thể bằng độ nhạy cảm của cơ bắp \& sự tập trung tinh tế của trí óc để dồn chúng vào quyền trong môn Quyền Anh \& quyền hoặc cước trong môn Tiệt Quyền Đạo 1 cách tối ưu nhất, i.e., tiết kiệm hao phí lực, tận dụng được tối ưu chuyển động của cơ thể, độ xoay bàn chân \& cả thân người, cắt tất cả các động tác thừa thãi khiến hao phí lực không cần thiết, khiến cho gần như toàn bộ năng lượng được dồn vào đòn đánh.
	
	{\sf Hồng [25; fluid mechanics researcher]}: {\sc Lý Tiểu Long} có học Cơ Học Chất Lỏng hay không mà tự biết phải cực tiểu hóa hàm tạo bởi toán tử xoáy
	\begin{equation}
		\label{curl cost functional}
		\tag{$J_\omega$}
		J_\omega(Q_T,{\bf v}(Q_T))\coloneqq\int_0^T\int_{\Omega_t} |\nabla\times{\bf v}(t,{\bf x})|^2\,{\rm d}{\bf x}\,{\rm d}t
	\end{equation}
	với $T$ là tổng thời gian luyện tập hoặc cận chiến còn $\Omega_t$ là vùng mà người ra đòn đang chiếm lĩnh trong không gian vật lý 3 chiều $\mathbb{R}^3$ tại thời điểm $t\in[0,T]$ \& thay đổi theo thời gian $t$?
	
	{\sf Nhân [26--?; Jeet Kune Do practitioner] \& Hồng [26--?; boxing practitioner]}: Chà, ông ta là thiên tài võ thuật mà. Cần chi phải học mấy cái đó. {\sc Lý Tiểu Long} chỉ cần cảm nhận theo bản năng tự nhiên của ông ta thôi. Chỉ cần cái bộ cảm biến cơ bắp siêu nhạy cảm (highly sensitive muscle system) cùng 1 trí tuệ kiểu tư duy phản biện triết học (philosophical critical thinking) đã là quá đủ cho sự ra đời của trường phái võ thuật Jeet Kune Do. Dù là nghiên cứu sinh học thuật, nhưng anh cũng nên thử qua vài món võ thuật như Jeet Kune Do hoặc Boxing, biết đâu lại cần.
\end{quote}

\subsubsection{Vigor -- Khí lực, sức mãnh liệt}

\subsubsection{Rigour -- Tính chặt chẽ}

\subsubsection{Visionary -- Nhìn xa trông rộng}

%------------------------------------------------------------------------------%

\section{Dark Psychology $\boldsymbol{\Psi}$ -- Tâm Lý Học Hắc Ám $\boldsymbol{\Psi}$}

\begin{flushright}\it
	``It takes a wolf to catch a wolf.'' -- Training Day
\end{flushright}
\& it takes a wise man to catch a wolf in sheep-clothing (?).

\noindent\textbf{\textsf{Resources -- Tài nguyên.}}
\begin{itemize}
	\item {\it In Sheep's Clothing: Understanding \& Dealing with Manipulative People} \cite{Simon2010} với bản dịch tiếng Việt {\it Sói Đội Lốt Cừu: Kẻ Hiếu Chiến Ngầm \& Các Thủ Thuật Thao Túng Tâm Lý}.
	\item {\it Character Disturbance: The Phenomenon of Our Age} \cite{Simon2011} của tác giả {\rm George Simon Jr.}.
\end{itemize}


\subsection{Psychological manipulation 101 -- Thao túng tâm lý cơ bản}
\begin{itemize}\small\sf
	\item {\it101} [a] relating to a university course that is intended as an introduction to a subject for people who have never studied it before; relating to the basic facts in a particular field or subject.
\end{itemize}
Điển hình 1 người học Toán, đặc biệt là Đại số (Algebra) có thể thành thạo các thao tác đại số (algebraic manipulations) như cộng $+$, trừ $-$, nhân $\cdot$, chia $:$, etc., nhưng liệu người đó có chuẩn bị đủ kiến thức về mặt tâm lý để đối đầu \& phòng chống với 1 dạng manipulation khác ở 1 địa hạt khác xa hoàn toàn với địa hạt của Khoa học nói chung \& của Toán học nói riêng, khi người đó trưởng thành \& bước vào môi trường làm việc không khi mà dạng manipulation ở địa hạt Tâm lý học không hề thuần logic \& có quy tắc, trật tự lớp lang như manipulation của Đại số nữa: {\it A pure form of Chaos Devil}.

\begin{question}
	What is ``psychological manipulation''? Why? \& how to detect \& conquer it? -- Thao túng tâm lý là gì? Tại sao? \& làm cách nào để phát hiện cũng như chống lại nó?
\end{question}
Kẻ thao túng tâm lý là kẻ tiểu nhân, nham hiểm, đê tiện. Tùy vào mức độ thao túng tâm lý thành thạo đến mức nào, mà hậu quả tâm lý, sức khỏe có thể nghiêm trọng đến ngưỡng tương đương, tùy vào sức chống chọi tâm lý \& nhận thức của nạn nhân. Các người sống sót (tạm gọi là {\it survivor}) sau các cuộc thao túng tâm lý nặng có thể trở thành những kẻ mạnh về tâm lý \& chia sẻ cho những người đang bị thao túng tâm lý khác để có thể trở thành những người sống sót mạnh mẽ như chính họ đã từng trải qua.

%************

\subsubsection{\texttt{{\color{RubineRed}\#include} ``basic terminologies''}}
I have not read yet, \& unintended to, all the heavy references mentioned in this section, but they seem necessary to be included for the sake of completeness.

\begin{definition}[Psychological manipulation \cite{Psychological Wiki/Psychological manipulation,Wikipedia/Pyschological manipulation}]
	\emph{Psychological manipulation} is a type of \href{https://psychology.wikia.org/wiki/Social_influence}{social influence} that aims to change the \href{https://en.wikipedia.org/wiki/Perception}{perception} or behavior of others through underhanded, \href{https://en.wikipedia.org/wiki/Deceptive}{deceptive}, or even \href{https://psychology.wikia.org/wiki/Abuse}{abusive} tactics (\cite{Braiker2004}). By advancing only the interests of the manipulator, often at the other's expense, such methods could be considered exploitative, abusive, devious, \& deceptive.
\end{definition}

\begin{definition}[Psychological abuse{\tt/}violence, emotional{\tt/}mental abuse \cite{Psychology Wiki/Abusive relationship,Wikipedia/Psychological abuse}]
	\emph{Psychological abuse}, often called \emph{emotional abuse}, is a form of \href{https://en.wikipedia.org/wiki/Abuse}{abuse}, characterized by a person subjecting or exposing another person to behavior that may result in \href{https://en.wikipedia.org/wiki/Psychological_trauma}{psychological trauma}, including \href{https://en.wikipedia.org/wiki/Anxiety_disorder}{anxiety}, \href{https://en.wikipedia.org/wiki/Chronic_depression}{chronic depression}, or \href{https://en.wikipedia.org/wiki/Post-traumatic_stress_disorder}{post-traumatic stress disorder} (\cite{Dutton1994,Dutton_Goodman_Bennett2000,Thompson_Kaplan1996}). It is often associated with situations of \href{https://en.wikipedia.org/wiki/Abusive_power_and_control}{power imbalance in abusive relationships}, \& may include \href{https://en.wikipedia.org/wiki/Bullying}{bullying}, \href{https://en.wikipedia.org/wiki/Gaslighting}{gaslighting}, \& \href{https://en.wikipedia.org/wiki/Workplace_bullying}{abuse in the workplace} (\cite{Dutton_Goodman_Bennett2000,Thompson_Kaplan1996}). It also may be perpetrated by persons conducting \href{https://en.wikipedia.org/wiki/Torture}{torture}, other \href{https://en.wikipedia.org/wiki/Violence}{violence}, acute or prolonged \href{https://en.wikipedia.org/wiki/Human_rights_abuse}{human rights abuse}, particularly without legal redress such as \href{https://en.wikipedia.org/wiki/Detention_without_trial}{detention without trial}, \href{https://en.wikipedia.org/wiki/False_accusation}{false accusations}, false convictions \& extreme \href{https://en.wikipedia.org/wiki/Defamation}{defamation} such as where perpetrated by state \& media. 
\end{definition}

\begin{definition}[Breaking point (psychology) \cite{Wikepedia/Breaking point (psychology)}]
	In human psychology, the \emph{breaking point} is a moment of \href{https://en.wikipedia.org/wiki/Stress_(medicine)}{stress} in which a person breaks down or a situation becomes critical.
\end{definition}
``The intensity of environmental stress necessary to bring this about varies from individual to individual.''

\begin{definition}[Social isolation \cite{Wikipedia/Social isolation}]
	\emph{Social isolation} is a state of complete or near-complete lack of contact between an individual \& \href{https://en.wikipedia.org/wiki/Society}{society}. It differs from \href{https://en.wikipedia.org/wiki/Loneliness}{loneliness}, which reflects temporary \& involuntary lack of contact with other humans in the world. Social isolation can be an issue for individuals of any age, though symptoms may differ by age group.
\end{definition}
``Social isolation has similar characteristics in both temporary instances \& for those with a historical lifelong isolation cycle. All types of social isolation can include staying home for lengthy periods of time, having no communication with family, acquaintances or friends, and{\tt/}or willfully avoiding any contact with other humans when those opportunities do arise.''

See also, e.g. \cite{NASEM2020,NhuTrang2020}.

\begin{definition}[Anxiety disorder \cite{Wikipedia/Anxiety disorder}]
	\emph{Anxiety disorders} are a group of \href{https://en.wikipedia.org/wiki/Mental_disorder}{mental disorders} characterized by significant feelings of \href{https://en.wikipedia.org/wiki/Anxiety_(mood)}{anxiety} \& \href{https://en.wikipedia.org/wiki/Fear}{fear}. Anxiety is a worry about future events, while fear is a reaction to current events. These feelings may cause physical symptoms, such as increased heart rate \& shakiness. There are several anxiety disorders, including \href{https://en.wikipedia.org/wiki/Generalized_anxiety_disorder}{generalized anxiety disorder}, \href{https://en.wikipedia.org/wiki/Specific_phobia}{specific phobia}, \href{https://en.wikipedia.org/wiki/Social_anxiety_disorder}{social anxiety disorder}, \href{https://en.wikipedia.org/wiki/Separation_anxiety_disorder}{separation anxiety disorder}, \href{https://en.wikipedia.org/wiki/Agoraphobia}{agoraphobia}, \href{https://en.wikipedia.org/wiki/Panic_disorder}{panic disorder}, \& \href{https://en.wikipedia.org/wiki/Selective_mutism}{selective mutism}. The disorder differs by what results in the symptoms. An individual may have more than one anxiety disorder.
\end{definition}
See the main reference: \cite{APA2013}.

\begin{definition}[Depression (mood) \cite{Wikipedia/Depression (mood)}]
	\emph{Depression} is a state of low \href{https://en.wikipedia.org/wiki/Mood_(psychology)}{mood} \& aversion to activity. It can affect a person's thoughts, behavior, motivation, \href{https://en.wikipedia.org/wiki/Feeling}{feelings}, \& \href{https://en.wikipedia.org/wiki/Subjective_well-being}{sense of well-being}. The core symptom of depression is said to be \href{https://en.wikipedia.org/wiki/Anhedonia}{anhedonia}, which refers to loss of interest or a loss of feeling of pleasure in certain activities that usually bring joy to people. Depressed mood is a symptom of some \href{https://en.wikipedia.org/wiki/Mood_disorders}{mood disorders} such as \href{https://en.wikipedia.org/wiki/Major_depressive_disorder}{major depressive disorder} or \href{https://en.wikipedia.org/wiki/Dysthymia}{dysthymia}; it is a normal temporary reaction to life events, such as the loss of a loved one; \& it is also a symptom of some physical diseases \& a \href{https://en.wikipedia.org/wiki/Side_effect}{side effect} of some drugs \& medical treatments. It may feature sadness, difficulty in thinking \& concentration \& a significant increase or decrease in appetite \& time spent sleeping. People experiencing depression may have feelings of dejection, hopelessness and, sometimes, suicidal thoughts. It can either be short term or long term. 
\end{definition}
See also, not e.g. as usual but especially this time, \cite{Solomon2015}.

\begin{verbatim}
	// * * * * * * * * * * * * * * * * * * * * * * * * * * * * * * * * * * * * * //
\end{verbatim}

\subsubsection{\texttt{{\color{cyan}{\it int}} {\color{YellowGreen} main} ({\color{cyan}{\it int}} argc, {\color{cyan}{\it char}} *argv[]) \{}}

\fbox{\textbf{1}} {\it Just another winter's day}$\ldots$

\begin{figure}[H]
	\centering
	\includegraphics[width=0.75\textwidth]{Berlin_winter2020}
	\caption{Winter 2020, Berlin, Germany.}
	\label{fig1}
\end{figure}
Light precipitation (rain $+$ snow) outside the window (see Fig. \ref{fig1})$\ldots$ The room is so quiet that I can even hear my own heartbeat when I lay down: So annoying (see Fig. \ref{fig2})$\ldots$

\begin{figure}[h]
	\centering
	\includegraphics[width=0.75\textwidth]{Arima_Kousei_you_are_in_my_way}
	\caption{Arima Kousei's emotional explosion when performing piano with Nagi, Episode 18: {\it Hearts Come Together}, {\it Your Lie in April} (2014--2015).}
	\label{fig2}
\end{figure}

\begin{quotation}
	``{\it I can't concentrate. I can still hear the sound of the piano. It's in my way. You're in my way, so$\ldots$ Get out of here!}'' - Arima Kousei, {\it Your Lie in April} (2014--2015).
\end{quotation}

I always have this kind of weird feeling\footnote{An indication of bad writing?} when I try to write something \& connect all, or as many as possible, the ideas in my flow of thoughts together, you know? This even does not sound right to me because the question is: {\it How the hell you can organize \& \href{https://en.wikipedia.org/wiki/Connect_the_dots}{connect the dots} when that flow in your head is a real shitty mess?}

In such a situation, as usual, a bad writing habit [no chronological order, switch randomly between English \& Vietnamese\footnote{{\it If the real purpose is to express your thoughts{\tt/}ideas{\tt/}emotions, why does it matter to write under a lot of unnecessary{\tt/}redundant chains{\tt/}constraints{\tt/}criteria?}}, etc.] can be partially accepted \& invoked as an unavoidable replacement to a ``standard'' one: {\it Just let it all out naturally$\ldots$}

\begin{verbatim}
	// * * * * * * * * * * * * * * * * * * * * * * * * * * * * * * * * * * * * * //
\end{verbatim}

\noindent
\fbox{\textbf{2}} Đã mấy tháng nay, mình thường cố ý đi làm trễ khoảng 1 tiếng, để đợi những người khác đến chỗ làm gần đủ hết, vì mình chỉ muốn tránh thủ tục chào hỏi xởi lởi{\tt/}xã giao ``khích lệ tinh thần'' phiền phức mỗi buổi sáng. Mình đánh hơi thấy mùi fake và bắt đầu trở lên im lặng như một năm về trước. Mình cũng dần dần hiểu được phần nào ý nghĩa của từ ``guồng'' mà các anh lớn hơn với thằng bạn mình hay nhắc tới: Mọi thứ [công việc + cuộc sống cá nhân] bắt đầu lặp đi lặp lại một cách đơn điệu (monotonously), màu sắc cứ dần dần mà phai nhạt ngày qua ngày$\ldots$

À xém chút quên ngữ cảnh, khi những dòng này được viết, mình đã nếm trải 10 tháng đầu tiên của hành trình làm nghiên cứu sinh ở Đức\footnote{Usually called ``PhD'' in general \& ``Doktorand'' in German.} bên Applied Math, cụ thể là về Tối Ưu Hình Dáng\footnote{Dịch thuật ngữ chuyên ngành toán ra tiếng việt thường mang lại cảm giác củ chuối như vậy! Sad$\ldots$} (\href{https://en.wikipedia.org/wiki/Shape_optimization}{Shape Optimization}) cho ống dẫn khí của động cơ đốt trong. Khi mới bắt đầu làm thì mình thấy thích topic này lắm, nhưng đến giữa chừng thì lại bắt đầu chán. {\it Là do mình hay công việc nghiên cứu nó vốn tẻ nhạt như vậy?}$\ldots$

\begin{verbatim}
	// * * * * * * * * * * * * * * * * * * * * * * * * * * * * * * * * * * * * * //
\end{verbatim}

\noindent
\fbox{\textbf{3}} Vào một ngày đông nào đó (vì bữa đó quá đỗi bình thường nên mình chả thèm{\tt/}thể nhớ chính xác là ngày nào) mình lại đạp xe đến chỗ làm, cố gắng tránh né ánh mắt của mọi người. Đến trưa thì đợi ông thầy phụ của mình đi rủ họ ăn chung với nhau, rồi mình lại xách hộp cơm tự nấu đem theo ra khỏi tòa nhà, đến công viên đối diện để ngồi ăn giữa trời tuyết: {\it lủi thủi} và {\it cô độc}. Vừa ăn trưa mình vừa ngắm chim (ở Berlin chim bồ câu đâu ra cả nùi nùi và đặc biệt là tụi chim này rất dạng, chắc đất chốn đây hẳn là lành lắm!). Một lát sau mình lại lết lên phòng làm việc. {\it Nặng nề} và {\it ủ dột}, mình vẫn phải tiếp tục làm những tasks được giao nhưng không hề cảm thấy thích như mọi khi mình được tự ý làm những thứ mình tự chế{\tt/}bày ra nữa. Haizz$\ldots$ Cố làm vậy, rồi từ từ cũng sẽ hết ngày, lại đi về, nấu cơm ăn, xem YouTube đến khi mỏi mắt, rồi ngủ, rồi ngày sau lại lặp lại y như thế: {\it Ôi cái guồng này nó làm mình chán phát điên mất} - mình vừa làm vừa nghĩ vậy$\ldots$

Đợi đến tầm xế chiều, thì cô bạn Maroc ngồi đối diện lưng-lưng (not mặt-mặt) với mình moi ra một hộp Chocolate để tặng mình. Hề cái là cô này phải đợi anh bạn người Đức chung phòng đi lấy{\tt/}nhả nước gì đấy thì cổ mới có dịp để tặng riêng cho mình (phòng mình là phòng toàn PhDs duy nhất của cả nhóm, gồm 3 đứa mình). Tỏ tình? Nah, nhìn lại cái thân hình đang ngày càng trì trệ\footnote{Ironically, I am currently working on Shape Optimization but my body shape gets far \& far away from ``optimal''.} mấy tháng qua của mày đi. Đùa chứ bản nói vì mình đã giúp bản rất nhiều thủ tục giấy tờ - mà mình đã phải vật lộn rất nhiều trong mấy tháng đầu nhờ đợt dịch - khi mới tới nên bản biết hơn (mình đổi từ ``cổ'' sang ``bản'' vì mình để ý tên email riêng của cổ: à, thì ra cô này bằng tuổi mình). Bản nói bản biết ơn mình và ông thầy phụ của mình lắm, trong mắt cô bạn mới tới thì đây là 2 người tử tế giúp đỡ bản nhiều nhất trong nhóm$\ldots$

{\it Ồ, thì ra vẫn có người trong nhóm xem mình là người tốt{\tt/}tử tế à?} Thật tình thì lúc mình giúp cô bạn thì mình chả nghĩ gì nhiều, giúp thì giúp thôi. Mình là kiểu người ngu dốt kiểu vậy. À không, còn tệ hơn nhiều cơ: {\it a giver}\footnote{\href{https://www.youtube.com/watch?v=YyXRYgjQXX0}{Are you a giver or a taker? | Adam Grant}, TED.}. Việc mình trở thành 1 giver chắc là do ảnh hưởng từ mẹ mình. Hồi nhỏ mẹ mình hay dạy là con cứ tốt với mọi người xung quanh thì họ sẽ tốt lại với con thôi, chứ đừng cứ sống ích kỷ rồi không ai chơi với con hết, đại loại vậy. Đó cũng là bài học đối nhân xử thế đầu tiên của mình và ảnh hưởng đến tính cách mình rất nhiều$\ldots$

Nghe cô bạn cảm ơn, mình nhìn hộp Chocolate, một chút vui thoáng qua nhưng cũng chợp tắt thật nhanh. Mình không cho phép mình vui lâu hơn vì chắc gì đó là những lời thật lòng, hay {\it chẳng qua chỉ là xã giao với nhau để được giúp nhiều hơn mà thôi?} Vì mình hiểu rằng bất kể mình có tốt bụng và tử tế cỡ nào đi chăng nữa, cô bạn này sau một thời gian nữa sẽ ghét và khinh mình vì những lời xuyên tạc từ những người khác trong những bữa ăn trưa tưởng chừng như thân thiện giữa các đồng nghiệp với nhau đó mà thôi$\ldots$

Một lúc sau thì anh bạn Đức về phòng. Anh này hiện đang chí mí cuối năm cuối PhD và đang bức tốc để chuẩn bị về đích{\tt/}lên đỉnh Olympia, à nhầm, bảo vệ. Mình ngỏ ý chia 50--50 với ảnh, để cho không khí trong phòng 3 đứa mình đỡ căng thẳng. Anh từ chối thẳng, yêu cầu mình để anh ấy yên, anh ta không cần bất cứ thứ gì từ mình. Uầy, vô tình chuốc thêm căng thẳng rồi. Ngu người thật. Đành ăn một mình vậy. Bụng mập càng thêm mập$\ldots$

Cô bạn nói mình lần sau sẽ mua loại ngon hơn, lần này thì chỉ còn loại đó. Có vẻ thật lòng - mình nghĩ. Xem ra mình còn chút hy vọng xót lại vào humanity trong mối quan hệ giữa người với người. Mà nếu cô bạn canh để tặng quà Giáng sinh thì bữa đó là ngày 18 tháng 12, ngày đi làm cuối cùng của năm 2020, giờ mình mới có manh mối để nhớ ra ngày đó. Mùa nghỉ đông đầu tiên kéo dài hơn 40 ngày cuối cùng cũng bắt đầu\ldots

\begin{verbatim}
	// * * * * * * * * * * * * * * * * * * * * * * * * * * * * * * * * * * * * * //
\end{verbatim}

\noindent
\fbox{\textbf{4}} Khoảng thời gian nghỉ đông này rất quý giá để mình đầu tư vào work-life balance lại, bằng cách$\ldots$ làm việc nhiều hơn và đặc biệt là chỉ làm những thứ mình muốn, chứ không phải những tasks nhàm chán, đôi khi ngu xuẩn, mà ông thầy phụ mình cứ thảy cho suốt mấy tháng qua. Đó là cũng là khoảng thời gian yên tĩnh để mình có thể suy nghĩ về những chuyện đã xảy ra trong những tháng qua. Có vài chuyện mình không tài nào hiểu nổi, đương nhiên vẫn là những câu chuyện muôn thuở về mối quan hệ giữa con người với con người. Mình mà lị: {\it The Trouble Boy}!!$\ldots$ Đại loại là, từ lúc mới đến Berlin tới giờ, bất kể mình giúp đỡ những người khác mỗi khi họ cần, họ đều cư xử toxic ngược lại, thậm chí có anh postDoc còn verbally bully mình. Thế mà mình cứ ảo{\tt/}hoang tưởng là họ hoan nghênh mình tới để làm việc, trao đổi kiến thức và cống hiến cơ đấy, nhưng thì ra suất Marie--Curie này không dễ ăn{\tt/}handled như mình vẫn tưởng bở. Môi trường học thuật mà thằng ngu này! Mình lại suy nghĩ quá đơn giản rồi. Chán cái đầu nông cạn{\tt/}ngây thơ (naive) của mình thật! Thôi cứ đành ngậm họng lại, tránh tiếp xúc với mọi người (đương nhiên là ngoại trừ boss bự với boss nhỏ, ngu thì ngu chứ chưa muốn bị đuổi!) để khỏi bị bắt nạt, và giữ năng lượng để tập trung làm việc vậy$\ldots$ Còn 2 năm mấy, 3 năm nữa cơ mà! Đừng có kiểu mới nhập Viện chưa bao lâu mà đã gục{\tt/}tịt ngòi chứ!$\ldots$

Ấy vậy mà, mặc dù đã tập trung làm việc miệt mài, tinh thần và cảm xúc mình vẫn bị kéo xuống liên tục, và các mối quan hệ với đồng nghiệp cứ thế càng ngày càng tệ: {\it Thà không có thì tốt biết mấy!} - Mình chợt nghĩ. Cứ thế này thì mình sẽ bị vắt kiệt (drained out{\tt/}\href{https://en.wikipedia.org/wiki/Occupational_burnout}{occupational burnout}) mất. Hóa ra hy vọng, niềm hân hoang, etc. chỉ là những thứ cảm xúc vô nghĩa mình tự tạo ra, để tự huyễn hoặc, và làm tiền đề cho sự chán chường nặng nề và thất vọng tột độ mà thôi$\ldots$

Được cái là ông thầy phụ (co-supervisor) của mình rất có tâm. Ổng giúp mình rất nhiều khi mình mới tới Berlin, từ cả nhà cửa, giấy tờ thủ tục, đến hầu hết tất cả các khía cạnh trong công việc. Nhận ra mình đang có dấu hiệu xuống tinh thần, ổng đôn đốc mình làm việc nề nếp, khoa học hơn, giao nhiều tasks hơn, và rồi cứ 4 giờ chiều mỗi ngày mình sẽ đến tận phòng để gặp và báo cáo ổng theo mệnh lệnh ở Thư viện chung, trước mặt của các nữ sinh viên đang học Master của nhóm. Thật tình thì đa số mấy tasks ổng giao mấy tháng gần đây có phần ngô nghê và vớ vẩn thiệt. {\it Hay lúc nào cũng nhảm vậy nhỉ? Hay do mình đã tiến bộ mà không tự nhận thức được?} Không phải chảnh, đơn giản vì mình là người trực tiếp upgrade software, line-by-line, nên mình hiểu rõ nhiều cái technicals mất dạy mà phải tốn có khi cả tháng mới vượt qua được, nên nhiều cái ổng không biết mà chỉ bậy{\tt/}bừa{\tt/}xàm cũng không phải là điều quá ngạc nhiên. Mình là người kiểu vậy: Không quan trọng supervisor khủng (vì họ chỉ tổ bận mà thôi), chỉ cần supervisor có tâm là được. Tình cảm thầy trò vẫn quan trọng hơn là những danh tiếng hào nhoáng (reputation) mà mình chỉ hưởng xoáy chứ không phải có được nhờ chính thực lực của mình.

Tasks ngu thì ngu nhưng kệ vậy, miễn ổng tốt với mình là được rồi, cứ làm nhiệm vụ được giao vậy. Rồi sẽ ổn cả thôi, dù ảm đạm và cô độc đến cỡ nào đi nữa$\ldots$

\begin{verbatim}
	// * * * * * * * * * * * * * * * * * * * * * * * * * * * * * * * * * * * * * //
\end{verbatim}

\noindent
\fbox{\textbf{5}} Mình chợt nhớ lại 10 tháng trước. {\it Mình làm ở đây mới có 10 tháng thôi mà đã thay đổi nhiều vậy cơ á?} 10 tháng trước, mình lao vào giữa vùng dịch, để làm lại từ đầu lần nữa. Quyết liệt và dữ dội, tưởng chừng không có gì cản trở được mình nữa, thì mình lại bắt đầu có dấu hiệu gục hệt như một năm trước lúc mình học Master về Pure\footnote{``Pure'', not ``poor''. But, actually, they sound similar. \& that is exactly the reason why I have switched from Pure Math to Applied Math.} Math ở Pháp. Thật tự đáng hổ thẹn và thất vọng$\ldots$

Khoảng thời gian nghỉ đông tĩnh lặng và lạnh lẽo giúp mình có thể tập trung suy nghĩ sâu hơn, về nhiều thứ cực kỳ mâu thuẫn đã và đang diễn ra khiến mình thật sự bối rối và bắt đầu nghi ngờ. Điều này làm mình chợt nhớ lại cuốn sách \cite{Bancroft2003} về ``{\it Kẻ thao túng tâm lý{\tt/}bạo hành tinh thần}'', mà mình biết tới nhờ một lần đọc một psychology post về ``{\it Kẻ bạo hành âm thầm}''. Well, thực tình mà nói đây là một cuốn sách về bạo hành gia đình, cả bạo hành thể chất lẫn bạo hành tâm lý nhưng chú trọng đặc biệt về bạo hành tâm lý, cụ thể là tác động từ người chồng gia trưởng lên người vợ lẫn những đứa con đáng thương của ông ta. {\it Daijoubu desu!} Mình vẫn còn con nít mà, nên vẫn có thể áp dụng vế sau lên chính trường hợp của mình được!

{\it Một luồng những suy nghĩ, câu hỏi tự vấn từ sâu trong nội tâm bắt đầu tuôn trào}\footnote{Actually, I have been seriously overthinking about my overthinking problem all the time!$\ldots$ {\it Wait, should I overthink about that also?!!}}$\ldots$

\begin{verbatim}
	// * * * * * * * * * * * * * * * * * * * * * * * * * * * * * * * * * * * * * //
\end{verbatim}

\noindent
\fbox{\textbf{6}} {\it Another flow of thoughts}$\ldots$
\begin{itemize}
	\setlength\itemsep{0em}
	\item {\it Why did almost all of my colleagues treat me so toxic like that despite the fact that I have helped them several times whenever they asked me for my hand. Don't they understand what the term ``grateful'' really means?}
	\item {\it Why does a good person like me have to suffer all these toxic behaviors in lab-environment \& even social isolation?}
	\item {\it Is this world designed for good{\tt/}kind people?}
	\item {\it Does there exist a place on this Earth for all human beings with full of purely good intentions to live \& work?}
	
	In a moment, it seems to me that there is no place for a person with low \href{https://en.wikipedia.org/wiki/Emotional_intelligence}{Emotional Intelligence} (EQ) \& very low \href{https://en.wikipedia.org/wiki/Social_intelligence}{Social Intelligence} (SI) but \href{https://hsperson.com/}{Highly Sensitive} (HSP) like me in this cruel world.    
	\item {\it Why do I have to avoid all these people \& develop psychological problems, especially \href{https://en.wikipedia.org/wiki/Impostor_syndrome}{imposter syndrome}\footnote{The \#1 mental illness in academic{\tt/}lab environment?}?} I have conducted a lot of tasks, even much more than my predecessors, then {\it why do I have to keep suffering?}
	\item {\it Why were I called ``The Lucky Guy''? Why does hardly anybody give a minimum level of respect to me even when I have been conducting such a huge amount of work? Is it because my skin is too yellow to be considered white enough?}
	\item {\it Why do I have to do these stupid{\tt/}bullshit tasks \& report daily to my co-supervisor. \& especially, whenever I refused to do such an unnecessary{\tt/}redundant task, he started to be wildly angry immediately?} But he is a truthfully good, kind, honest, \& enthusiastic co-supervisor, because he has helped me a lot through these tough times in this strange year.
\end{itemize}
I really do not understand: {\it absolutely confused} \& {\it overwhelmed}$\ldots$

{\it Tại sao mình có cảm giác bị đâm lén sau lưng liên tục nhưng không biết chính xác là ai cả? Tại sao ai cũng đối xử mình như thứ rác rưởi mặc dù mình giúp họ khá nhiều? Tại sao mình cảm giác bản thân không đủ tốt và làm việc chưa đủ cần cù trong khi mình đã làm được những điều mà những người tiền nhiệm ở vị trí của mình chưa từng làm được: ví dụ nâng cấp phần mềm được phát triển và sử dụng suốt 10 năm trong vòng chưa tới 10 tháng?}

Yeah, everything is so weird, very weird, super weird. This kind of feeling is exactly like there is an invisible ``Black Hole of Truth'' somewhere behind your back. You can feel its existence as a blur, but you cannot see it or touch it directly. So annoying$\ldots$ Very tired$\ldots$ Completely exhausted$\ldots$

- {\it Can you figure out the reason why I hate such a ``Black Hole'' that much?}

- Is it because that hole is black?

- No, you racist fuck! Because$\ldots$ Like the way every matter, e.g. light (`photon') is rapidly sucked towards the `singularity' at the center of a real \href{https://en.wikipedia.org/wiki/Black_hole}{Black Hole}, all truths are also sucked towards this ``Truth Destroyer'' \& then demolished{\tt/}consumed by it.

\begin{verbatim}
	// * * * * * * * * * * * * * * * * * * * * * * * * * * * * * * * * * * * * * //
\end{verbatim}

\noindent
\fbox{\textbf{7}} {\it Có phải mình đang cố đổ lỗi lên mọi người xung quanh để chối bỏ mọi chuyện tệ hại là lỗi của mình không?}

- Không hề.

- Xì{\tt/}hừm{\tt/}hứ{\tt/}etc., thằng nhãi ranh đó lại thế, chả bao giờ biết tự nhận{\tt/}chịu trách nhiệm (take the responsibility), chỉ biết đổ lỗi lên đầu người khác. Xem ra việc làm nghiên cứu sinh tưởng chừng sẽ giúp nó trưởng thành, chín chắn hơn phần nào nhưng giờ nó lại thế, vẫn thế, có khi còn tệ hơn! Đảm bảo nó sẽ mãi là đứa con nít không bao giờ lớn nổi.

- Ừm, có khi nói vậy cũng đúng phần nào$\ldots$

Tới đây thì bạn cần hiểu một điều rằng: {\it khi bạn tự thấu hiểu chính bản thân từ bên trong, thì mọi lời phán xét{\tt/}đàm tiếu bên ngoài không còn quan trọng nữa}$\ldots$

Mình không thể nào nói là mình không quan tâm (I do not care!) tới những lời phán xét như vậy, bởi vì thực sự là mình có quan tâm. Mình quan tâm là vì mình muốn bảo vệ sự thật (lại một hành động ngu xuẩn nữa trong chuỗi dài những hành động ngu xuẩn: scientist wannabe, huh?). Và một trong những bước cốt lõi để bảo vệ sự thật đó là phát hiện những lời giả dối{\tt/}đàm tiếu{\tt/}vu khống (Lie Detector) và điều chỉnh những lời phát xét bị biến dạng (distorted judgments) do tác dụng của chính cái ``Black Hole of Truth'' gây ra$\ldots$

{\it Bạn biết vì sao nhiều người thích phán xét thế không?} Đó là vì việc phán xét rất dễ: không cần phải suy nghĩ{\tt/}động não nhiều, không cần tốn quá nhiều sức, chỉ cần hướng thẳng tới việc công kích{\tt/}đả thương người khác là bất cứ ngôn từ nào phát ra từ miệng một người tự nhiên{\tt/}tự động{\tt/}mặc định trở thành lời phát xét tiêu cực (đôi khi chết người) ngay lập tức. {\it Thế thì nghe tiêu cực nhỉ? Nhưng nếu bản chất là tiêu cực vậy thì tại sao còn những người muốn làm điều như vậy?} Nhiều là đằng khác! Đó là vì: phán xét tuy rất dễ nhưng đem lại rất nhiều cảm giác sảng khoái cho ``người thẩm phán'': cảm giác tha hồ nếm trải quyền lực{\tt/}sự thống trị  lấn át lên trên nạn nhân để thỏa mãn cái tôi mà không cần tốn nhiều sức lực. Trái lại, việc bình tĩnh suy nghĩ một cách thấu đáo để thực sự hiểu một vấn đề lại khó hơn rất nhiều: đòi hỏi đầu tư nhiều thời gian và công sức hơn, và đặc biệt là những kinh nghiệm tích lũy mà thông thường chỉ những người từng trải và{\tt/}hoặc đủ khôn ngoan{\tt/}sáng suốt mới có được. Mặc dù cần đầu tư nhiều như vậy, nhưng điều nực cười là việc thấu cảm (also thấu hiểu, cảm thông) thường không nhận được bất kỳ sự ủng hộ nào từ số đông do ảnh hưởng của tư tưởng bầy đàn{\tt/}hiệu ứng đám đông\footnote{Điều này là đối tượng nghiên cứu chính của nhánh chuyên ngành Tâm lý học đám đông của Tâm lý học xã hội.}. {\it Again, another sad reality}$\ldots$

Uòi uòi uòi, đoạn này hơi bị tâm đắc, dịch lại phát nữa cho bớt sung vại:

{\it Do you know why so many people love judging?} That is because judging is super easy{\tt/}simple: no need to think{\tt/}consider too much \& can be generated effortlessly, the only attempt needed is that you have to concentrate on hurting{\tt/}bullying verbally your victim(s) violently{\tt/}brutally{\tt/}vulnerably, then every single word coming from your mouth will be automatically toxic immediately \& then they form together a deadly judgment. {\it Sound negative, huh? But if the nature of judging is negativity, then why so many people keep doing that?} Although judging is easy, however, in return, it helps them gain a lot of strength \& power: the feeling of tasting the dominance \& several psychological powers over their victim(s) without pending a considerable amount of time \& effort: A {\it hyper-win-lose} situation! On the contrary,
\begin{quotation}
	``{\it In order to be able to think, you have to risk being offensive}.'' - {\sc Jordan B. Peterson}
\end{quotation}
the ability to calm yourself in order to consider{\tt/}contemplate{\tt/}examine events{\tt/}statements{\tt/}situations consciously \& logically\footnote{In reality, the set of non-logical events{\tt/}statements is dense in the set of all events{\tt/}statements, while the set of logical ones, though infinite, has zero measure in that \href{https://en.wikipedia.org/wiki/Universal_set}{universal set}. This pair exists together, both mutually exclusive \& supporting each other simultaneously. Like the way the pair of the rational set $\mathbb{Q}$ \& the irrational one $\mathbb{R}\backslash\mathbb{Q}$ does:
	\begin{itemize}
		\item[(i)] The universal set here is $\mathbb{R}$. The supporting \& mutually exclusive relationships just mentioned can be expressed mathematically by:
		\begin{align*}
			\mathbb{Q}\cup(\mathbb{R}\backslash\mathbb{Q}) = \mathbb{R},\ \mathbb{Q}\cap(\mathbb{R}\backslash\mathbb{Q}) = \emptyset.
		\end{align*}
		\item[(ii)] The set of rationals has Lebesgue measure zero \& their \href{https://en.wikipedia.org/wiki/Cardinality}{cardinality} is the 1st transfinite number, \href{https://en.wikipedia.org/wiki/Aleph_number}{aleph-null} ($\aleph_0$):
		\begin{align*}
			\operatorname{m}_1(\mathbb{Q}) = \operatorname{m}_1(\mathbb{N}) = 0,\ \operatorname{card}(\mathbb{Q}) = \operatorname{card}(\mathbb{N}) = \aleph_0.
		\end{align*}
		\item[(iii)] The cardinality of the set of irrationals is continuum:
		\begin{align*}
			\operatorname{m}_1((\mathbb{R}\backslash\mathbb{Q})\cap[a,b]) = \operatorname{m}_1([a,b]) = b - a,\ \forall a,b\in\mathbb{R}, \mbox{ \& } \operatorname{card}(\mathbb{R}\backslash\mathbb{Q}) = \operatorname{card}(\mathbb{R}) = \mathfrak{c} = 2^{\aleph_0}.
		\end{align*}
\end{itemize}} to connect as many dots as possible, requires much more time \& efforts, especially the expertise{\tt/}experiences, which are usually possessed{\tt/}gained only by hardened and{\tt/}or wise enough individuals. Ironically, in spite of all of these serious investigations, sympathy, empathy, \& humanity have been almost all the time being undervalued{\tt/}underrated{\tt/}underestimated, even ignored from the majority of human beings, which is the prototypical \& fundamental effect of \href{https://dictionary.apa.org/majority-influence}{majority influence}.

\begin{verbatim}
	// * * * * * * * * * * * * * * * * * * * * * * * * * * * * * * * * * * * * * //
\end{verbatim}

\noindent
\fbox{\textbf{8}} - {\it Ừa thì đành tạm chấp nhận là mày không cố đỗ lỗi lên người khác, thế nhưng mày đang muốn nói điều gì?}

Điều mình đang muốn nói ở đây là: Trong bất kỳ mối quan hệ hay tương tác giữa người với người nào, điển hình ở đây là môi trường làm việc (workplace), khi những trải nghiệm về thực tại của bạn trở nên cực kỳ tệ và độc hại, trái ngược hẳn với tất cả những viễn cảnh mà bạn đã và đang mong đợi so với những gì bạn cho đi, kể cả sau khi đã tự nghi ngờ lòng tốt{\tt/}trắc ẩn lẫn kỹ năng{\tt/}thực lực của bản thân (self-doubts, imposter syndrome) đủ lâu để biết chắc rằng bạn tuyệt đối không phải là vấn đề, thì bạn cần phải hiểu một điều rằng: {\it You are not the problem. Instead, you are just another victim}$\ldots$

{\it But the victim of what? I do not understand.} Like the metaphor about the Black Hole of Truth mentioned above, the conflict between your genuine perception \& the toxic reality in any relationships (see, especially, \cite{Bancroft2003}) and{\tt/}or any working environments, except the case of \href{https://en.wikipedia.org/wiki/Schizophrenia}{schizophrenia}\footnote{Tâm thần phân liệt.} of course, as a matter of fact, indicates the existence of such an invisible truth-destroyer within your environment, which is usually very close to you: {\it a psychological manipulator}$\ldots$

\begin{quotation}
	``{\it The first step in solving a problem is to recognize that it does exist}.'' - Zig Ziglar
\end{quotation}

Actually, I met \& dealt with such a person in the past, 1 year ago when I studied my Master in Pure Math in France. Ironically, although I already had some experience or thought to have it, it took me a lot of months{\tt/}efforts to detect who is the real monster in my current working environment. Sadly, I even paid a lot of trusts \& respects to that manipulator for a long time$\ldots$ Because the situation at this time is much worse for me: a psychological manipulator with a master manipulation skills in a very, very different level$\ldots$
\begin{itemize}
	\item {\it Why do, no matter how many seeds of kindness I have planted \& distributed to others, almost of them still treat me like trash{\tt/}garbage by verbally bulling, distributing toxic behaviors in return, \& even dehumanizing me?}
\end{itemize}
It seems to me that the main objective of a typical psychological manipulator is to destroy his{\tt/}her victim(s)'s confidence \& the ability to work 1st, then ruin his{\tt/}her{\tt/}(their) life later, not physically but psychologically \& mentally. However, in fact, our mental \& physical health are strongly connected \& correlated with each other. Thus, if one of them gets destroyed, then both will be shattered{\tt/}demolished{\tt/}messed up{\tt/}screwed up{\tt/}fucked up at the end anyway. \& this typical psychological manipulator seems to desire to enslave his{\tt/}her victim(s). The interesting point of this process is that it is conducted not immediately, but gradually instead. This makes sense to me. Because if someone suddenly comes to you \& forces you to form a master-slave relationship, you will certainly refuse{\tt/}deny it immediately!\footnote{Won't you?}\,\footnote{However, if that person offers you a chance to become his{\tt/}her master, the situation is far different from the one I am trying to explain here though!} That is exactly the strategy of psychological manipulators: they tighten the rope so slowly{\tt/}gently{\tt/}gradually \& carefully\footnote{Typically, they will, e.g., give some favor(s){\tt/}do some good{\tt/}kind{\tt/}nice stuff(s) whenever you seem, or they thought you seem, to start doubting them and{\tt/}or recognize their pattern of invisible destructive actions{\tt/}behaviors{\tt/}tricks.} that you cannot recognize{\tt/}detect it during a long time\footnote{Or just because I am the only sucker here, guys? \texttt{.\_.} Guys?!!}. \& then one day, you become his \texttt{private object}\footnote{Instead of the \texttt{public class} you deserve!} without any self-awareness left - the special \& critical piece of your character which has been being tricked since a long time ago. That explains the reason why psychologists use the term ``manipulation{\tt/}manipulator'' instead of ``destruction{\tt/}destroyer'' for such kind of action{\tt/}actor{\tt/}actress in this particular context.

The relationship with a psychological manipulator is like a never-ending{\tt/}endless nightmare that no matter how hard you have been struggling to wake up, you just keep going deeper \& deeper to the rock bottom, and$\ldots$
\begin{quotation}
	``{\it I've had a lot of what I thought were rock bottoms, only to discover another, rockier bottom underneath}.'' - BoJack Horseman, {\it BoJack Horseman} (2014--2020).
\end{quotation}
Under several psychological effects of a psychological manipulator, like a black hole, everything gets sucked in \& destroyed brutally. \& even worse, that kind of person absorbs all the lights in our daily life \& then spits all darkness which each of us never wants to experience any single moment in it. All the truths are deformed{\tt/}perturbed{\tt/}twisted via the dark-magic mouth of a psychological manipulator to gradually become a germ of harsh realities in the past of somebody \& nightmare scenarios of life in his{\tt/}her future. Nothing can grow{\tt/}thrive{\tt/}flourish. The necessity to deal with this situation becomes urgent, critical, \& even vital!

\begin{verbatim}
	// * * * * * * * * * * * * * * * * * * * * * * * * * * * * * * * * * * * * * //
\end{verbatim}

\noindent
\fbox{\textbf{9}} {\it - I still do not understand. Your arguments do not make sense to me.}

- Daijoubu! You will understand such a thing (only) when you really suffer from it. Moreover, you should be aware of the fact that even if you do not understand something, or it does not make any sense to you, does not imply it is wrong, because: Not relative like gravity, {\it truth is absolute}.

Back to the question of dealing with a psychological manipulator: {\it What should I do then?} It seems to me at the time I write these lines that there are only 2 options to ``handle'' psychological manipulations:
\begin{enumerate}
	\item {\it Funnily enough, you can actually solve a lot of problems in your life by ignoring them!} Or,
	\item {\it You have to spot that monster{\tt/}parasite{\tt/}Shadow King (Fig. \ref{fig3}) out \& face{\tt/}fight{\tt/}deal with it.}
\end{enumerate}

\begin{figure}[H]
	\centering
	\includegraphics[width=0.75\textwidth]{Shadow_King}
	\caption{Amahl Farouk (Shadow King), {\it Legion} (2017--2019).}
	\label{fig3}
\end{figure}
So the next question is: {\it Which option{\tt/}strategy should I choose?} My intensely personal answer is that your choice will{\tt/}should depend on your own personality trait\footnote{See also, e.g.,
	\begin{itemize}
		\item \href{https://en.wikipedia.org/wiki/Big_Five_personality_traits}{Wikipedia/Big 5 personality traits}.
		\item \href{https://en.wikipedia.org/wiki/Extraversion_and_introversion}{Wikipedia/Extraversion \& introversion}.
		\item \href{https://psychology.wikia.org/wiki/Extraversion_and_introversion}{Psychology Wiki/Extraversion \& introversion}.
		\item \href{https://psychology.wikia.org/wiki/Ambiversion}{Psychology Wiki/Ambiversion}.
		\item \href{https://www.youtube.com/watch?v=pCceO_D4AlY}{Jordan B. Peterson. 2017 Personality 14: Introduction to Traits/Psychometrics/The Big 5}.
		\item \href{https://www.youtube.com/watch?v=qYvXk_bqlBk}{Who are you, really? The puzzle of personality | Brian Little}. TED.
\end{itemize}}.
\begin{quotation}
	``{\it Life is a matter of choices, \& every choice you make makes you}.'' - John C. Maxwell
\end{quotation}
The most important question to me, in the matter of choice in life, is: {\it Are you introvert or extrovert, or in between: ambivert?}

Back to the turning point where we are trying to make a move, the former seems to be the favorite of wise human beings, meanwhile individuals with high integrity absolutely{\tt/}definitely choose the latter\footnote{Cf. Người khôn ngoan vs. kẻ chính trực - câu chuyện muôn thuở không bao giờ có hồi kết.}.

Let me assume that you have just chosen the 1st option: {\it Just ignore all the matters around you, keep working \& living happily in your own}. Sound idealistic, huh? Yeah, everything usually seems naive, simple, idealistic in the 1st glance like that$\ldots$
\begin{quotation}
	``{\it When there is no enemy within, the enemy outside can do you no harm}.'' - African proverb
	
	``{\it It's not always necessary to be strong, but to feel strong}.'' - Chris McCandless, {\it Into the Wild} (2007)
\end{quotation}
Provided you are{\tt/}feel strong inside, you do not care about psychological manipulator(s) anymore: {\it No enemy inside $\Rightarrow$ No enemy outside}. Sound perfect then! But the right{\tt/}real question is: {\it Can you always be strong{\tt/}keep being strong all the time in your life?}
\begin{quotation}
	``{\it Ups \& downs in life are very important to keep us going, because a straight line even in an ECG means we are not alive}.'' - Ratan Tata
\end{quotation}
It does not really matter you want this situation or not, you know? You have to deal with it when you already got involved. {\it Can you ignore the negative, even deadly{\tt/}fatal, effects of psychological \& emotional manipulation? Are you able to work, live, or just even breathe in that toxic atmosphere{\tt/}environment?} When you are{\tt/}feel strong enough, the ignoring-strategy seems to work effortlessly but charmingly. However, the critical point(s) is{\tt/}(are) located in the duration when you are in your downward spiral: {\it Does ignorance really help anymore?} Even assume you can ignore{\tt/}bear it a couple of times: {\it How long then}? {\it How many chances you have left before the ultimate mental breakdown comes to you?}

Furthermore, it is highly likely psychological manipulator(s) will perfectly choose your darkest time to attack \& mess your brain up: {\it A real nightmare begins}$\ldots$

\begin{verbatim}
	// * * * * * * * * * * * * * * * * * * * * * * * * * * * * * * * * * * * * * //
\end{verbatim}

\noindent
\fbox{\textbf{10}} It is kind of scary \& full of crazy shit, you know? Like a mind game$\ldots$ Chottomatte! No, it is actually a mind game, which should not exist in this humanity in the 1st place.

- {\it But why does this kind of redundant mind game exist in the human race?}

- Because of the term ``race'' in that phrase.

It seems that psychological manipulation is just another by-product on the \href{https://en.wikipedia.org/wiki/Human_evolution}{human evolution} at the 1st glance: {\it generally{\tt/}mostly unexpected} \& {\it unnecessary}. But look closer, think deeper: it actually can be {\it totally predicted} \& {\it absolutely unavoidable}.

Hence, this kind of mind game is compulsory{\tt/}mandatory then. The next couple of questions should be:

- {\it Do you want to play it or not (struggle or surrender)? \& if you do, how are you going to play it?}

- The answer lies in the 2nd option: {\it You must hunt them down!$\ldots$}

\begin{verbatim}
	// * * * * * * * * * * * * * * * * * * * * * * * * * * * * * * * * * * * * * //
\end{verbatim}

\noindent
\fbox{\textbf{11}} Let me postulate{\tt/}posit from now on the following assumption:

\begin{assumption}
	You target is to hunt the psychological manipulator(s) down in order to protect your mental health \& emotional{\tt/}physical{\tt/}psychological well-being.
\end{assumption}
So, you really want to reverse this predator-prey situation because you are too tired \& exhausted after having been being a poor prey for a long time, right? {\it How to hunt the predator back down then?}

My answer, at the moment this paragraph is written, is to think harder \& deeper \& keep questioning yourself: Because like in most predator-prey situations, you are the only person who has all the necessary pieces of the big picture. One of the simplest but most effective questions to spot out a psychological manipulator (PM, for short\footnote{But please, do not confuse this abbreviation with {\it Project Manager}, although they coincide in some situations!}) in your living{\tt/}working environment is exactly the title of the famous psychology book \cite{Bancroft2003}: {\it Why does he do that?}

It seems to me the most fundamental motive of psychologically manipulative behaviors is to gain the dominance \& powers from the victims by making them suffer both mentally \& psychologically, \& then, of course, physically. 

- I am still confused, {\it why is there such a PM in our civilized society?}

- Didn't you see it yet? Because:
\begin{quotation}
	{\it Psychological manipulation is the only way to make a coward become a fucking hero!!!}$\ldots$
\end{quotation}

\begin{quotation}
	``{\it I hate weak people. Weak people$\ldots$ never fight face to face. \& poison well. Despicable. Weak people$\ldots$ They lack patience. They will quickly reap what they sow. I killed people with these `protecting fists'. I blooded my master's previous Soryuu style. I couldn't respect my father's last words. That's right. That's who I wanted to kill}.'' - Akaza, before committing suicide, Chap. 155, {\it Demon Slayer: Kimetsu No Yaiba} (2016--2020).
	
	``{\it Sometimes people try to destroy you, precisely because they recognize your power - not because they don't see it, but because they see it \& they don't want it to exist}.'' - Bell Hooks
\end{quotation}
Ironically also,

\begin{quotation}
	``{\it I see it in the people that do the real work, \& what's sad in a way is that the people that are the most giving, hardworking, \& capable of making this world better, usually don't have the ego \& ambition to be a leader}.'' - Celine, {\it Before Sunset} (2004)
	
	``{\it The world is fucked by unemotional, rational men deciding shit}.'' - Celine, {\it Before Midnight} (2013)
\end{quotation}
Let me introduce the following table to ``approximately model'' the last 2 quotes:\footnote{It will be much clearer{\tt/}more logical{\tt/}more reasonable if you add the 3rd dimension to this two-dimensional personality-trait table: \textbf{Integrity}.}

\begin{table}[h]
	\centering
	\begin{tabular}{|c|c|c|}
		\hline
		{\sc Personalities}& \textbf{Ungifted{\tt/}untalented} & \textbf{Gifted{\tt/}Talented} \\
		\hline
		\textbf{Kind{\tt/}Generous} & an ordinary but lovable giver & an ideal{\tt/}perfect leader (e.g. Hồ Chí Minh) \\
		\hline
		\textbf{Selfish{\tt/}Egocentric} & a psychological manipulator & another fucking Adolf Hitler wannabe!!! \\
		\hline
	\end{tabular}
	\caption{Some combinations of personality traits \& their associated potentially developed characters{\tt/}roles in our ``civilized'' society.}
\end{table}
{\it Nhưng ta có nên đánh giá cả một con người chỉ bằng mỗi tiêu chí rằng cá thể đó có tài hay bất tài hay không?} Nếu câu trả lời là có một cách đồng loại, thì thế giới này sẽ ngày càng trở nên thực dụng và tình người hay sự thấu cảm rồi cũng sẽ bị phủ lên một màu xám xịt u ám và đơn điệu$\ldots$

Thực ra thì 1 câu hỏi khác nên xuất hiện sớm hơn trước câu hỏi kiểu dehumanizing đó: {\it Chúng ta có nên đánh giá{\tt/}phán xét ai đó hay không?}
\begin{quotation}
	``{\it You know what's more destructive than a nuclear bomb?$\ldots$ Words}.'' - Kim Jong-un, {\it The Interview} (2014)
\end{quotation}
I am not trying to complicate matters here, alright? Everything is already complicated like that by its own nature. Like the mysterious but fascinating depth in the nature of anything, we cannot measure{\tt/}understand them exactly{\tt/}directly{\tt/}absolutely, we can only try our best to create models which seem reasonable enough to us to reflect partially their insights in a deeper \& clearer way$\ldots$

\begin{quotation}
	``{\it This is real life. It's not perfect, but it's real}.''\footnote{There is no ultimately complete{\tt/}perfect person, since if that person existed, he{\tt/}she would be lack of the incompleteness part(s) at least.} - Jesse, {\it Before Midnight} (2017).
\end{quotation}
Another couple of questions arises: {\it How can we understand the depth of anything in life if they are so complicated like that? Perhaps it would be much better if we stop questioning ourselves?}\footnote{``{\it Life is like riding a bicycle. To keep your balance, you must keep moving}.'' - Albert Einstein}

$\ldots$

- {\it Is this life worth living?}$\ldots$ - Yes!

- {\it Can we atone{\tt/}make amends later for all the bad{\tt/}wrong things we have done in this life after we move into the other side?}
\begin{quotation}    
	BoJack Horseman: {\it ``Is it terrifying?''}
	
	Herb Kazzaz: {\it ``No. I don't think so. It's the way it is, you know? Everything must come to an end, the drip finally stops.''}
	
	BoJack Horseman: ``{\it See you on the other side}.''
	
	Herb Kazzaz: ``{\it Oh, BoJack, no, there is no other side. This is it}.'' - BoJack Horseman vs. Herb Kazzaz,  Episode: {\it The View from Halfway Down}, {\it BoJack Horseman} (2014--2020)
\end{quotation}
$\ldots$ {\it Just keep thinking \& digging, deeper \& harder. There will be a day that monster will be revealed \& the hidden truth will be exposed}$\ldots$

\begin{verbatim}
	// * * * * * * * * * * * * * * * * * * * * * * * * * * * * * * * * * * * * * //
\end{verbatim}

\noindent
\fbox{\textbf{12}} Cuối cùng thì sau 10 tháng đầu tiên của hành trình nghiên cứu sinh đầy rẫy gian nan và thử thách, mình cũng giải quyết được vấn đề tâm lý mà mình chịu đựng suốt thời gian qua. Chỉ tiếc là một bài toán Phổ thông nho nhỏ cho tụi học sinh Đức mà mình tự chế bị ầm thầm phá hoại và không được published online. Tiếc là một thư mục khác trong laptop cá nhân của mình thì bị xóa (do bị ngó password lúc thảo luận mà mình không hề{\tt/}chịu đề phòng) nhằm tạo ra mâu thuẫn giữa 3 đứa PhDs cùng phòng. Và đáng tiếc là một vài đóng góp khác bị phủ nhận công sức tuyệt đối mặc dù mình đã dành nhiều tâm huyết để thực hiện (see, e.g., \cite{phill2018,Young-Powell2018}).

Nhưng giờ thì mọi thứ đã trở nên rõ ràng và hợp lý hơn rất nhiều. Mình cũng chả trách gì, xem ra cái suất Marie--Curie này chả dễ ngốn tí nào (see the similar situation for this Marie--Curie PhD \cite{FundamentalPessimist2020}). Mình cũng đã giải thích sự tình và hòa giải với anh bạn Đức khó tánh ở đoạn \fbox{3}, vì mình cảm nhận, thông qua vài hành động của ảnh, cả 2 thằng mình đều là những kẻ chính trực ngu ngốc, khờ dại, và đáng thương hại. Cả 2 thằng đều bị kẻ thao túng tâm lý giở trò{\tt/}đâm thọt sau lưng, tánh ảnh nóng nên quạo, còn tánh mình hiền nên quỵ.

Ai dè cứ nói thẳng ra là mọi chuyện lại đâu vào đấy. Sau đó căng thẳng của phòng 3 đứa PhDs bọn mình cũng được giải tỏa, và chìa khóa để giải quyết vấn đề giữa người với người (có lẽ) nằm ở lòng chân thành và sự thấu cảm. Haizz$\ldots$ chỉ tiếc là cái CV của mình mất đi 1 dòng trong mục Referees. Nhưng vui cái là sau khi giải thích cặn kẽ sự tình với ông thầy chính ({\it real boss}) của mình, bắt đầu từ bây giờ mình sẽ được làm trực tiếp với ổng luôn. Với tại mình nói mình đã upgrade xong software trong thời gian ngắn dị thường (và chỉ ra luôn là nó chạy sai!), thành ra ổng tưởng thằng này bố láo nên bắt mình viết báo cáo! Mặc dù biết là {\it sẽ vả}\footnote{\textbf{sex var}(-iable) in \texttt{Unicode/Telex}.} hơn rất nhiều vì phải báo cáo thường xuyên hơn và phải làm một mình mà không bị ai phá, à nhầm, phụ, nhưng chỉ cần không bị đì{\tt/}kiểm soát{\tt/}thao túng nữa, là cỡ nào mình cũng ráng chơi{\tt/}quất{\tt/}tém hết! Không phải khoe, à mà khoe cũng chả sao, vì ổng là Viện trưởng, phải điều hành rất nhiều thứ nên quá bận, thành ra mới tạo cơ hội cho kẻ hám quyền lộng hành. Mọi thứ làm mình vô cùng confused trước đây cuối cùng cũng đã trở nên sáng tỏ và hợp lý hẳn!

Nên nhớ rằng, như một phép ví von dở tệ và làm nổi da gà, không cần biết bạn gieo bao nhiêu hạt giống tốt đẹp vào những mối quan hệ với người khác, kẻ thao túng tâm lý sẽ luôn dùng tiểu xảo{\tt/}thủ đoạn để chúng chết trước khi kịp nảy nầm, sau đó sai khiến{\tt/}thao túng khiến mọi người xung quanh trở nên kỳ thị{\tt/}xung đột và giẫm đạp một cách tàn nhẫn lên cả bạn và miếng vườn của bạn. Cách duy nhất để vượt qua điều này là bạn phải tìm{\tt/}săn lùng kẻ phá hoại, nếu không thì tất cả viễn cảnh tươi đẹp mà bạn hy vọng sẽ dần dần tàn lụi và lần lượt tiêu biến, bao gồm luôn cả bạn!$\ldots$

{\it - Why were I called the Lucky Guy? Were I really lucky?}

- Yeah, I were lucky enough to experience the worst nightmare of a PhD student: the situation in which his{\tt/}her (co-)supervisor is a real master psychological manipulator who always tries to find the way to ruin his{\tt/}her work and{\tt/}or life secretly.

Nevertheless, in return, I have learned a lot from such kind of person.

\begin{quotation}
	``{\it A wise man can learn more from his enemies than a fool from his friends}.'' - Niki Lauda, {\it Rush} (2013)
\end{quotation}
Love, trust, \& respect are so special. I cannot force{\tt/}manipulate anyone to give these things to me for free: Honestly, I do not have any right to do that \& these things actually need {\it building{\tt/}giving}, instead of {\it forcing{\tt/}taking}$\ldots$ {\it The harder you force someone to love{\tt/}trust{\tt/}respect you, the stronger{\tt/}deeper they hate{\tt/}distrust (mistrust){\tt/}disgust (disrespect) you instead.}

{\it It is all about a well-deserved victory indeed}$\ldots$

\begin{verbatim}
	// * * * * * * * * * * * * * * * * * * * * * * * * * * * * * * * * * * * * * //
\end{verbatim}

\subsubsection{{\color{gray}/* Quick sketches */}}
\fbox{\textbf{13}} I have been using \href{https://github.com}{GitHub} during the beginning of my PhD journey. It seems to me at this moment that \href{https://git-scm.com/}{git} is a perfect tool for blogging! This \texttt{git add-commit-push} working style is so perfect to motivate me \& push my workflow in a more disciplined way, which is highly suitable for an undisciplined shitty person like me.

{\it - How many commits have you attempted to push your limit today, man?}

- {\it You mean suicide?} Yeah$\ldots$ a couple of times.

- No, WTF, dude! I mean commit in the \texttt{git} sense. {\it What's wrong with your brain, bruh?}

\begin{verbatim}
	// * * * * * * * * * * * * * * * * * * * * * * * * * * * * * * * * * * * * * //
\end{verbatim}

\noindent
\fbox{\textbf{14}} {\it - Why did you remove a necessary supervision from your PhD journey? Don't you want another Vietnamese to be able to get this kind of ``prestigious'' fellowship?}

\begin{quotation}
	``{\it I can't concentrate. I can still hear the sound of the piano. It's in my way. You're in my way, so$\ldots$ Get out of here!}'' - Arima Kousei, {\it Your Lie in April}  (2014--2015).
\end{quotation}

- No, I did that because I really want it (Fig. \ref{fig4}).
\begin{figure}[h]
	\centering
	\includegraphics[width=0.75\textwidth]{Arima_Kousei_emotional_explosion}
	\caption{Arima Kousei's emotional explosion, {\it Your Lie in April} (2014--2015).}
	\label{fig4}
\end{figure}

{\it - Do you feel ashamed of that?}

\begin{quotation}
	``{\it It doesn't matter where you are, it's who you are, that's not gonna change whether you're in California or Maine or New Mexico. You know, you can't escape you}.'' - Penny Carson, {\it BoJack Horseman} (2014--2020)
\end{quotation}

- At the moment when you admit that you are stupid, you have eliminated a lot of inner sources of fear\footnote{\href{https://www.youtube.com/watch?v=bb9g9mtDHZo}{How To Beat Fear \& Anxiety | Jordan Peterson | Powerful Life Advice}.} \& shame from yourself: {\it You escaped you}.

{\it - But what if you fail then?}

- {\it Then what? So what?} As long as I have already tried the best of myself, nothing else matters$\ldots$

\begin{verbatim}
	// * * * * * * * * * * * * * * * * * * * * * * * * * * * * * * * * * * * * * //
\end{verbatim}

\noindent
\fbox{\textbf{15}} {\it - Is this world a place for good{\tt/}kind people?}

- No, it is the place for all people, particularly good{\tt/}kind people, but the problem is: if you are, especially, a good person, you have to struggle to protect yourself from the other predators. {\it Do you want to be a prey, which sounds similar to ``afraid'', all your life?} 

My mom used to teach me that if I do good things to other people, they will definitely give good things back. {\it But is this true?} I doubt it. This human race makes me doubt it. However, the question following that one should not be like: {\it Should I become a good{\tt/}kind person or not?} Instead, it should be: {\it I definitely want to be a good{\tt/}kind person, but good{\tt/}kind with whom?}

Like the working principle of \texttt{git}, we will encounter a lot of complicated events{\tt/}situations{\tt/}troubles in our life, \& for each of them, we have our own list of choices{\tt/}options, like \texttt{branches}. At each turning point, we are able to develop several different future versions of ourselves, of our personality, \& this also affects a lot to our action{\tt/}reaction{\tt/}behaviors{\tt/}character in future.

{\it Then which set of choices is the best? Which branch is the optimal?} Nobody knows the answer (right?). We have to choose according to our considerations{\tt/}perceptions, which are usually based on our own experiences, at that very moment. If that choice works, it will become background{\tt/}fundamental (sometimes self-delusion!) step for the next turning point. \& if it does not, we also have an additional experience for the next turn instead. {\it Nothing is absolutely good or bad, just our perception towards a particular event{\tt/}idea{\tt/}problem{\tt/}situation.}

\begin{verbatim}
	// * * * * * * * * * * * * * * * * * * * * * * * * * * * * * * * * * * * * * //
\end{verbatim}

\noindent
\fbox{\textbf{16}} It seems to me that individuals with high integrity usually attracts mysteriously a lot of troubles, especially in social relationships. So I call each of them, including myself, a ``drama-magnet''.

\begin{quotation}
	``{\it Nhưng cũng có lúc mọi thứ không như là những gì ta muốn. Thế giới này vận hành theo cái cách luôn ghì ta xuống}.'' - Đen Vâu, {\it Bài này chill phết} (2019)
	
	``{\it You pray for rain, you gotta deal with the mud too. That's part of the deal}.'' - Denzel Washington
\end{quotation}
Do not blame yourself as a drama-magnet, you should be proud of it instead!
\\

- Mà bạn êi, bạn nghiên cứu sinh ngành toán sao không lo học toán, mà đi đọc mấy cái tâm lý này chi cho phí thời gian vại?

- À, là vì toán và tâm lý là 2 ngành học dễ gây trầm cảm nhất, cộng với việc chuyên ngành của mình là toán nữa, nên mình cày luôn tâm lý cho đủ bộ!

- Sugoiii!!!

\begin{verbatim}
	// * * * * * * * * * * * * * * * * * * * * * * * * * * * * * * * * * * * * * //
\end{verbatim}

\subsubsection{\texttt{{\color{RubineRed}return} {\color{Purple}<?>}; \}}}
\texttt{\small You expect this text to return what?} \texttt{0?}

\paragraph{{\it Emotions}: good{\tt/}bad? {\it Emotionless?}}
\noindent
\fbox{\textbf{17}} We start loving something{\tt/}someone, or being attracted by something{\tt/}someone, because that thing{\tt/} person seems to give us a lot of energy{\tt/}happiness{\tt/}emotions\footnote{E.g., music, movie, sport, subject, relationship, career, etc.} and{\tt/}or reduce{\tt/}release our pain{\tt/}suffering\footnote{E.g., lyrics, sex, alcohol, drugs, etc.}. But after playing with it{\tt/}him{\tt/}her, consuming it{\tt/}him{\tt/}her for a while, we get bored, we find another one, a new source of emotion generators\footnote{Oh my fucking god, Japanese! You always knew it, didn't you?}. \& that will be repeated \& repeated. These emotions are exactly the things that make each of us a human being (Fig. \ref{fig5})$\ldots$

{\it Is emotion a type of energy?} We are always carried{\tt/}driven by our own emotions, deliberately{\tt/}intentionally sometimes \& unconsciously some other times: {\it We live for it, work for it, fight for it, also being taken down because of it}$\ldots$ {\it We are just emotion's slaves: emotion hiders and{\tt/}or emotion seekers}$\ldots$
\begin{quotation}
	``{\it If you place your hopes in anything, they will be betrayed. Promises will go unfulfilled \& faith will let you down}.'' - Mitsuru{\tt/}326, {\it {\sc Darling} in the {\sc Franxx}} (2018)
	
	``{\it Damned human wannabes}.'' - Klaxosaur Princess{\tt/}001 about VIRM after killing Tarsier, {\it {\sc Darling} in the {\sc Franxx}} (2018)
	
	``{\it Is this what living is for you? We believed that by abandoning our ties \& embracing solitude, we could perfect ourselves, make ourselves stronger. Very well. I shall give you every ounce of strength that remains within me. Whether or not you can take over the controls will be up to you. I stake this planet's future on you two}.'' - Klaxosaur Princess{\tt/}001, {\it {\sc Darling} in the {\sc Franxx}} (2018)
	
	``{\it Perhaps some lives only shine when in unison with others}.'' - Klaxosaur Princess{\tt/}001 about Hiro \& Zero Two's bond, {\it {\sc Darling} in the {\sc Franxx}} (2018)
	
	``{\it Decide whether you want to fight or accept your ruin}.'' - Klaxosaur Princess{\tt/}001's last words before she sacrifices herself \& entrusts the world's fate to Hiro \& Zero Two, {\it {\sc Darling} in the {\sc Franxx}} (2018)
\end{quotation}

\begin{figure}[h]
	\centering
	\includegraphics[width=0.75\textwidth]{002_more_human}
	\caption{``{\it Hey, darling, do I seem a little more human now?}'' - Zero Two, {\it {\sc Darling} in the {\sc Franxx}} (2018).}
	\label{fig5}
\end{figure}
{\it Greed, Wrath, Pride, Lust, Envy, Sloth, Gluttony - \href{https://en.wikipedia.org/wiki/Seven_deadly_sins}{Seven deadly sins} (7 mối tội đầu)}$\ldots$

{\it Do you{\tt/}I really need that much to be able to \href{https://www.imdb.com/title/tt0454921/}{pursue happiness}?}$\ldots$

To be happy with others, you should be happy in the relationship with yourself 1st\footnote{Not (only) masturbation!}: {\it Happiness can be found not only from the outside, but also from the inside}.

{\it If the ultimate goal of a human being is happiness, can we adjust our attitude and{\tt/}or lower our hopes \& expectations towards anything in order to achieve that in an easier \& more pleasant way?} Perhaps you will have much less achievements, but, in return, the way that you achieve them is much happier, more pleasant \& desirable. {\it Don't you want it that way? Is that a way to increase our Emotional Intelligence ourselves?}

\begin{verbatim}
	// * * * * * * * * * * * * * * * * * * * * * * * * * * * * * * * * * * * * * //
\end{verbatim}

\paragraph{Importance of {\it understanding} \& then {\it applying psychology} in {\it life}}

\noindent\fbox{\textbf{18}} Several people misunderstand strength as angriness and{\tt/}or dominant aggression. But the real strength actually is the {\it calmness under the storm} (Fig. \ref{fig6})$\ldots$ {\it Listening \& observing to understand instead of watching \& judging}$\ldots$

\begin{figure}[h]
	\centering
	\includegraphics[width=0.75\textwidth]{boat_under_storm}
	\caption{Robin Williams's Prized Memento from {\it Good Will Hunting} (1997).}
	\label{fig6}
\end{figure}
{\it Have you ever noticed that there is the phrase ``logical'' at the tail end of the term ``psychological'', which seems to treat nonsense{\tt/}non-logical arguments{\tt/}events{\tt/}situations from human bullshit after its logical counterpart seems so useless?} If anything is not logical{\tt/}does not make sense in the 1st place, you should consider it with the psychological point of view instead, then perhaps it will be even more ``logical''\footnote{Like {\it strong solutions} vs. {\it generalized solutions} in the field of \href{https://en.wikipedia.org/wiki/Partial_differential_equation}{Partial Differential Equations} (PDEs).}.

{\it At the end, what really matters for our life?}
\begin{quotation}
	Diane Nguyen: {\it ``It's too late. What's done is done.''}
	
	BoJack Horseman: {\it ``No.''}
	
	Diane Nguyen: {\it ``There's nothing I can do, BoJack. I'm not real. None of this is.''}
	
	BoJack Horseman: {\it ``So, what do I do now?''}
	
	Diane Nguyen: {\it ``BoJack, it doesn't matter.''}
	
	BoJack Horseman: {\it ``Well, if it doesn't matter, can I stay on the phone with you at least?''}
	
	Diane Nguyen: {\it ``Okay.''}
	
	BoJack Horseman: {\it ``How was your day?''}
	
	Diane Nguyen: {\it ``Good.''}
	
	BoJack Horseman: {\it ``Yeah?''}
	
	Diane Nguyen: {\it ``Yeah. My day was good.''} - BoJack Horseman vs. Diane Nguyen,  Episode: {\it The View from Halfway Down}, {\it BoJack Horseman} (2014--2020)
\end{quotation}
Watch e.g., \href{https://www.youtube.com/watch?v=Pt21dU5Pu8g}{BoJack Horseman - The View from Halfway Down}.\hfill$\square$

\begin{verbatim}
	// ************************************************************************* //
\end{verbatim}

\begin{flushright}
	{\sc Berlin, Germany}. Jan 2021.
	
	This text is a part of my {\it personal project}: 
	
	\texttt{\#Series}: {\sc Lost in Germany}.
\end{flushright}

\begin{quote}
	{\sf Hồng [27--?; psychologist]}: Để tôi tìm cách mô tả cấu trúc của 1 nhóm bắt nạt dẫn đầu \& thao túng bởi 1 kẻ thao túng tâm lý (psychological manipulator) ở mức nguy hiểm vừa cho anh hiểu.
	
	{\sf Hồng [11--?; chess player]}: Tôi chả biết gì nhiều ngoài cờ vua hay cờ tướng cả. Có thể sử dụng cờ để cắt nghĩa không? Nếu có thể.
	
	{\sf Hồng [27--?; psychologist]}: Được chứ. Tưởng tượng 1 bàn cờ vua. Vua 1 bên là kẻ thao túng tâm lý (psychological manipulator in the role of predator). Vua bên còn lại sẽ là con mồi (prey) của kẻ thao túng tâm lý. Vị tiểu bạo chúa (the King), tức kẻ thao túng tâm lý sẽ chọn các người có nhân cách thuần tốt, chỉ biết cắm đầu học hoặc làm, dạng ngây thơ, làm con tốt thí, i.e., quân Tốt (Pawn) \& để che chắn khi bất cứ quân nào của đối phương tấn công hắn. Hắn sẽ phải tìm 1 quân Hậu (Queen). Đúng nghĩa \& đúng chức năng của quân Hậu, đây phải là 1 người con gái hoặc phụ nữ có tính tình cộc cằn, thô lỗ, có thể ăn nói ngang ngược, càn quét ngang dọc, ăn xiêng nói xéo, để tấn công con mồi. Ngoài ra cần vài quân Xe (Rook), những người có chút đặc tính tốt \& biết ơn, dạng trung thành, nhưng chửi thề kiểu thô tục như cách đi ngang dọc của quân cờ này. Ngoài ra cần ít nhất 1 quân Ngựa (Knight, or Horse better in this context), 1 con ả dốt tới mức không biết mình dốt, cực kỳ dễ xoay theo chiều gió, ẹo ẹo trước mặt người khác, khiêu vũ cái điệu nhảy ``đàn ông nhạy cảm là đồ hèn, còn phụ nữ sẽ có quyền vì là phái yếu''. [...] Câu hỏi sẽ là: {\it Liệu anh có thể thắng ván cờ này hay không nếu kẻ thao túng tâm lý được phép đi trước, tức là đã nhắm anh làm con mồi \& tấn công anh trước?}
	
	{\sf Hồng [11--?; chess player]}: Tôi không chắc. Nhưng trò chơi trí não (mind game) này có vẻ khó \& hình như không đáng để tôi chơi?
	
	{\sf Hồng [23--?; psychological manipulated survivor]}: Chính xác. Cái thời điểm kẻ thao túng tâm lý nhấc con tốt đầu tiên thì hắn sẽ phải thua. Hắn có thể sẽ thắng ván cờ thao túng, khiến con mồi đau khổ, đôi khi đến mức sống dở chết dở, muốn tự tử. Nhưng hắn sẽ thua trên 1 Bàn Cờ lớn hơn, đó là bàn cờ về lương tâm. Tại sao ư? Hắn đã chơi cái trò chơi thao túng quá nhiều ván cờ, nên những ai đủ trưởng thành, chính chắn đã biết tỏng cái mánh hèn hạ của hắn, nên họ sẽ cẩn thận với hắn hơn, không dễ mắc mưu cái trò thao túng bẩn thỉu của hắn nữa. Xét về nghĩa tạo ra các người sống sót điềm tĩnh (calm survivor) \& nghĩa xây dựng cộng đồng tích cực, thì kẻ thao túng tâm lý là 1 kẻ bại trận thực sự: {\it a real loser in the battles of positivity \& building psychologically stable positive communities}. Check Mate. Assholes.
\end{quote}

\subsubsection{Signs of psychological manipulators}

\begin{itemize}
	\item Kẻ thao túng tâm lý là kẻ xảo quyệt, luôn rình rập con mồi, nên bạn có thể vô tình hoặc hữu ý tìm thấy hắn núp ở 1 góc kín, góc khuất tầm nhìn nào đấy để theo dõi con mồi của hắn hoặc theo dõi bạn trong trường hợp bạn đang bị hắn nhắm tới.
	\item Kẻ thao túng tâm lý là bậc thầy trong việc tìm điểm yếu (Grand Master of seeking weaknesses) của đối phương, \& cực kỳ tài giỏi trong việc khơi dậy những khía cạnh đen tối \& xấu xa nhất của bản chất con người. Điều khá hay là chả có ai dạy 1 cách đàng hoàng cho kẻ thao túng tâm lý cả. Thường hắn sẽ quan sát cách đối xử vợ của cha hắn \& bắt chước theo, rồi phát triển các thủ thuật lên thành kỹ năng thao túng riêng của hắn. Tất cả những sự phát triển theo chiều hướng độc hại của kẻ thao túng tâm lý là dựa vào bản năng sinh tồn tự nhiên của kẻ săn mồi 1 cách hoàn toàn tự nhiên \& thuần túy về mặt sinh học của thuyết tiến hóa Darwin.
\end{itemize}
Khi bạn đối mặt với tay sai của 1 kẻ thao túng tâm lý, bạn phải giết hoặc làm suy yếu chính cái phần tốt lành \& phần ngây thơ của bạn, chúng sẽ hồi phục hoặc hồi sinh lại sau khi cuộc chiến đã kết thúc nên bạn cứ yên tâm, bạn phải giáng các đòn tâm lý vào cái phần xấu xí, độc ác mà kẻ thao túng tâm lý khơi gợi \& thúc đẩy sự phát triển ở những tay sai này. Có nghĩa là bạn không hại họ, mà bạn chỉ tấn công cái phần xấu của họ, khiến họ giác ngộ, để cho trống chỗ trong tâm trí \& linh hồn của họ, để cái phần tốt tánh có chỗ để sinh sôi nảy nở trở lại. Bạn thuần hóa tay sai của kẻ thao túng bằng cách đó. Còn nếu họ không còn phần tốt nào, hãy xem tay sai đó ngang hàng với kẻ thao túng tâm lý \& cứ thẳng tay tấn công các đòn đánh tâm lý.

\subsection{The Dark Triad: narcissist, sociopath, \& psychopath -- Bộ 3 đen tối: kẻ ái kỷ, kẻ chống đối xã hội, \& kẻ thái nhân cách}

\subsubsection{Narcissists -- Những kẻ ái kỷ}
\textbf{\textsf{Resources -- Tài nguyên.}}
\begin{itemize}
	\item \cite{MacKenzie2015} {\sc Jackson MacKenzie} (tác giả gay trẻ). {\it Psychopath Free: Recovering from Emotionally Abusive Relationships With Narcissists, Sociopaths, \& Other Toxic People}.
	\item \cite{Hare1999} {\sc Robert D. Hare}. {\it Without Conscience: The Disturbing World of the Psychopaths Among Us} (tạm dịch: {\it Không Lương Tâm: Thế Giới Xáo Trộn Của Những Kẻ Thái Nhân Cách Lẫn Trong Chúng Ta}).
\end{itemize}
\begin{flushright}
	Spring 2021. Berlin, Germany. -- Xuân 2021. Thủ đô Berlin, Đức.
\end{flushright}
Hắn thấy mình đang ngồi ăn gà với (chị) Trinh, để kết thúc 1 ngày đi chơi chung, sau khi ghé vài chỗ ở Friedrichstrasse, gần cái Đồng Hồ Thế Giới -- World Clock\footnote{\href{https://en.wikipedia.org/wiki/World_Clock_(Alexanderplatz)}{Wikipedia{\tt/}World Clock (Alexanderplatz).}} ở Alexanderplatz. Cả 2 ăn gà, tám chuyện, có vài cái đùi nhìn ngon ơi là ngon, cả đùi Trinh trong đó. 
\begin{quote}
	{\sf Trinh [26{\tt/}27?; Amazon junior software developer; covert aggressive narcissist, demanding taker]}: Không biết Minh Toàn giờ sao rồi ta?
	
	{\sf Hồng [25; mathematics PhD student; still agreeable giver]}: À, hồi năm 2 em có đi thi chung Olympic Toán sinh viên với ảnh ngoài Huế: ảnh siêu giỏi Toán, cả Giải tích lẫn Đại số, Vật Lý nữa, với đọc rất nhiều sách. Chị Thảo mà cho mượn cuốn nào là ảnh bắt xe tới KTX em lấy liền, ngay, \& luôn. Em chưa có điều kiện tiếp xúc nhiều nhưng có nghe chị Thảo nói ảnh đọc (bản dịch tiếng Việt từ bản gốc tiếng Nhật) {\it Rừng Na Uy} \cite{Murakami_rung_Na_Uy} (bản dịch tiếng Anh {\it Norwegian Wood} \cite{Murakami_Norwegian_wood}) của {\sc Haruki Murakami} nữa. Chắc chắn là ảnh hướng nội, giống giống em, mà trùm hơn nhiều mặt.
\end{quote}
Hắn thoáng nghĩ, mặc dù không có nhiều thông tin lắm về anh Toàn trước hắn 1 khóa, nhưng dễ dàng nhận ra anh ta là kiểu người tốt, tuýp thích mấy thứ tư duy sâu sắc, hướng nội như hắn (maybe much more introverted but no way less), ít nói nhưng khi nói thì chắc lắm, nên chả thể nào rơi vào cạm bẫy của con bitch nói chuyện hời hợt thích lợi dụng bằng cách quyến rũ ngấm ngầm ở tầm thông minh bậc cao này được, dẫu 2 người nếu có làm ở chỗ thầy Dũng với nhau hay không đi chăng nữa. Mà tại sao phải làm vậy nhỉ? Có 1 lần hắn sơ suất nộp trễ học bổng rồi lúc lủi thủi về vô tình gặp anh Toàn, ảnh dặn mốt phải làm sớm hơn nên hắn biết ảnh tốt \& muốn trả ơn. Mà thôi, anh Toàn là người đủ thông minh \& sâu sắc nên chắc chả sao, hắn chả cần mắc công cảnh báo giúp chi cho mệt. Hắn xua cái ý nghĩ đó đi rồi ngắm mấy cái đùi mà chèm chẹp mút tiếp: {\it Juiciest Korean angry chickens in Berlin}! Tiềm thức của hắn luôn để ý tới ánh mắt của Trinh ở tầm spherical vision của hắn. Hắn quả thực có 1 spherical vision khá rộng, nên nhiều khi người ta cứ tha hồ liếc hắn mà cứ tưởng hắn đang quay chỗ khác hoặc không hề{\tt/}không thể chú ý đến họ, đặc biệt là những người đeo kính cận -- ánh mắt của họ hằn trên 2 miếng kính trong spherical vision của hắn. Nói chung thì hắn chả hiểu tại sao cái tiềm thức vô lý của hắn cứ chương chướng kiểu khó ở vậy. Mà lúc đó, sau khi check sơ bộ vài thông tin trên Facebook của chị Trinh cũng như chị ta đã check siêu kỹ thông tin của hắn để có thể dễ dàng quá trình ``đồng bộ hóa'' -- ``synchronization'' nhằm khai thác sự đồng cảm 1 cái khéo léo đối với nhiều người khác nhưng lại vụng về đối với hắn, hắn vẫn chả hiểu tại sao chả có ai bằng tuổi của hắn hoặc ít hơn, thêm kính ngữ `chị' (hay `senpai' như tụi Nhật) như hắn khi kêu con mụ này cả. Mãi nửa năm sau đó hắn mới hiểu: {\it 1 khi bạn chơi dơ, chơi bẩn có tiếng thì chả có đách đứa nào thèm nể bạn cả, dẫu cặp đùi bạn có ngon \& mọng nước cỡ nào đi chăng nữa}.

\paragraph{Signs -- Các dấu hiệu}
\begin{itemize}
	\item Cách xưng hô khá quan trọng. Nếu 1 người sử dụng xưng hô ``ta'' \& ``mi'' với bạn, coi thường, hạ thấp bạn, \& tính tình siêu ít kỷ, chỉ biết quan tâm đến lợi ích của bản thân, có thể đó là 1 dấu hiệu cho biết người đó là kẻ ái kỷ.
	\item Thường ưa dùng các từ dễ thương có vần ``um'', e.g., ``khum'' thay cho ``không'', nghe giống ``cum'', thường thích mượn các con vật dễ thương như ``cừu'' để che đậy bản chất là con sói gian manh, nham hiểm. Có thể phim \href{https://www.imdb.com/title/tt0102926}{\it The Silence of the Lambs} (1991) với tựa Việt {\it Sự Yên Lặng Của Bầy Cừu} có liên quan đến đặc điểm hành vi này.
	\item Đôi khi ra lệnh kẻ khác giúp như cướp giật, bốc lột (dấu hiệu của demanding taker). Nếu người được nhắm đến có ý định từ chối hoặc ngó lơ, thì kẻ ái kỷ sẽ ra vẻ dễ thương, quyến rũ, e.g., vài tấm hình để lộ 1 phần cơ thể (nhưng không quá hở hang, mà nửa kín nửa hở) hoặc gương mặt gợi tình nhưng lạnh lùng, hấp dẫn kiểu bí ẩn, etc., để dụ con mồi vào cạm bẫy, lưới tình, những cuộc kết đôi se duyên lý tưởng, nhưng sẽ không bao giờ xảy ra hoặc nếu có xảy ra thì cũng sẽ không có tình cảm chân thành vì sự thiếu vắng 1 phần chức năng quan trọng trong cấu trúc não của kẻ ái kỷ liên quan đến sự cảm thông, lòng thấu cảm với người khác, \& đặc biệt là sự thiếu vắng lương tâm (conscience) trong suy nghĩ, động cơ hành động, \& các hành vi biểu hiện ra bên ngoài.
\end{itemize}

\subsubsection{Sociopaths -- Những kẻ chống đối xã hội}
\textbf{\textsf{Resources -- Tài nguyên.}}
\begin{itemize}
	\item Quyển {\it The Sociopath Next Door} \cite{Stout2006} với bản dịch tiếng Việt {\it Kẻ Ác Cạnh Bên} \cite{Stout2019} của nhà tâm lý học \& tác giả người Mỹ {\sc Martha Stout}.
\end{itemize}

\subsubsection{Psychopaths -- Những kẻ thái nhân cách}
\textbf{\textsf{Resources -- Tài nguyên.}}
\begin{itemize}
	\item Quyển {\it Psychopath Free: Recovering from Emotionally Abusive Relationships With Narcissists, Sociopaths, \& Other Toxic People} \cite{MacKenzie2015} của tác giả gay trẻ {\sc Jackson MacKenzie}.
\end{itemize}

\subsection{Defense Against the Dark Triad 101 -- Lớp học phòng chống bộ 3 đen tối cơ bản}
Lấy cảm hứng từ bộ môn {\it Phòng chống Nghệ thuật Hắc ám -- Defence Against the Dark Arts} -- dạy các phương pháp \& kỹ thuật chống lại Nghệ thuật Hắc ám \& các sinh vật Hắc ám trong bộ truyện tuổi thơ nổi tiếng {\it Harry Potter} của nhà văn, nhà từ thiện, nhà sản xuất phim \& truyền hình, nhà biên kịch người Anh {\sc J. K. Rowling}, phần này được dùng để mô tả các bước cơ bản để phòng chống nghệ thuật thao túng hắc ám.

%------------------------------------------------------------------------------%

\section{On Teaching -- Bàn Về Việc Dạy}
\label{sect: teaching}

\begin{question}
	Bạn sẽ làm gì khi tình cờ dạy 1 đứa giỏi hơn bạn?
\end{question}
Tôi gặp Hồng [nam, 27 tuổi] đang loay hoay viết về buổi trò chuyện của hắn với các thầy cô giáo cũ dưới quê.

\begin{quote}
	{\sf Hồng [28; writer]}: Anh định viết thế nào?
	
	{\sf Hồng [27; NS teacher]}: Khó. Chả dễ. Viết lung tung cho đủ ý thì dễ, mà cho hay, cho trơn tru, đọc bắt tai thì khó quá xá.
	
	{\sf Hồng [28; writer]}: Nếu dạng trò chuyện, tâm sự thì anh có thể tham khảo phong cách đối thoại trong 2 cuốn sách {\it Dám Bị Ghét} \cite{Ichiro_Fumitake2022a} \& {\it Dám Hạnh Phúc} \cite{Ichiro_Fumitake2022b} của 2 tác giả Nhật Bản {\sc Kishimi Ichiro \& Koga Fumiake}.
	
	{\sf Hồng [27; NS teacher]}: Nội dung gì nhỉ?
	
	{\sf Hồng [28; writer]}: Bàn về thuyết tâm lý học trường phái Adlerian. Nguyên bản là cuốn {\it The Science of Living} \cite{Adler2013} của {\sc Alfred Alder}.
	
	{\sf Hồng [27; NS teacher]}: Để tui đọc thử. Hy vọng không phải mấy cái học thuyết nhảm địt chỉ lý thuyết suông mà không tí thực tế.
\end{quote}

Nghe lời tôi khuyên, hắn bắt đầu viết. Cụ thể như sau.

\subsection{Teaching kids in countryside -- Dạy trẻ vùng quê}
Tui tình cờ trò chuyện với Nhân [nam, 26 tuổi], 1 gia sư dạy các môn Tự nhiên như Toán Lý Hóa Tin, bên cạnh công việc nghiên cứu chưa đâu vào đâu của hắn, ở 1 vùng quê hẻo lánh.
\begin{quote}
	{\sf Hồng [27; writer]}: Thế anh thích dạy? Thích công việc gõ đầu trẻ?
	
	{\sf Nhân [26; Natural Science (NS for short) tutor -- gia sư Khoa học Tự nhiên]}: Cũng không hẳn. Không thích cũng không ghét. Thích vài cái \& cũng ghét 1 đống cái. Ban đầu nghe lời chị nên thử dạy, do công việc nghiên cứu bế tắc, hết đường tiến nên tạm lui về. Bế tắc sao thì sau tui sẽ kể chi tiết. Giờ tập trung vô việc dạy cái đã. Không kể liền có khi mất hồi nào không hay.
	
	Trước tui có về trường cấp 3 cũ để tham gia dạy đội tuyển học sinh giỏi của tỉnh nhà, hồi năm nhất, năm 2 Đại học. Đội tuyển chỉ có 6 đứa, chứ chưa được 8 hay 10 như của mấy tỉnh mạnh như Sài Gòn hay Hà Nội. Mà được cái 6 đứa giỏi, ngoan, chịu làm bài. Tui thích lắm, với hồi trước mấy thầy có phụ tiền cho tui lúc tui bị bệnh nên coi như là báo đáp cái ơn.
	
	Dưới quê thì khác hẳn. Mẹ nó cái vùng không có khỉ để ho mà cò cũng chả thèm gáy. Đa số học sinh không được giỏi cho lắm, toàn mấy dạng báo báo, mà chả dạng báo nào giống dạng báo nào. Những đứa vừa giỏi vừa ngoan, đủ trình để tui dạy hết sức, chắc hiếm như đếm số ngón tay của 1 đứa bị cùi.
	
	{\sf Hồng [27; writer]}: Anh cứ bình tĩnh, việc gì phải xỉ vả thế?
	
	{\sf Nhân [26; NS tutor]}: Tùy vào dạy ai, dạy cái gì, \& dạy ở mức độ nào.	
\end{quote}

\begin{quote}
	{\sf Students[6--18]}: Sao mấy thầy cô cứ khó chịu chuyện yêu đương? Người lớn chả hiểu gì cả.
	
	{\sf Hồng [27; NS teacher]}: Đấy là góc nhìn của bạn. Còn trẻ tui cũng thế. [...]
	
	Rồi cô ấy tự sát với đứa con trong bụng. Thế bạn còn muốn sống thử không?
\end{quote}

\begin{question}
	Bạn sẽ làm gì, trong tư thế \& với tư cách của 1 người thầy, người cô, nếu phát hiện 1 cuộc bắt nạt học đường trong lớp bạn dạy?
\end{question}

\begin{question}
	Bạn sẽ làm gì, trong tư thế \& với tư cách của 1 người thầy, người cô, nếu phát hiện 1 đứa học sinh ăn cắp đồ của người khác hoặc của chính bạn?
\end{question}

Yondu Rabit

\subsection{Parental role in child development -- Vai trò của cha mẹ trong sự phát triển của con cái}
Tôi dò hỏi Nhân về blackground của gia đình \& sự dạy dỗ của cha mẹ tác động lên sự phát triển của con cái họ thế nào.
\begin{quote}
	{\sf Hồng [28; psychologist]}: Tôi muốn thu thập vài dữ liệu để thống kê về tác động của cha mẹ lên sự phát triển tích cách \& tâm sinh lý của con cái. Anh có ý gì không?

	{\sf Nhân [26; NS tutor]}: Không hiếm trong dân gian những câu thành ngữ: ``Cha nào con nấy'', ``Hổ phụ sinh hổ tử'' hay English idiom ``Like father like son''. Không vào hang cọp sao bóp được cọp con. Kiểu vậy.
	
	{\sf Hồng [28; psychologist]}: Anh liều mạng vào hang cọp chỉ để bóp con cọp con thôi á? Sao anh liều cả tính mạng để đổi lấy cái ít giá trị như vậy?
	
	{\sf Nhân [26; NS tutor]}: Nếu bắt trúng con cọp dễ thương thì cũng xứng đáng lắm á. Mà chỉ ví von thui, anh đừng đi vô cái hang đầy cọp đó sâu quá. Dơ với nguy hiểm lắm. Trong lúc dạy, tui có soi, à nhầm, quan sát nhẹ mối tương quan của cha mẹ với con cái để đưa ra vài kết luận sau.
	
	{\sf Hồng [28; psychologist]}: Nhưng đấy chỉ là 1 vài trường hợp, có nghĩa là cỡ dữ liệu mẫu rất nhỏ (data sample with small size), thông tin cực kỳ hạn chế. Sao mà anh đưa đến kết luận được?
	
	{\sf Nhân [26; NS tutor]}: Cái hay là ở chỗ đó, cái trí nhớ điên khùng của tui giúp tui liên hệ lại khá nhiều các bạn bè, các anh chị hồi thời tui còn đi học. Phải chi nó bớt tệ đi, có khi tui lại nhớ được hết, nếu chịu để ý, không chừng. Nên thành ra cỡ mẫu có khi lại khá lớn, ít nhất là đủ lớn để bóp, không, để xài. Nói chung là nên thay kết luận thành phỏng đoán cho hợp lý \& không bị bắt bẻ.
	
	{\sf Hồng [28; pychologist]}: Anh hạn chế giúp tôi mấy trò đùa lại. Nghiêm túc lên. Cái phỏng đoán anh vừa nói là dự đoán thống kê. Thế vẫn là cái trí nhớ ảnh như trên phim ấy à.
	
	{\sf Nhân [26; NS tutor]}: Không, trí nhớ ảnh (photographic memory) ấy chỉ là khái niệm do phim ảnh thuần túy tưởng tượng ra thôi. Thực tế chỉ có trí nhớ điện tử (eidetic memory). Anh có xem phim Hannibal chứ? Cảnh ... tìm đến tên
	
	{\sf Hồng [28; pychologist]}: Có vẻ không thích hợp cho lắm khi bàn về phim ăn thịt người trong lúc làm khảo sát về sự phát triển của trẻ em, anh nhể?
	
	{\sf Nhân [26; NS tutor]}: À ừ nhỉ. Sở thích của tui hơi quái. Này nhé. Theo tui quan sát.
	\begin{itemize}
		\item Những đứa trẻ lớn lên trong gia đình nghèo, nhưng được giáo dục tốt, thường sẽ rất để ý chuyện tiền bạc, đặc biệt là tiền nợ người khác.
		\item Những đứa trẻ lớn lên trong các gia đình kinh doanh, giàu có, tài phiệt, lại thường giả sử mọi người phải phục vụ cho mình. Kiểu ta đây là vua chúa. Chúng ép tui phải dạy miễn phí, không là report, hoặc gieo tiếng xấu, tiếng oán cho lũ bạn để không ai dám vô học tui nữa.
		\item 
	\end{itemize}
	{\sf Hồng [28; pychologist]}: Là sao nhỉ? Người giàu thì mặc định phải phóng khoáng về mặt tiền bạc? Còn người nghèo thì thường tiết kiệm chứ hả?
	
	{\sf Nhân [26; NS tutor]}: Cái ngược ngạo ở chỗ đấy. 
	
	But, there is always an important jump right there. Like the jump over discontinuous in the shock wave solution of hyperbolic equations. You can see it right? Or at least feel it. Let me demonstrate it in an easier way so that you can comprehend.
	
	Excellent movies are always some steps ahead of science, why? Because they are made by gifted people, people with high sensitivity, they measure by their own sensitive functions, something deep inside their brain, their cognitive structure brain.
	
	You scientists always doubt them, consider them as mad men. But the right times will come, they are so damn right. Cinematic is the higher revelation of human nature.
	
	{\sf Nhân [26; NS tutor]}: Chính tui đây cũng là 1 ví dụ cho việc cha mẹ xây dựng nền tảng tính cách nào cho con cái. Cha tui là dân lao động chân tay ít học, dù có làm tới trưởng khoa 1 bệnh viện thời chiến nhưng chỉ học lớp bình dân học vụ. Tính toán vài phép tính đơn giản cũng sai. Nói chung là ổng mà làm kế toán công ty nào thì công ty đó xác định. Chả hiểu sao con của ổng kế thừa gen từ ai mà mê toán. Tui ngó kỹ hết mấy ông hàng xóm rồi nên cũng yên tâm. Được cái cha tui ổng thích vẽ, không đẹp đẽ gì nhưng vẽ vời nhiều, ổng thích nói lái, chế thơ tục, mấy câu đùa dâm dục. \& anh xem giờ tui như thế này đây. Chưa kể về xu hướng tính dục nữa. Tui còn nhỏ ngây thơ nên xem thử các đĩa phim trong túp lều của cha tui, \& giờ đây tui như thừa Testosterone vậy. Nản. Toàn bị gái dụ, đặc biệt là mấy con ái kỷ (narcissist) lợi dụng cái xu hướng tính dục đó mà quay vòng vòng như cái vibrator châu Phi rung lắc như The Rumble trong Attack on Titan vậy. Anh có thể tưởng tượng nổi không?
		
	{\sf Hồng [28; pychologist]}: Tôi từ chối tưởng tượng. Mặc dù không phải kiến thức có thể đưa vào sách giáo dục phổ thông hay tài liệu giáo dục giới tính cho lắm nhưng rất bổ ích.
\end{quote}
Hắn nói rồi ghi chép lại vào quyển sổ tay của hắn như gã trai drop-out Đại học trong \href{https://www.imdb.com/title/tt0159145/}{\it Golden Boy} (1995--1996) trên chiếc xe đạp của hắn, làm đủ nghề bán thời gian để học hỏi \& từ chối cám dỗ bởi phụ nữ: All healthy sex jokes are good for life. Healthy jokes, even dark ones, are vital for life.

\begin{example}[\sc Blaise Pascal]\rm
	Trong quyển {\it Pense\'e} \cite{Pascal2003}
\end{example}

\begin{example}[\sc Mike ``Iron'' Tyson]\rm
	Trong quyển {\it Undisputed Truth} \cite{Tyson_Sloman2013}, tay đấm thép Mike ``Iron'' Tyson
\end{example}

\subsection{Dare to take responsibility -- Dám lãnh trách nhiệm}
1 trong những điều khó nhất trong giáo dục tính cách của trẻ là hình thành tính dám nhận trách nhiệm trong mọi việc. Nếu làm được, có thể xem như việc giáo dục đã thành công theo nghĩa nhân sinh.

\begin{question}
	Khi nào nên nhận trách nhiệm? Phải nhận trách nhiệm ở mức bao nhiêu? Khi nào nên từ chối trách nhiệm?
\end{question}

\subsection{On teaching growth -- Dạy về trưởng thành}

\begin{quote}
	{\sf Student [13; 7th grader, Hitler's fan]}: Dạ con đủ 18 tuổi là trưởng thành rồi muốn làm gì làm phải hông thầy?
	
	{\sf Hồng [28; mathematics teacher; philosopher wannabe]}: Bậy bậy. Mà bạn đang nói về dạng trưởng thành nào? Trưởng thành về thể xác hay về tâm sinh lý? Trưởng thành về thể xác dễ, chỉ cần ăn uống, ngủ nghỉ đầy đủ là tới tuổi sẽ tự lột xác. Còn trưởng thành về tâm sinh lý khó hơn nhiều. Bạn sẽ dễ dàng gặp những người 40 mấy, 50 mấy tuổi mà tính cách vẫn trẻ con. Nhưng bạn cũng có thể gặp những đứa trẻ ra đời sớm để bương chải do hoàn cảnh gia đình, những đứa đó va chạm với đời sớm nên thường có cơ hội trưởng thành sớm hơn.
	
	{\sf Student [13; 7th grader, Hitler's fan]}: Vậy có cách nào mà trưởng thành nhanh lên không thầy?
	
	{\sf Hồng [28; mathematics teacher; philosopher wannabe]}: Có chứ. Bản chất của trưởng thành là gì? Là đau khổ. Nhưng chỉ đúng 1 phần chứ chưa đúng hẳn. Có những người vẫn đau khổ tới già vẫn chưa trưởng thành. Tôi nghĩ đau khổ là điều kiện cần của trưởng thành. Còn điều kiện đủ là người đó chiêm nghiệm để học được bài học gì từ những đau khổ đó để có kinh nghiệm để mà có thể sống tiếp 1 cách ngày càng trọn vẹn hơn. Nói 1 cách ví von là như thế này, có nhiều cây cảnh trong nhà, dẫu bạn có tưới chúng bằng những loại nước tinh khiết nhất, bón chúng bằng những loại phân tốt nhất, cho chúng hưởng ánh ban mai, chúng vẫn yếu ớt \& dần dần lụi tàn. Nhưng có nhiều cây, chả cần quan tâm chăm sóc gì cả, cứ quăng đại ra mưa gió, sấm chớm ầm ầm, có khi chúng lại bắt đầu nảy mầm. Đấy là cái sức sống mãnh liệt vượt qua mọi sự cản trở \& nghịch cảnh. 
	
	{\sf Student [13; 7th grader, Hitler's fan]}: Con người có thể trưởng thành như cây cối không thầy?
	
	{\sf Hồng [28; mathematics teacher; philosopher wannabe]}: [Trời đang mưa tầm tả] Cũng có. Bạn có thể chạy thẳng ra ngoài trời, 1 cú sét sẽ giúp bạn trưởng thành 1 cách nhanh chóng. Như 1 cú sốc trời giáng khiến mọi tế bào bắt buộc phải lớn lên.
	
	{\sf Student [13; 7th grader, Hitler's fan]}: Trưởng thành theo phong cách người thực vật.
	
	{\sf Hồng [28; mathematics teacher; philosopher wannabe]}: Rau củ được nấu nhừ bằng dòng điện tích electron tươi ngon đúng nghĩa. Just joking. Don't do it.
\end{quote}

\subsection{On specialization \& generalization -- Bàn về đặc biệt hóa \& tổng quát hóa}

\begin{itemize}\sf\small
	\item \textbf{specialization} [n] the process of becoming an expert in a particular area of work, study or business; the fact of spending more time on one area of work, etc. than on others; specialization (in something) a particular area of work, study or business which somebody spends more time on than on other areas.
	\item \textbf{generalization} [n] a general statement that is based on only a few facts or examples; the act of making such statements.
\end{itemize}

\noindent\textbf{\textsf{Resources -- Tài nguyên.}}
\begin{itemize}
	\item \cite{Polya2014}. George Polya. {\it How to Solve It}.
\end{itemize}
\begin{quote}
	{\sf Hồng [27; STEM teacher]}: Khi dạy học sinh bất cứu mảng gì mới, tôi bắt chúng đặc biệt hóa (specialization) để xét các trường hợp đặc biệt nếu chưa rõ cấu trúc 1 bài toán. Nếu đã giải được bài toán đó rồi, thì sau đó tổng quát ``bài toán'' đó lên. ``Bài toán'' ở đây được hiểu là 1 vấn đề của môn Toán sơ cấp, Vật lý sơ cấp, Hóa học sơ cấp, hoặc Tin học lập trình cơ bản. Tổng quát thế nào? Đầu tiên thay các số cụ thể bởi các tham số. Xem liệu bài toán có còn giải chung cách giải như vừa làm hay không. Sau đó, tổng số số đối tượng của bài toán lên thành $n\in\mathbb{N}$, e.g., nếu bài toán có 2 điểm hay 3 người, tôi yêu cầu học sinh tổng quát lên thành $n$ điểm \& $m$ người, với $m,n\in\mathbb{N}^\star$. Tiếp theo là tổng quát lên các tập số, e.g., nếu bài toán chỉ cho trên tập số tự nhiên $\mathbb{N}$, tôi sẽ hỏi học sinh là bài toán này nếu tổng quát lên tập số nguyên $\mathbb{Z}$, tập số hữu tỷ $\mathbb{Q}$, tập số thực $\mathbb{R}$, thậm chí là tập số phức $\mathbb{C}$ nếu đấy là học sinh cấp 3, lớp 11, 12, thì có còn đúng \& sử dụng được phương pháp giải hay phương pháp chứng minh của bài toán gốc hay không. Cấu trúc của bài toán là gì \& bản chất của bài toán cùng bản chất của phương pháp bạn vừa chứng minh là gì? Có mở rộng, nhờ bản chất vừa tìm được, các chứng minh đó cho các bài toán khác hay không.
\end{quote}


%------------------------------------------------------------------------------%

\section{On Learning -- Bàn Về Việc Học}
\label{sect: learning}

\begin{quote}
	{\sf Hồng [18--?; self-learner]}: It is never about how many books you read. It is about how you comprehend some of the most useful ones \& how deep you can connect them together in the grandiose theme.
	
	Chưa bao giờ là việc bạn đọc nhiều sách cỡ nào. Mà điều quan trọng là bạn có thể hiểu sâu những cuốn hữu ích nhất như thế nào \& kết nối chúng vào 1 chủ đề vĩ đại, lớn lao hơn ở mức độ sâu sắc nào nữa.
\end{quote}


\begin{question}
	Bạn sẽ làm gì khi học 1 người có kiến thức mà bạn tự cho là yếu hơn bạn?
\end{question}

\begin{question}
	Đích đến cuối cùng của việc học là gì?
\end{question}

\noindent\textbf{\textsf{Resources -- Tài nguyên.}}
\begin{itemize}
	\item \cite{Long2021}. Vũ Hoàng Long -- Người Kể Chuyện. {\it Học Trường Chuyên -- Những Góc Nhìn Đa Chiều}.
\end{itemize}

A decomposition of knowledge -- 1 phân hoạch của sự hiểu biết:
\begin{itemize}
	\item Những điều ta chưa biết.
	\begin{itemize}
		\item Những điều ta chưa biết nhưng biết là ta chưa biết.
		\item Những điều ta chưa biết \& chưa biết là chưa biết. Dốt tới mức chưa biết mình dốt thuộc phân loại này.
	\end{itemize}
	\item Những điều ta đã biết.
	\begin{itemize}
		\item Những điều ta đã biết \& tin tưởng ta đã biết rõ.
		\item Những điều ta biết đôi chút \& còn nhiều nghi ngờ chưa sáng tỏ về điều đó.
		\item Những điều ta nghĩ ta đã biết nhưng thật ra ta chưa biết gì cả. Illusion.
	\end{itemize}
\end{itemize}

\begin{question}
	Where does the real wisdom belong to?
\end{question}

\begin{quote}\it
	``Các em có thể trở thành những người lao động chân chính, những nhà kỹ thuật có chuyên môn giỏi, những nhà nghiên cứu thành công, những doanh nghiệp tầm cỡ, những nhà lãnh đạo xuất sắc, $\ldots$ nhưng trước hết phải là những người tử tế.'' -- {\sc Văn Như Cương}
\end{quote}

%------------------------------------------------------------------------------%

\section{A Bullshit Theory on Living -- 1 Thuyết Nhảm Nhí Về Việc Sống}
\label{sect: bullshit theory on live}

\begin{quote}\it
	``No matter what anybody tells you, words \& ideas can change the world.'' -- {\sc Tom Schulman}, \href{https://www.imdb.com/title/tt0097165}{Dead Poets Society} (1989)
\end{quote}
{\it Tạm dịch}: Bất luận ai nói gì với bạn đi nữa, ngôn từ \& ý tưởng có thể thay đổi cả thế giới.

\subsection{Love, Death, Robots, \& Artificial Intelligence -- Tình yêu, cái chết, người máy, \& Trí tuệ nhân tạo}

\subsection{Heavens on Earth -- Thiên đàng trên thế gian}

\begin{question}
	What is living? Why do we live? What to live for? Who to live with? How to live?
	
	-- Sống là gì? Tại sao ta sống? Ta sống vì điều gì? Ai mà ta sống chung với? Sống thế nào?
\end{question}
The authors devote this section to extend some aspects in some humble senses, with all respects, the theory of {\it Adlerian psychology} proposed by the Austrian psychiatrist {\sc Alfred Adler}.

The main sources are \cite{Adler2013}.

\begin{Rule}[On stupidity -- Bàn về sự ngu dốt]
	Take responsibility of your stupidity. Do not let your stupidity, no matter if you are aware of it or not, harm or even destroy any person.
	
	Chịu trách nhiệm cho sự ngu dốt của bạn. Đừng để sự ngu dốt của bạn, dù bạn có nhận thức được nó hay không đi nữa, làm hại hoặc hủy hoại người khác.
\end{Rule}


Maps of Meanings.

Let's hunt some low self-esteem.

Let's cook.

\begin{center}\it
	{\sc Life's recipe}
	
	A little bit of trusts. A little bit of betrayals.
	
	A little bit of loves. A little bit of denials.
\end{center}


\begin{Rule}[On deserving responsibility -- Bàn về sự xứng đáng về trách nhiệm] 
	Only take responsibility for whom deserved your responsibility. -- Chỉ chịu trách nhiệm cho ai xứng đáng với trách nhiệm của bạn.
\end{Rule}

\subsection{The last moonwalk in the Inferno -- Điệu nhảy moonwalk cuối cùng ở Hỏa Ngục}

\section{Miscellaneous -- Linh tinh}

\section{Acknowledgment -- Lời tri ân}
We, the authors, apologize to save this part for the last. Set the acknowledgment switch variable is $\delta_{\rm ack}$. Đặt biến công tắc cảm ơn là $\delta_{\rm ack}$.

\begin{verbatim}
if project_evaluation = fail
    take_all_responsibility = enabled;
else {
    take_all_credit = disabled;
    distribute_credit;
    acknowledgment;
}
\end{verbatim}
Xin cảm ơn các bài viết sâu sắc về giáo dục của thầy{\tt/}Dr. {\sc Trần Nam Dũng}, về giáo dục \& nghiên cứu Toán cao cấp của thầy{\tt/}Prof. {\sc Nguyễn Hữu Việt Hưng}, Prof. {\sc Ngô Bảo Châu}, Prof. {\sc Hà Huy Khoái}, Prof. {\sc Phan Thành Nam}. 2 tác giả không hề có bất cứ trao đổi nào với các người thầy này, nhưng bài viết của họ đã đủ để các tác giả bắt tay vào chấp bút cho 1 số ý tưởng của phần bàn về việc học \& bàn về việc dạy. Nếu có bất cứ lỗi nào xảy ra trong tiểu thuyết, đó hoàn toàn là lỗi của 2 tác giả đã hiểu sai ý từ các bài viết gốc. Hoàn toàn không liên quan đến ý tốt hay các tư tưởng của những người thầy này.

Xin chân thành những người đã từng giúp, đã từng thương hoặc thương hắn đến tận ngày hôm nay. Cuộc hành trình mưu cầu hạnh phúc \& khám phá tri thức sẽ thật khó khăn nếu không có họ củng cố niềm tin của 2 tác giả. Xin cảm ơn chị {\sc Lê Thị Minh Thảo}, anh {\sc Hoàng Công Đức}, anh{\tt/}thầy {\sc Lê Phúc Lữ}, anh{\tt/}thầy {\sc Khang}, chị {\sc Hương}, thầy{\tt/}anh {\sc Trà Quốc Khanh}, thầy{\tt/}anh {\sc Lê Văn Chánh}, cô {\sc Lê Thị Thanh Lĩu} đã truyền động lực trên hành trình học Toán của 2 tác giả. Xin cảm ơn cô {\sc Hạnh}, cô {\sc Đặng Thị Bích Thư} đã viếng đám tang của cha 2 tác giả, để họ không bị nỗi sợ bị bỏ rơi nuốt chửng. Xin cảm ơn thầy {\sc Lê Hoàng Minh}, thầy {\sc Nguyễn Thanh Tài}, thầy {\sc Võ Văn Huynh}, thầy {\sc Đệ}, thầy {\sc Lê Thanh Hải}, thầy {\sc Trần Thanh Liêm}, thầy {\sc Nguyễn Văn Quí}\footnote{Rest in peace, my respected mathematics teacher.}, thầy {\sc Trần Nam Dũng}, thầy {\sc Nguyễn Tấn Trung}, thầy {\sc Đặng Đức Trọng}, đã là những mẫu hình người thầy tuyệt vời để 2 tác giả bắt chước theo.

Bạn đọc có thể dễ dàng nhận thấy đây là 1 chủ đề vô cùng nhạy cảm, cực kỳ nguy hiểm. Nên nếu dự án này thất bại, thì đều do sự non nớt của 2 tác giả với ngòi bút chưa đủ sắc bén nên không thể giải phẫu vấn đề này đến nơi đến chốn, không liên quan gì đến bất cứ ai \& khẳng định là không ai giúp 2 tác giả trong tiểu thuyết này. Các tác giả xin hoàn toàn chịu trách nhiệm về sự thất bại nếu có.

%------------------------------------------------------------------------------%

\appendix

\section{Lists -- Các Danh Sách}
We, the authors, love making \st{love} lists. Chúng tôi, 2 tác giả, thích tạo ra các danh sách.

\subsection{A summary list of principles -- Danh sách tổng hợp các nguyên lý}

\subsection{A summary list of rules -- Danh sách tổng hợp các nguyên tắc}

\subsection{A summary list of $\Psi$-theorems, $\Phi$-theorems -- Danh sách tổng hợp các ``định lý'' về tâm lý \& triết học}

\section{Authors Bibliography -- Đôi điều về các tác giả}

\begin{itemize}
	\item {\sc Nguyễn Quản Bá Hồng} (1996--?): Tự thân vận động để đậu lớp chuyên Toán khóa 2011--2014 dưới sự ngăn cấm thi trường chuyên của cha mẹ. Kẻ hủy diệt ngôn ngữ khiến gần như tất cả các giáo viên dạy Văn của trường chuyên đều ghét, hoặc ít nhất là không ưa. Kẻ dốt tiếng Anh nên du học mà chưa có bằng Ielts hay Toefl. Học Thạc sĩ năm cuối, bỏ Thạc sĩ năm đầu (Master 2, skipped Master 1) ở Đại học Rennes 1 (Universit\'e de Rennes 1, France) trong lớp ENS của Rennes \& ENS Paris trộn lại. Tiến sĩ Toán Tối Ưu {\it hụt} tại Đại học Humboldt ở Berlin, Đức.
	
	Yêu chó (dog(gy) lover). Yêu \st{sex} sách (book lover). Yêu sự chân thành, ghét sự giả dối. Sẵn sàng \st{từ bỏ xu hướng tính dục \& dục vọng} dùng vài năm không bon chen để kiếm việc lương cao, sẵn sàng thí nghiệm tâm lý lên chính bản thân để đạt được mục tiêu riêng nhưng có thể có ích chung, nhốt mình ở nhà để viết 1 cuốn \st{tự truyện (autobiography)} tiểu thuyết hư cấu (fictional novel) để khóa mõm các kẻ thích bắt nạt trí tuệ \& nghiên cứu về các nhân cách độc hại, có tác động xấu đến người khác \& các cộng đồng.
	
	Hắn là kiểu sinh ra trong gia cảnh nghèo mà có tính cầu toàn khó ưa. Hắn có thể phá hủy cả 1 tác phẩm do hắn làm (hắn tôn trọng bất cứ tác phẩm nào của người khác nên không có chuyện hắn phá tác phẩm của người khác) chỉ đơn giản vì hắn không thích 1 (vài) chi tiết nào trong tác phẩm đó mà hắn không thể sửa được. Điểm này giống nhân vật {\sc Rust Cohl} trong series film \href{https://www.imdb.com/title/tt2356777/}{True Detective}:
	\begin{quote}\it
		``Rust would pick a fight with the sky if he didn't like its shade of blue. But when we finally got him over to the house - this is when that case was hot - the bastard looks like he was on his way to firing squad.'' -- \href{https://www.imdb.com/title/tt2356777/}{True Detective} (2014--) [S1.E1].
	\end{quote}
	\item {\sc Nguyễn Quản Trung Nhân} (?--?): Classified information -- thông tin tuyệt mật. You can read some from this book, but there is no information given further.
\end{itemize}

\subsection{Price -- Giá bán}
Tiểu thuyết này cùng các tài liệu liên quan được bán với giá $x$ VND{\tt/}euro, với $x\in[0,\infty)$ là 1 số thực không âm, i.e., muốn trả bao nhiêu cũng được, tùy vào đánh giá độ hay theo nhận thức của bạn đọc.

{\it Lý do}: Nếu $x < 0$ thì bỏ 1 đống công sức ra phá sản thì không được thông minh cho lắm, không xứng đáng để dạy Toán \& đầu tư kinh doanh. Nên các tác giả từ chối bán với giá là 1 số thực âm. Còn nếu $x\in\mathbb{C}\backslash\mathbb{R}$, i.e., $x = a + bi$ với $a,b\in\mathbb{R}^\star$ thì phải xem lại `đồng tiền ảo' như Bitcoin, hay tệ hơn là Luna Coin này xài được hay lâu dài hay không.

%------------------------------------------------------------------------------%

\printbibliography[heading=bibintoc]
	
\end{document}